% Réduction des conséquences néfastes des activités humaines sur les ressources naturelles — SVT 1ère D — Cours signé OnBuch+
% Programme officiel MINESEC. Compiler avec : tectonic ND07-reduction-consequences-activites-humaines-ressources-d.tex
\newcommand{\DOCMATIERE}{SVT}
\newcommand{\DOCNIVEAU}{1ère D}
\newcommand{\DOCMODULE}{Module 1 : Le Monde Vivant}
\newcommand{\DOCLECON}{ND07}
\newcommand{\DOCTITRE}{Réduction des conséquences néfastes des activités humaines sur les ressources naturelles}
% OnBuch+ — préambule « cours signés » ENRICHI (compilé avec tectonic / XeLaTeX)
% Système visuel « Pop » d'OnBuch+ : fond crème (jamais blanc), bordures encre
% épaisses, ombres « dures » décalées, titres Archivo Black en capitales.
% Version MATHÉMATIQUES : graphes (pgfplots/tikz), boîtes pédagogiques
% (définition, propriété, méthode, attention, le savais-tu, exercice résolu,
% exercice, TP, bilan), page de garde et signature OnBuch+ ancrée en bas.
%
% Macros attendues dans main.tex AVANT \input{preamble} :
%   \DOCMATIERE  \DOCNIVEAU  \DOCMODULE  \DOCLECON (numéro)  \DOCTITRE
\documentclass[11pt]{article}

\usepackage{fontspec}
\usepackage[french]{babel}
\usepackage[a4paper,margin=2cm,top=3.4cm,bottom=2.8cm,headheight=1.4cm,footskip=1.3cm]{geometry}
\usepackage{amsmath,amssymb,amsfonts}
\usepackage{mathtools} % \xmapsto, \coloneqq... souvent utilisés par les modèles
\usepackage{cancel} % simplifications barrées (bilans de forces)
% \abs{z} (module, valeur absolue) et \norm{u} (norme), utilisés par les modèles
\providecommand{\abs}{}\let\abs\relax\DeclarePairedDelimiter{\abs}{\lvert}{\rvert}
\providecommand{\norm}{}\let\norm\relax\DeclarePairedDelimiter{\norm}{\lVert}{\rVert}
\usepackage[table]{xcolor} % [table] : fournit aussi \rowcolors (alternance de lignes)
\usepackage{pagecolor}
\usepackage{tikz}
\usetikzlibrary{arrows.meta,positioning,calc,shapes.geometric,shapes.misc,shapes.arrows,decorations.pathmorphing,decorations.pathreplacing,decorations.markings,patterns,shadows,mindmap,trees,fit,backgrounds,matrix,intersections,angles,quotes}
\usepackage{pgfplots}
\usepackage[version=4]{mhchem} % équations nucléaires \ce{^{235}_{92}U}
\usepackage[european,siunitx]{circuitikz} % schémas électriques
\pgfplotsset{compat=1.18}
\usepgfplotslibrary{fillbetween,groupplots} % aires entre courbes (calcul intégral)
\usepackage{enumitem}
\usepackage{fancyhdr}
% [most] charge toutes les bibliothèques tcolorbox (listings, minted, raster,
% documentation...) et gonfle inutilement la mémoire TeX sur un document long
% (~300 boîtes) : on ne charge que ce qui est réellement utilisé.
\usepackage[skins,breakable]{tcolorbox}
\usepackage{graphicx}
\usepackage{colortbl}
\usepackage{booktabs}
\usepackage{tabularx}
\usepackage{diagbox} % case d'en-tête diagonale des tableaux à double entrée
% Colonnes extensibles centrée / gauche / droite (C, L, R), souvent utilisées par les modèles
\newcolumntype{C}{>{\centering\arraybackslash}X}
\newcolumntype{L}{>{\raggedright\arraybackslash}X}
\newcolumntype{R}{>{\raggedleft\arraybackslash}X}
\usepackage{array}
\usepackage{multicol}
\usepackage{siunitx}
% parse-numbers=false : accepte aussi \SI{2,5 \times 10^{-3}}{...}
\sisetup{locale=FR,output-decimal-marker={,},group-minimum-digits=4,parse-numbers=false}
\DeclareSIUnit\ohm{\ensuremath{\symup{\Omega}}} % U+2126 absent de la police
\DeclareSIUnit\division{div} % réglages d'oscilloscope (ms/div, V/div)
\DeclareSIUnit\tr{tr} % tours (vitesse de rotation, ex. \SI{1500}{\tr\per\minute})
\DeclareSIUnit\ch{ch} % cheval-vapeur (puissance moteur, unité usuelle hors SI)
\DeclareSIUnit\franccfa{FCFA} % franc CFA (coûts, contexte camerounais)
\DeclareSIUnit\dioptre{\ensuremath{\delta}} % vergence d'une lentille (dioptrie)
\usepackage{qrcode}
% \degree : commande fréquemment utilisée par les modèles pour un angle ou une
% température en dehors de \SI{}{} (non fournie par les paquets chargés).
\providecommand{\degree}{^{\circ}}
% \qed (fin de démonstration), utilisé par les modèles sans amsthm
\providecommand{\qed}{\hfill$\square$}
% \barre : double barre des valeurs interdites dans un tableau de variations (tkz-tab)
\providecommand{\barre}{\ensuremath{\|}}
% environnement proof (amsthm non chargé) utilisé par les modèles
\ifdefined\proof\else\newenvironment{proof}[1][Démonstration]{\par\noindent\textit{#1.}\ }{\qed\par}\fi
\usepackage{lastpage}
\usepackage{hyperref}
\hypersetup{hidelinks}

% ============================================================
% Palette « Pop » OnBuch+ — mêmes hex que la classe P (plus_ui.dart)
% ============================================================
\definecolor{poporange}{HTML}{F59321}
\definecolor{popdark}{HTML}{E07A0C}
\definecolor{popcream}{HTML}{FFF6E8}
\definecolor{popink}{HTML}{1C1714}
\definecolor{poppurple}{HTML}{7A5AE0}
\definecolor{popgreen}{HTML}{2E9E6A}
\definecolor{poppink}{HTML}{E0457B}
\definecolor{popblue}{HTML}{2F7DE1}
\definecolor{popgold}{HTML}{C98A12}
\definecolor{popmuted}{HTML}{8A7F76}
\definecolor{popline}{HTML}{EADFD2}
\definecolor{popgray}{HTML}{8A7F76}
\colorlet{popgrey}{popgray}
% Alias tolérants (le modèle nomme parfois les couleurs différemment)
\colorlet{popwhite}{white}
\colorlet{popblack}{popink}
\colorlet{popred}{poppink}
\colorlet{popyellow}{popgold}
% Teintes claires pour fonds de tableaux / zones
\colorlet{popblueL}{popblue!10!white}
\colorlet{poporangeL}{poporange!14!white}
\colorlet{popgreenL}{popgreen!12!white}
\colorlet{poppurpleL}{poppurple!12!white}
\colorlet{poppinkL}{poppink!12!white}
\colorlet{popgoldL}{popgold!14!white}

\pagecolor{popcream}

% ============================================================
% Typographie
% ============================================================
\defaultfontfeatures{Ligatures=TeX}
\setmainfont{PlusJakartaSans-Medium.ttf}[Path=../../../fonts/,BoldFont=PlusJakartaSans-Bold.ttf]
% Maths en sans-serif (Fira Math) avec chiffres et lettres droites de Plus
% Jakarta, pour que les formules se fondent dans le texte.
\usepackage[math-style=ISO,bold-style=ISO]{unicode-math}
\setmathfont{FiraMath-Regular.otf}[Path=../../../fonts/]
\setmathfont{PlusJakartaSans-Medium.ttf}[Path=../../../fonts/,range={up/num,up/{Latin,latin},\mathup}]
\usepackage{icomma} % virgule décimale sans espace en mode math
% Caractères mathématiques Unicode absents des polices de texte (Plus Jakarta,
% Archivo Black) : rendus par leur équivalent mathématique, en texte comme en
% formule — sinon ils s'affichent en carré vide (ex. « Arithmétique dans ℤ »).
\usepackage{newunicodechar}
\newunicodechar{ℝ}{\ensuremath{\mathbb{R}}}
\newunicodechar{ℤ}{\ensuremath{\mathbb{Z}}}
\newunicodechar{ℂ}{\ensuremath{\mathbb{C}}}
\newunicodechar{ℕ}{\ensuremath{\mathbb{N}}}
\newunicodechar{ℚ}{\ensuremath{\mathbb{Q}}}
\newunicodechar{𝔻}{\ensuremath{\mathbb{D}}}
\newunicodechar{α}{\ensuremath{\alpha}}
\newunicodechar{β}{\ensuremath{\beta}}
\newunicodechar{γ}{\ensuremath{\gamma}}
\newunicodechar{δ}{\ensuremath{\delta}}
\newunicodechar{ε}{\ensuremath{\varepsilon}}
\newunicodechar{θ}{\ensuremath{\theta}}
\newunicodechar{λ}{\ensuremath{\lambda}}
\newunicodechar{μ}{\ensuremath{\mu}}
\newunicodechar{σ}{\ensuremath{\sigma}}
\newunicodechar{τ}{\ensuremath{\tau}}
\newunicodechar{φ}{\ensuremath{\varphi}}
\newunicodechar{ψ}{\ensuremath{\psi}}
\newunicodechar{ω}{\ensuremath{\omega}}
\newunicodechar{Σ}{\ensuremath{\Sigma}}
% majuscules produites par \MakeUppercase dans les titres (α-aminés -> Α)
\newunicodechar{Α}{\ensuremath{\alpha}}
\newunicodechar{Β}{\ensuremath{\beta}}
\newunicodechar{Γ}{\ensuremath{\Gamma}}
\newunicodechar{Δ}{\ensuremath{\Delta}}
\newunicodechar{Ε}{\ensuremath{\varepsilon}}
\newunicodechar{Θ}{\ensuremath{\Theta}}
\newunicodechar{Λ}{\ensuremath{\Lambda}}
\newunicodechar{Μ}{\ensuremath{\mu}}
\newunicodechar{Τ}{\ensuremath{\tau}}
\newunicodechar{Φ}{\ensuremath{\Phi}}
\newunicodechar{Ψ}{\ensuremath{\Psi}}
\newunicodechar{Ω}{\ensuremath{\Omega}}
\newunicodechar{↦}{\ensuremath{\mapsto}}
\newunicodechar{∈}{\ensuremath{\in}}
\newunicodechar{∉}{\ensuremath{\notin}}
\newunicodechar{∧}{\ensuremath{\wedge}}
\newunicodechar{⇔}{\ensuremath{\Leftrightarrow}}
\newunicodechar{⟺}{\ensuremath{\Longleftrightarrow}}
\newunicodechar{⇒}{\ensuremath{\Rightarrow}}
\newunicodechar{⟹}{\ensuremath{\Longrightarrow}}
\newunicodechar{∘}{\ensuremath{\circ}}
\newunicodechar{∪}{\ensuremath{\cup}}
\newunicodechar{∩}{\ensuremath{\cap}}
\newunicodechar{≡}{\ensuremath{\equiv}}
\newunicodechar{⊂}{\ensuremath{\subset}}
\newunicodechar{⊥}{\ensuremath{\perp}}
\newunicodechar{∃}{\ensuremath{\exists}}
\newunicodechar{∀}{\ensuremath{\forall}}
\newunicodechar{∛}{\ensuremath{\sqrt[3]{\,}}}
\newunicodechar{∜}{\ensuremath{\sqrt[4]{\,}}}
\newunicodechar{⃗}{\ensuremath{{}^{\to}}}
\newunicodechar{ⁿ}{\ifmmode^{n}\else\textsuperscript{\ensuremath{n}}\fi}
\newunicodechar{ˣ}{\ifmmode^{x}\else\textsuperscript{\ensuremath{x}}\fi}
\newunicodechar{ᵇ}{\ifmmode^{b}\else\textsuperscript{\ensuremath{b}}\fi}
\newunicodechar{ʳ}{\ifmmode^{r}\else\textsuperscript{\ensuremath{r}}\fi}
\newunicodechar{ᵘ}{\ifmmode^{u}\else\textsuperscript{\ensuremath{u}}\fi}
\newunicodechar{ʸ}{\ifmmode^{y}\else\textsuperscript{\ensuremath{y}}\fi}
\newunicodechar{ᵃ}{\ifmmode^{a}\else\textsuperscript{\ensuremath{a}}\fi}
\newunicodechar{ᵉ}{\ifmmode^{e}\else\textsuperscript{\ensuremath{e}}\fi}
\newunicodechar{ᵏ}{\ifmmode^{k}\else\textsuperscript{\ensuremath{k}}\fi}
\newunicodechar{ᵖ}{\ifmmode^{p}\else\textsuperscript{\ensuremath{p}}\fi}
\newunicodechar{ᴺ}{\ifmmode^{N}\else\textsuperscript{\ensuremath{N}}\fi}
\newunicodechar{ᶜ}{\ifmmode^{c}\else\textsuperscript{\ensuremath{c}}\fi}
\newunicodechar{ʰ}{\ifmmode^{h}\else\textsuperscript{\ensuremath{h}}\fi}
\newunicodechar{ᵠ}{\ifmmode^{q}\else\textsuperscript{\ensuremath{q}}\fi}
\newunicodechar{ₙ}{\ifmmode_{n}\else\textsubscript{\ensuremath{n}}\fi}
\newunicodechar{ₐ}{\ifmmode_{a}\else\textsubscript{\ensuremath{a}}\fi}
\newunicodechar{ᵢ}{\ifmmode_{i}\else\textsubscript{\ensuremath{i}}\fi}
\newunicodechar{ᵣ}{\ifmmode_{r}\else\textsubscript{\ensuremath{r}}\fi}
\newunicodechar{ₖ}{\ifmmode_{k}\else\textsubscript{\ensuremath{k}}\fi}
\newunicodechar{ᵦ}{\ifmmode_{\beta}\else\textsubscript{\ensuremath{\beta}}\fi}
\newunicodechar{ₓ}{\ifmmode_{x}\else\textsubscript{\ensuremath{x}}\fi}
\newfontfamily\popdisplay{ArchivoBlack-Regular.ttf}[Path=../../../fonts/]
\newfontfamily\poplogo{SpaceGrotesk-Bold.ttf}[Path=../../../fonts/]
\newfontfamily\popsemi{PlusJakartaSans-SemiBold.ttf}[Path=../../../fonts/]
\newfontfamily\popxbold{PlusJakartaSans-ExtraBold.ttf}[Path=../../../fonts/]
\newfontfamily\popxboldls{PlusJakartaSans-ExtraBold.ttf}[Path=../../../fonts/,LetterSpace=3.0]

\setlist[enumerate]{leftmargin=1.7em,itemsep=5pt,topsep=4pt}
\setlist[itemize]{leftmargin=1.7em,itemsep=4pt,topsep=4pt,label={\color{poporange}\textbullet}}
\setlength{\parindent}{0pt}
\setlength{\parskip}{5pt}
\renewcommand{\arraystretch}{1.25}

% pgfplots : style Pop par défaut
\pgfplotsset{
  every axis/.append style={axis line style={popink,line width=1pt},tick style={popink},
    grid style={popline,line width=0.5pt},label style={font=\small},tick label style={font=\footnotesize}},
  popcurve/.style={poporange,line width=1.8pt,no markers,smooth},
}
% Même style disponible dans un \draw TikZ simple (hors pgfplots)
\tikzset{popcurve/.style={poporange,line width=1.8pt}}
\tikzset{popgrid/.style={popline,very thin}}
\colorlet{popcurve}{poporange} % le nom du style est parfois utilisé comme couleur

% ============================================================
% Pill / squiggle
% ============================================================
\newcommand{\poppill}[3]{%
  \tikz[baseline=-0.55ex]\node[draw=popink,line width=1.2pt,rounded corners=2.6pt,%
    fill=#2,inner xsep=6.5pt,inner ysep=3pt]{%
    \popxboldls\fontsize{7}{7}\selectfont\color{#3}\MakeUppercase{#1}};%
}
\newcommand{\popsquiggle}[1]{%
  \tikz[baseline=-0.4ex]{
    \draw[#1,line width=2.6pt,line cap=round]
      plot [smooth] coordinates {(0,0.09) (0.34,0.01) (0.68,0.21) (1.02,0.01) (1.36,0.10)};
  }%
}
% Mot-clé surligné (marqueur orange sous le texte)
\newcommand{\cle}[1]{{\popxbold\color{popdark}#1}}

% ============================================================
% En-tête / pied
% ============================================================
\pagestyle{fancy}
\fancyhf{}
\renewcommand{\headrulewidth}{0pt}
\renewcommand{\footrulewidth}{0pt}
\lhead{\raisebox{-0.15ex}{\poplogo\fontsize{13}{13}\selectfont\color{popink}OnBuch\color{poporange}+}}
\rhead{\poppill{\DOCMATIERE\ \textbullet\ \DOCNIVEAU\ \textbullet\ Leçon \DOCLECON}{white}{popink}}
\lfoot{\popxbold\fontsize{7.6}{7.6}\selectfont\color{popmuted}\MakeUppercase{Cours signé OnBuch+}\ \textbullet\ {\popsemi\DOCTITRE}}
\rfoot{\tikz[baseline=-0.6ex]\node[rounded corners=8.5pt,fill=popink,minimum height=17pt,minimum width=17pt,inner xsep=5pt,inner sep=0]{\popxbold\fontsize{8}{8}\selectfont\color{popcream}\thepage\,/\,\pageref*{LastPage}};}

% ============================================================
% Titres
% ============================================================
\newcommand{\chaptitre}[2]{%
  \begin{tcolorbox}[enhanced,colback=poporange,colframe=popink,boxrule=2.6pt,arc=11pt,
      drop shadow=popink,left=15pt,right=15pt,top=12pt,bottom=11pt,before skip=3pt,after skip=18pt]
    \poppill{#2}{popink}{popcream}\par\vspace{9pt}
    {\raggedright\hyphenpenalty=10000\exhyphenpenalty=10000\popdisplay\fontsize{22}{24}\selectfont\color{popcream}\MakeUppercase{#1}\par}
  \end{tcolorbox}
}
% Section numérotée : pastille orange avec le numéro + titre Archivo
\newcounter{popsec}
\newcommand{\coursec}[1]{%
  \stepcounter{popsec}\setcounter{popsub}{0}%
  \par\vspace{16pt}\noindent
  \begin{minipage}{\linewidth}
  \tikz[baseline=-0.7ex]\node[draw=popink,line width=1.6pt,fill=poporange,rounded corners=4pt,minimum size=22pt,inner sep=0]{\popdisplay\fontsize{12}{12}\selectfont\color{popcream}\thepopsec};%
  \hspace{8pt}{\popdisplay\fontsize{15}{17}\selectfont\color{popink}\MakeUppercase{#1}}\par
  \vspace{4pt}\noindent{\color{poporange}\rule{36pt}{3.6pt}}
  \end{minipage}\par\nopagebreak\vspace{8pt}
}
\newcounter{popsub}
\newcommand{\courssub}[1]{%
  \stepcounter{popsub}%
  \par\vspace{9pt}\noindent{\popxbold\fontsize{11.5}{13}\selectfont\color{popink}{\color{poporange}\thepopsec.\thepopsub}\hspace{6pt}#1}\par\nopagebreak\vspace{4pt}
}

% ============================================================
% Boîtes pédagogiques — même recette : fond blanc, bordure épaisse colorée,
% ombre dure encre, badge plein. Titre optionnel [#1].
% ============================================================
\tcbset{popbase/.style={breakable,enhanced,colback=white,boxrule=2.3pt,arc=9pt,
  shadow={3pt}{-3pt}{0pt}{popink},left=14pt,right=14pt,top=11pt,bottom=11pt,before skip=11pt,after skip=15pt,
  fontupper=\popsemi\fontsize{10.6}{14.5}\selectfont\color{popink},
  fontlower=\popsemi\fontsize{10.6}{14.5}\selectfont\color{popink}}}
% Coche dessinée (le glyphe ✓ n'existe pas dans les polices du document)
\newcommand{\popcheck}[1]{\tikz[baseline=-0.5ex]\draw[#1,line width=1.6pt,line cap=round,line join=round](-0.13,0.0)--(-0.04,-0.09)--(0.13,0.1);}
\providecommand{\coche}{\popcheck{popgreen}}
\providecommand{\resultat}[1]{\par\smallskip\noindent\fcolorbox{popgreen}{popgreenL}{\parbox{\dimexpr\linewidth-2\fboxsep-2\fboxrule\relax}{\textbf{Résultat.} #1}}\par\smallskip}
\newcommand{\popboxhead}[3]{\poppill{#1}{#2}{white}\ifx\relax#3\relax\else\hspace{6pt}{\popxbold\fontsize{10.6}{12}\selectfont\color{popink}#3}\fi\par\vspace{7pt}}

\newtcolorbox{aretenir}[1][]{popbase,colframe=popgreen,before upper={\popboxhead{À retenir}{popgreen}{#1}}}
\newtcolorbox{exemplebox}[1][]{popbase,colframe=poporange,before upper={\popboxhead{Exemple}{poporange}{#1}}}
\newtcolorbox{definition}[1][]{popbase,colframe=popblue,before upper={\popboxhead{Définition}{popblue}{#1}}}
\newtcolorbox{propriete}[1][]{popbase,colframe=poppurple,before upper={\popboxhead{Propriété}{poppurple}{#1}}}
\newtcolorbox{methode}[1][]{popbase,colframe=popgold,before upper={\popboxhead{Méthode}{popgold}{#1}}}
\newtcolorbox{attention}[1][]{popbase,colframe=poppink,before upper={\popboxhead{Attention}{poppink}{#1}}}
\newtcolorbox{experience}[1][]{popbase,colframe=popblue,colback=popblueL,before upper={\popboxhead{Expérience}{popblue}{#1}}}
% « Le savais-tu ? » : carte violette pleine (contexte, culture, Cameroun)
\newtcolorbox{savaistu}[1][]{popbase,colback=poppurple,colframe=popink,
  fontupper=\popsemi\fontsize{10.4}{14.2}\selectfont\color{white},
  before upper={\poppill{Le savais-tu\,?}{white}{poppurple}\ifx\relax#1\relax\else\hspace{6pt}{\popxbold\color{white}#1}\fi\par\vspace{7pt}}}
% Exercice résolu : énoncé en haut, \tcblower puis la solution sur fond crème
\newcounter{popexo}\newcounter{popexr}
\newtcolorbox{exoresolu}[1][]{popbase,colframe=poporange,colbacklower=popcream,
  segmentation style={popink,line width=1.2pt,dashed},
  before upper={\stepcounter{popexr}\popboxhead{Exercice résolu \thepopexr}{poporange}{#1}},
  before lower={\poppill{Solution}{popink}{popcream}\par\vspace{6pt}}}
% Exercice (non résolu en place) — la correction est dans la partie Corrigés
\newtcolorbox{exercice}[1][]{popbase,colframe=popink,
  before upper={\stepcounter{popexo}\popboxhead{Exercice \thepopexo}{popink}{#1}}}
% Corrigé
\newtcolorbox{corrige}[1][]{popbase,colframe=popgreen,colback=popgreenL,
  before upper={\popboxhead{Corrigé}{popgreen}{#1}}}
% Objectifs / prérequis / bilan
\newtcolorbox{objectifs}{popbase,colframe=popink,colback=poporangeL,
  before upper={\popboxhead{Objectifs de la leçon}{popink}{}}}
\newtcolorbox{prerequis}{popbase,colframe=popink,colback=popblueL,
  before upper={\popboxhead{Prérequis}{popblue}{}}}
\newtcolorbox{bilan}{popbase,colframe=popink,colback=white,boxrule=2.8pt,
  before upper={\popboxhead{Fiche bilan}{poporange}{L'essentiel à connaître pour l'examen}}}
% Figure encadrée avec légende
\newtcolorbox{popfigure}[1][]{enhanced,breakable=false,colback=white,colframe=popline,boxrule=1.4pt,arc=8pt,
  left=10pt,right=10pt,top=10pt,bottom=8pt,before skip=10pt,after skip=12pt,halign=center,
  lower separated=false,halign lower=center,
  fontlower=\popsemi\fontsize{9}{11.5}\selectfont\color{popmuted},
  #1}
\newcommand{\legende}[1]{\par\vspace{6pt}{\popsemi\fontsize{9}{11.5}\selectfont\color{popmuted}#1}}

% ============================================================
% Signature OnBuch+
% ============================================================
\newcommand{\poplogoinline}[1]{{\poplogo\fontsize{#1}{#1}\selectfont\color{popink}OnBuch\color{poporange}+}}
\newcommand{\signatureonbuch}{%
  \begin{tcolorbox}[enhanced,colback=popink,colframe=popink,boxrule=2.2pt,arc=10pt,
      drop shadow=poporange,left=14pt,right=14pt,top=10pt,bottom=10pt,before skip=0pt,after skip=0pt]
    \tikz[baseline=-0.7ex]\node[circle,fill=popgreen,draw=popcream,line width=1.3pt,minimum size=22pt,inner sep=0]{\popcheck{white}};%
    \hspace{9pt}\begin{minipage}[c]{0.62\linewidth}
      {\popxbold\fontsize{12}{13}\selectfont\color{popcream}Signé\ \textbullet\ {\poplogo OnBuch\color{poporange}+}}\par\vspace{2pt}
      {\raggedright\popsemi\fontsize{8}{10.5}\selectfont\color{popline}\MakeUppercase{Équipe pédagogique OnBuch+}\\\MakeUppercase{Conforme au programme officiel MINESEC}\par}
    \end{minipage}\hfill
    {\poplogo\fontsize{22}{22}\selectfont\color{popcream}OnBuch\color{poporange}+}
  \end{tcolorbox}%
}

% ============================================================
% Page de garde
% #1 titre · #2 matière · #3 niveau · #4 module · #5 n° leçon · #6 durée ·
% #7 description / objectifs (texte ou itemize)
% ============================================================
\newcommand{\pagegarde}[7]{%
  \thispagestyle{empty}
  \newgeometry{a4paper,margin=1.7cm,top=1.6cm,bottom=1.4cm}%
  \begin{tikzpicture}[remember picture,overlay]
    \fill[poporange!18] ([xshift=-3.2cm,yshift=-3.6cm]current page.north east) circle (4.6cm);
    \fill[poppurple!14] ([xshift=2.4cm,yshift=7.6cm]current page.south west) circle (3.8cm);
  \end{tikzpicture}%
  {\poplogo\fontsize{30}{30}\selectfont\color{popink}OnBuch\color{poporange}+}\hfill
  \raisebox{0.6ex}{\poppill{Cours signé \textbullet\ Premium}{popink}{popcream}}\par
  \vspace{1pt}\popsquiggle{poppurple}\par
  \vspace{8pt}
  \begin{tcolorbox}[enhanced,colback=poporange,colframe=popink,boxrule=2.8pt,arc=13pt,
      drop shadow=popink,left=18pt,right=18pt,top=12pt,bottom=13pt,after skip=10pt]
    \poppill{#2 \textbullet\ #3}{popink}{popcream}\hspace{5pt}\poppill{#4}{white}{popink}\par\vspace{8pt}
    {\popdisplay\fontsize{14}{14}\selectfont\color{popink}LEÇON #5}\par\vspace{4pt}
    {\raggedright\hyphenpenalty=10000\exhyphenpenalty=10000\popdisplay\fontsize{25}{27}\selectfont\color{popcream}\MakeUppercase{#1}\par}
    \vspace{8pt}
    {\popxbold\fontsize{9.5}{9.5}\selectfont\color{popink}\MakeUppercase{Durée conseillée : #6}}
  \end{tcolorbox}
  \begin{tcolorbox}[enhanced,colback=white,colframe=popink,boxrule=2.2pt,arc=10pt,
      drop shadow=popink,left=16pt,right=16pt,top=10pt,bottom=10pt,after skip=8pt]
    \poppill{Dans cette leçon}{poppurple}{white}\par\vspace{5pt}
    \popsemi\fontsize{9.3}{12.6}\selectfont\color{popink}#7
  \end{tcolorbox}
  \noindent\begin{minipage}[t]{0.487\linewidth}
    \begin{tcolorbox}[enhanced,colback=popgreen,colframe=popink,boxrule=2.2pt,arc=10pt,
        drop shadow=popink,left=11pt,right=11pt,top=10pt,bottom=10pt,height=3.1cm,valign=center]
      \tcbox[colback=white,colframe=popink,boxrule=1.3pt,arc=5pt,boxsep=0pt,nobeforeafter,
        left=3pt,right=3pt,top=3pt,bottom=3pt,valign=center]{\qrcode[height=1.9cm]{https://play.google.com/store/apps/details?id=com.luvvix.onbuch&hl=fr}}%
      \hspace{8pt}\begin{minipage}[c]{0.5\linewidth}
        {\popdisplay\fontsize{11}{12.5}\selectfont\color{white}\MakeUppercase{Télécharge OnBuch}}\par\vspace{4pt}
        {\raggedright\popsemi\fontsize{8.4}{11}\selectfont\color{white}Gratuit \textbullet\ Google Play\\Tuteur IA, annales, concours, résultats\par}
      \end{minipage}
    \end{tcolorbox}
  \end{minipage}\hfill
  \begin{minipage}[t]{0.487\linewidth}
    \begin{tcolorbox}[enhanced,colback=poppurple,colframe=popink,boxrule=2.2pt,arc=10pt,
        drop shadow=popink,left=11pt,right=11pt,top=10pt,bottom=10pt,height=3.1cm,valign=center]
      \tikz[baseline=-0.7ex]\node[circle,fill=white,draw=popink,line width=1.3pt,minimum size=1.6cm,inner sep=0]{\popdisplay\fontsize{24}{24}\selectfont\color{poppurple}+};%
      \hspace{8pt}\begin{minipage}[c]{0.6\linewidth}
        {\popdisplay\fontsize{11}{12.5}\selectfont\color{white}\MakeUppercase{Passe à OnBuch+}}\par\vspace{4pt}
        {\raggedright\popsemi\fontsize{8.4}{11}\selectfont\color{white}Tous les cours signés, Tuteur IA illimité, fiches PDF — onglet OnBuch+ de l'app\par}
      \end{minipage}
    \end{tcolorbox}
  \end{minipage}
  \vspace{8pt}
  \popcontenu
  \vfill
  \signatureonbuch
  \restoregeometry
  \newpage
}

% Rangée « Ce cours contient » (page de garde)
\newcommand{\poptile}[3]{% #1 couleur #2 gros texte #3 légende
  \begin{tcolorbox}[enhanced,colback=#1,colframe=popink,boxrule=2pt,arc=9pt,shadow={2.5pt}{-2.5pt}{0pt}{popink},
      left=6pt,right=6pt,top=9pt,bottom=9pt,halign=center,valign=center,height=1.75cm,width=\linewidth,nobeforeafter]
    {\popdisplay\fontsize{12.5}{13}\selectfont\color{white}\MakeUppercase{#2}}\par\vspace{3pt}
    {\centering\popsemi\fontsize{7.8}{9.5}\selectfont\color{white}#3\par}
  \end{tcolorbox}}
\newcommand{\popcontenu}{%
  \par\noindent{\popxboldls\fontsize{7.5}{7.5}\selectfont\color{popmuted}CE COURS SIGNÉ CONTIENT}\par\vspace{7pt}
  \noindent\begin{minipage}[t]{0.235\linewidth}\poptile{popblue}{Cours}{complet, illustré, pas à pas}\end{minipage}\hfill
  \begin{minipage}[t]{0.235\linewidth}\poptile{poporange}{Méthodes}{et exercices résolus}\end{minipage}\hfill
  \begin{minipage}[t]{0.235\linewidth}\poptile{poppink}{Exercices}{type examen + corrigés}\end{minipage}\hfill
  \begin{minipage}[t]{0.235\linewidth}\poptile{popgreen}{Fiche bilan}{l'essentiel à retenir}\end{minipage}\par}

% Sommaire
\newtcolorbox{sommaire}{enhanced,colback=white,colframe=popink,boxrule=2.2pt,arc=10pt,
  drop shadow=popink,left=18pt,right=18pt,top=13pt,bottom=13pt,before skip=6pt,after skip=20pt,
  before upper={\poppill{Sommaire}{poppurple}{white}\par\vspace{9pt}\popsemi\fontsize{10.6}{14.5}\selectfont\color{popink}}}

% Signature de fin de cours — poussée tout en bas de la dernière page
\newcommand{\findecours}{%
  \par\vfill\nopagebreak
  \begin{center}\popsquiggle{poporange}\end{center}\vspace{4pt}
  \signatureonbuch
}
\AtBeginDocument{\def\llbracket{\mathopen{[\mkern-3.5mu[}}\def\rrbracket{\mathclose{]\mkern-3.5mu]}}} % intervalles d'entiers (glyphe ⟦ absent de la police)
% Glyphes absents de Fira Math : redéfinis à partir de glyphes disponibles
\AtBeginDocument{%
  \def\setminus{\mathbin{\backslash}}\let\smallsetminus\setminus
  \def\vdots{\mathord{\vbox{\baselineskip3.2pt\lineskiplimit0pt\kern4pt\hbox{.}\hbox{.}\hbox{.}}}}%
  \def\ddots{\mathinner{\mkern1mu\raise6.5pt\hbox{.}\mkern2mu\raise3.6pt\hbox{.}\mkern2mu\raise0.7pt\hbox{.}\mkern1mu}}%
  \def\triangle{\mathord{\scalebox{0.9}{$\symup{\Delta}$}}}%
  \def\nabla{\mathord{\raisebox{\depth}{\scalebox{1}[-1]{$\symup{\Delta}$}}}}%
  \def\checkmark{\popcheck{popdark}}%
  \def\star{\mathbin{\ast}}\def\bigstar{\mathord{\ast}}%
  \def\diamond{\mathbin{\scalebox{0.6}{$\Diamond$}}}%
  \def\bigcup{\mathop{\mathchoice{\raisebox{-0.3em}{\scalebox{1.7}{$\cup$}}}{\scalebox{1.25}{$\cup$}}{\cup}{\cup}}\displaylimits}%
  \def\bigcap{\mathop{\mathchoice{\raisebox{-0.3em}{\scalebox{1.7}{$\cap$}}}{\scalebox{1.25}{$\cap$}}{\cap}{\cap}}\displaylimits}%
  \def\lesseqqgtr{\mathrel{\lessgtr}}\def\bigoplus{\mathop{\oplus}\displaylimits}%
}
\providecommand{\ssi}{si et seulement si} % abréviation fréquente
\AtBeginDocument{\providecommand{\wideparen}{\overparen}} % arc de cercle

%% FIGFIX-BEGIN v5 — correctifs globaux des figures (docs/onbuch-generation/tools/figfix)
\usepackage{adjustbox}\usepackage{etoolbox}\tcbuselibrary{raster}
% toute figure TikZ/pgfplots plus large que la ligne est réduite à la largeur disponible
\BeforeBeginEnvironment{tikzpicture}{\begin{adjustbox}{max width=\linewidth}}
\AfterEndEnvironment{tikzpicture}{\end{adjustbox}}
% légendes pgfplots lisibles par-dessus les courbes
\pgfplotsset{every axis/.append style={legend style={fill=white,fill opacity=0.92,text opacity=1,draw=black!15,inner sep=3pt}}}
% étiquettes d'arêtes (syntaxe edge["..."]) avec fond blanc
\tikzset{every edge quotes/.append style={fill=white,inner sep=1.2pt}}
%% FIGFIX-END
\begin{document}

\pagegarde{Réduction des conséquences néfastes des activités humaines sur les ressources naturelles}{SVT}{1ère D}{Module 1 : Le Monde Vivant}{ND07}{12 h}{Cette leçon présente les principales ressources naturelles du Cameroun (eau, sols, forêts, biodiversité), les pressions que les activités humaines exercent sur elles, et les mesures de gestion durable permettant de les préserver. Elle introduit la notion de développement durable et s'appuie sur des exemples camerounais concrets (forêts communautaires, parcs nationaux, lutte contre la désertification).
\vspace{4pt}
\begin{multicols}{2}\small
\begin{itemize}
\item À la fin de cette leçon, je sais définir une ressource naturelle et distinguer ressources renouvelables et non renouvelables
\item À la fin de cette leçon, je sais citer les principales ressources naturelles du Cameroun et leur répartition
\item À la fin de cette leçon, je sais identifier, à partir d'un document, les impacts d'une activité humaine donnée sur une ressource naturelle
\item À la fin de cette leçon, je sais définir le développement durable et ses trois piliers
\item À la fin de cette leçon, je sais expliquer le principe de la gestion raisonnée (gestion durable) d'une ressource
\item À la fin de cette leçon, je sais proposer des mesures de gestion durable adaptées à une situation locale camerounaise
\end{itemize}\end{multicols}}

\begin{sommaire}
\begin{multicols}{2}
\begin{enumerate}
\item Les ressources naturelles : définition et typologie
\item Les pressions des activités humaines sur les ressources
\item Le développement durable : notion et principes
\item La gestion durable des ressources au Cameroun
\item Analyse de dossiers documentaires et synthèse
\item Activité d'intégration
\item Méthodes et exercices résolus
\item Exercices
\item Corrigés des exercices
\item Fiche bilan
\end{enumerate}
\end{multicols}
\end{sommaire}

\begin{objectifs}
\begin{itemize}
\item À la fin de cette leçon, je sais définir une ressource naturelle et distinguer ressources renouvelables et non renouvelables
\item À la fin de cette leçon, je sais citer les principales ressources naturelles du Cameroun et leur répartition
\item À la fin de cette leçon, je sais identifier, à partir d'un document, les impacts d'une activité humaine donnée sur une ressource naturelle
\item À la fin de cette leçon, je sais définir le développement durable et ses trois piliers
\item À la fin de cette leçon, je sais expliquer le principe de la gestion raisonnée (gestion durable) d'une ressource
\item À la fin de cette leçon, je sais proposer des mesures de gestion durable adaptées à une situation locale camerounaise
\item À la fin de cette leçon, je sais analyser un dossier documentaire portant sur l'exploitation d'une ressource naturelle au Cameroun
\item À la fin de cette leçon, je sais citer des exemples camerounais de conservation (aires protégées, forêts communautaires, reboisement)
\end{itemize}
\end{objectifs}

\begin{prerequis}
\begin{itemize}
\item Notion d'écosystème et de chaîne alimentaire (classes antérieures)
\item Notion de biodiversité et de classification du vivant (début de Première)
\item Relations entre l'Homme et son environnement (Seconde / enseignements antérieurs)
\item Lecture et exploitation de documents (graphiques, tableaux, cartes)
\end{itemize}
\end{prerequis}

\newpage

\coursec{Les ressources naturelles : définition et typologie}

\textbf{Situation d'accroche.}\par Dans la région de l'Extrême-Nord du Cameroun, le lac Tchad a perdu plus de 90 \% de sa superficie en cinquante ans, privant des milliers de pêcheurs de leur ressource. Dans le Sud, la forêt du bassin du Congo recule chaque année sous l'effet de l'exploitation forestière et de l'agriculture itinérante. Pourtant, ces ressources — eau, sols, forêts, biodiversité — sont la base de la vie et de l'économie camerounaises. Comment l'Homme dégrade-t-il ces ressources ? Et surtout, comment peut-on continuer à les exploiter sans les épuiser ? C'est tout l'enjeu du développement durable que cette leçon va explorer.

\textbf{Problématique.}\par Qu'appelle-t-on exactement une ressource naturelle ? Toutes les ressources se comportent-elles de la même façon face à l'exploitation humaine ? Pour répondre à ces questions, nous devons d'abord définir précisément la notion de ressource naturelle, puis établir une classification rigoureuse qui nous servira de fil conducteur dans toute la leçon.

\courssub{Qu'est-ce qu'une ressource naturelle ?}

\textbf{Observation.}\par Un pêcheur de Logone-Birni prélève des poissons dans le fleuve Logone. Une paysanne de l'Ouest laboure son champ et y sème le maïs. Une entreprise forestière de l'Est abat des ayous et des sapellis pour en faire du bois d'œuvre. Une compagnie pétrolière extrait du pétrole au large de Kribi. Qu'ont en commun toutes ces activités ?

Dans chaque cas, l'Homme \cle{prélève} un élément de la \cle{nature} (poisson, sol fertile, arbre, pétrole) et l'\cle{utilise} pour satisfaire un \cle{besoin} : se nourrir, se loger, se chauffer, produire des biens, gagner un revenu. C'est exactement ce qui définit une ressource naturelle.

\begin{definition}[Ressource naturelle]
Une \cle{ressource naturelle} est tout élément de la nature (minéral, végétal, animal, hydrique, énergétique) que l'Homme prélève dans son environnement et utilise pour satisfaire ses besoins vitaux (se nourrir, boire, se loger) ou économiques (production, commerce, industrie).
\end{definition}

\begin{exemplebox}[Des ressources naturelles dans ta vie quotidienne]
Sans le savoir, tu utilises chaque jour des dizaines de ressources naturelles :
\begin{itemize}
\item l'\textbf{eau} de la borne-fontaine ou du forage (ressource hydrique) ;
\item le \textbf{bois} de chauffe ou le charbon de bois pour cuisiner (ressource forestière) ;
\item le \textbf{sol} qui a produit le manioc, l'igname ou le maïs de ton repas (ressource pédologique) ;
\item le \textbf{pétrole} transformé en carburant pour les motos-taxis et les cars (ressource minérale et énergétique) ;
\item la \textbf{biodiversité} : poissons du marché, plantes médicinales comme le prunier d'Afrique (\emph{Prunus africana}) du Mont Oku.
\end{itemize}
\end{exemplebox}

\begin{attention}[Ne pas confondre ressource naturelle et produit manufacturé]
Une ressource naturelle est un élément \emph{prélevé directement dans la nature}. L'eau du forage est une ressource naturelle ; l'eau minérale embouteillée est un \emph{produit transformé}. Le bois d'œuvre est une ressource ; la chaise fabriquée avec ce bois est un bien manufacturé. Au Probatoire, on te demandera souvent de distinguer, dans un document, la ressource naturelle de ses produits dérivés.
\end{attention}

\courssub{Ressources renouvelables et non renouvelables : la grande distinction}

\textbf{Une expérience de pensée.}\par Imagine deux greniers. Le premier contient un stock de mil qui ne se renouvelle jamais : chaque sac pris diminue définitivement le stock. Le second est rempli chaque année par la récolte : on peut en prélever chaque année, à condition de ne pas prendre plus que ce que la récolte rapporte. Le premier grenier illustre les ressources \cle{non renouvelables} ; le second, les ressources \cle{renouvelables}.

\begin{definition}[Ressource renouvelable]
Une \cle{ressource renouvelable} est une ressource naturelle capable de se \cle{régénérer} à l'échelle d'une vie humaine (quelques années à quelques décennies), à condition que son exploitation reste modérée. Exemples : l'eau, les forêts, les sols, la biodiversité, les stocks de poissons.
\end{definition}

\begin{definition}[Ressource non renouvelable]
Une \cle{ressource non renouvelable} (ou \emph{épuisable}) est une ressource dont la formation dans la nature demande des durées géologiques (des millions d'années), très supérieures à l'échelle humaine. Tout prélèvement diminue donc irréversiblement le stock disponible. Exemples : le pétrole, le gaz naturel, les minerais (fer, bauxite, or, cobalt).
\end{definition}

\textbf{La condition de renouvelabilité : un point capital.}\par Dire qu'une ressource est renouvelable ne signifie pas qu'elle est inépuisable. Une ressource renouvelable ne se perpétue que si une condition stricte est respectée :

\begin{aretenir}[La règle d'or de la renouvelabilité]
Une ressource renouvelable reste disponible durablement si et seulement si :
\[ \text{rythme d'exploitation} \leq \text{rythme de régénération} \]
Si l'Homme prélève plus vite que la nature ne renouvelle, la ressource s'épuise \emph{localement}, même si elle est renouvelable en théorie.
\end{aretenir}

\begin{attention}[Erreur fréquente au Probatoire]
Beaucoup d'élèves écrivent : « l'eau et la forêt sont renouvelables, donc on ne peut pas les épuiser ». C'est \textbf{faux}. L'eau et la forêt sont renouvelables \emph{à condition} que l'exploitation ne dépasse pas la régénération. Une nappe phréatique pompée plus vite qu'elle ne se recharge s'épuise ; une forêt coupée plus vite qu'elle ne repousse disparaît. Le mot-clé à utiliser est : renouvelable \emph{mais vulnérable}.
\end{attention}

\begin{exoresolu}[Classer des ressources camerounaises]
Énoncé. On prélève chaque année au Cameroun : (a) du pétrole au large de Kribi ; (b) du bois dans la forêt de l'Est ; (c) de l'eau du fleuve Sanaga pour le barrage de Song-Loulou ; (d) de la bauxite à Minim-Martap ; (e) des poissons dans le lac Tchad. Pour chaque cas, indique si la ressource est renouvelable ou non renouvelable, en justifiant.
\tcblower
Solution détaillée.
\begin{itemize}
\item (a) \textbf{Pétrole} : non renouvelable. Sa formation a demandé des millions d'années (décomposition de plancton marin enfoui) ; tout baril extrait diminue définitivement le stock.
\item (b) \textbf{Bois de la forêt de l'Est} : renouvelable, car les arbres repoussent en quelques décennies, \emph{mais seulement} si le nombre d'arbres coupés ne dépasse pas la croissance naturelle (ou le reboisement).
\item (c) \textbf{Eau de la Sanaga} : renouvelable, grâce au cycle de l'eau (évaporation, précipitations) qui recharge le fleuve chaque saison des pluies.
\item (d) \textbf{Bauxite} : non renouvelable. C'est un minerai formé par l'altération des roches sur des millions d'années ; le gisement ne se reforme pas à l'échelle humaine.
\item (e) \textbf{Poissons du lac Tchad} : renouvelable (les poissons se reproduisent), mais vulnérable : la surpêche et le rétrécissement du lac ont fait chuter les captures, montrant que le rythme de prélèvement a dépassé le rythme de régénération.
\end{itemize}
\end{exoresolu}

L'arbre de classification ci-dessous résume cette typologie, avec des exemples camerounais.

\begin{popfigure}
\begin{center}
\begin{tikzpicture}[
  grow=right,
  level 1/.style={sibling distance=3.6cm, level distance=4.2cm},
  level 2/.style={sibling distance=1.05cm, level distance=4.6cm},
  every node/.style={font=\small},
  racine/.style={draw=popdark, fill=popblueL, rounded corners, thick, align=center, font=\small\bfseries},
  cat/.style={draw=popblue, fill=popblue!15, rounded corners, thick, align=center, font=\small\bfseries},
  cat2/.style={draw=poporange, fill=poporangeL, rounded corners, thick, align=center, font=\small\bfseries},
  ex/.style={draw=popmuted, fill=popcream, rounded corners, align=center, font=\footnotesize},
  edge from parent/.style={draw=popmuted, thick}
]
\node[racine, align=center]{Ressources\\naturelles}
  child { node[cat2, align=center]{Non renouvelables\\(formation : millions\\d'années)}
    child { node[ex, align=center]{Minerais : bauxite\\(Minim-Martap),\\fer (Mbalam), or (Est)} }
    child { node[ex, align=center]{Pétrole et gaz\\(bassin du Rio del Rey,\\offshore de Kribi)} }
  }
  child { node[cat, align=center]{Renouvelables\\(régénération à\\l'échelle humaine)}
    child { node[ex, align=center]{Biodiversité : faune,\\flore, stocks de poissons\\(littoral, lacs)} }
    child { node[ex, align=center]{Forêts : bassin du Congo\\(Sud, Est, Centre)} }
    child { node[ex, align=center]{Sols : terres agricoles\\(Ouest, Adamaoua)} }
    child { node[ex, align=center]{Eau : fleuves, lacs,\\nappes souterraines} }
  };
\end{tikzpicture}
\end{center}
\legende{Arbre de classification des ressources naturelles. Les ressources renouvelables (à gauche en bleu) se régénèrent à l'échelle d'une vie humaine si l'exploitation reste modérée ; les ressources non renouvelables (en orange) s'épuisent définitivement à chaque prélèvement.}
\end{popfigure}

\courssub{Les quatre grandes ressources étudiées dans cette leçon}

Le programme de Première D retient quatre ressources naturelles majeures, qui seront le fil conducteur de toute la leçon.

\textbf{1. L'eau.}\par L'eau est indispensable à la vie (boisson, hygiène), à l'agriculture (irrigation), à l'industrie et à la production d'électricité. Le Cameroun, surnommé « l'Afrique en miniature », est aussi un véritable « château d'eau » : ses fleuves (Sanaga, Wouri, Logone, Bénoué) alimentent les barrages hydroélectriques de Song-Loulou, Édéa et Lagdo, qui fournissent l'essentiel de l'électricité du pays. Pourtant, l'accès à l'eau potable reste inégal, et certaines nappes sont surexploitées ou polluées.

\textbf{2. Les sols.}\par Le sol est une ressource formée très lentement : il faut environ \num{100} à \num{1000} ans pour former un centimètre de sol fertile par altération de la roche et accumulation d'humus. Les sols fertiles des hautes terres de l'Ouest et de l'Adamaoua supportent l'agriculture vivrière (maïs, tubercules) et les cultures de rente (café, cacao). L'érosion, la déforestation et les cultures intensives sans jachère appauvrissent ces sols.

\textbf{3. Les forêts.}\par La forêt fournit du bois, des produits alimentaires (gibier, fruits, chenilles), des plantes médicinales, et joue un rôle climatique majeur : elle stocke le carbone atmosphérique et régule les pluies.

\begin{savaistu}[Le Cameroun, deuxième massif forestier d'Afrique]
Le Cameroun possède le \cle{deuxième massif forestier d'Afrique} après la République démocratique du Congo, avec environ \textbf{22 millions d'hectares} de forêts, soit près de 46 \% du territoire national. Ces forêts appartiennent au bassin du Congo, le deuxième poumon vert de la planète après l'Amazonie. Elles abritent des espèces emblématiques : gorilles, chimpanzés, éléphants de forêt, et des arbres précieux comme l'ayous, le sapelli et l'ébène.
\end{savaistu}

\textbf{4. La biodiversité.}\par La biodiversité désigne l'ensemble des espèces vivantes (animales, végétales, microbiennes) et de leurs milieux. Grâce à la diversité de ses reliefs et de ses climats, le Cameroun est l'un des pays les plus riches d'Afrique en biodiversité : plus de \num{9000} espèces de plantes, environ \num{900} espèces d'oiseaux, plus de \num{300} espèces de mammifères. Cette biodiversité est une ressource alimentaire (pêche, chasse), médicinale et touristique (parcs nationaux de Waza, de la Bénoué, de Korup).

\begin{exemplebox}[Le lac Tchad : une ressource renouvelable qui s'effondre]
Le lac Tchad, partagé entre le Cameroun, le Tchad, le Niger et le Nigeria, illustre parfaitement la condition de renouvelabilité. En 1960, il couvrait environ \SI{25000}{km^2} ; aujourd'hui, il n'en couvre plus que \SI{2500}{km^2} environ, soit une perte de plus de 90 \% de sa superficie. Les causes sont combinées : sécheresses répétées (diminution de la régénération naturelle) et prélèvements croissants pour l'irrigation (augmentation de l'exploitation). Résultat : les stocks de poissons se sont effondrés, les pêcheurs camerounais de la région de l'Extrême-Nord ont perdu leur ressource, et des conflits pour l'accès à l'eau et aux pâturages se multiplient. Une ressource \emph{renouvelable} a été \emph{épuisée localement} parce que le prélèvement a durablement dépassé la régénération.
\end{exemplebox}

\courssub{La répartition régionale des ressources naturelles au Cameroun}

Le Cameroun présente une remarquable diversité de milieux, du fait de son étirement latitudinal (du golfe de Guinée au lac Tchad) et de ses reliefs. On distingue trois grands ensembles :

\begin{itemize}
\item \textbf{Le Sud forestier} (régions du Sud, de l'Est, du Centre et du Littoral) : climat équatorial humide, grande forêt dense du bassin du Congo, sols fertiles sous couvert forestier, et au large, les gisements de pétrole du bassin du Rio del Rey et de Kribi ;
\item \textbf{Le Centre et l'Ouest} (Adamaoua, Ouest, Nord-Ouest) : hautes terres aux sols volcaniques très fertiles, zones d'agriculture intensive et d'élevage ;
\item \textbf{Le Nord soudano-sahélien} (Nord et Extrême-Nord) : savanes, zones d'élevage bovin, cultures de coton et de mil, ressources en eau fragiles (lac Tchad, fleuve Logone, barrage de Lagdo sur la Bénoué) ;
\item \textbf{Le littoral} : ressources halieutiques (pêche maritime au large de Douala et de Limbe), mangroves de l'estuaire du Wouri, pétrole offshore.
\end{itemize}

\begin{popfigure}
\begin{center}
\begin{tikzpicture}[scale=1, font=\footnotesize]
% Contour simplifié du Cameroun
\draw[thick, popdark, fill=popcream]
  (1.2,0) -- (2.6,0.3) -- (3.6,0.1) -- (4.6,0.6) -- (5.2,1.6) -- (5.0,2.8)
  -- (4.4,3.6) -- (4.2,4.6) -- (3.4,5.4) -- (3.0,6.4) -- (3.4,7.4) -- (3.0,8.4)
  -- (2.2,8.0) -- (1.8,7.0) -- (1.4,6.0) -- (1.0,5.0) -- (0.6,4.0) -- (0.4,3.0)
  -- (0.2,2.0) -- (0.6,1.0) -- cycle;
% Zones climatiques / ressources
\fill[popgreen!35, rounded corners] (0.6,0.6) -- (2.6,0.5) -- (4.4,0.9) -- (4.9,1.8) -- (4.6,2.9) -- (0.7,2.9) -- cycle;
\fill[popgold!30, rounded corners] (0.7,3.0) -- (4.6,3.0) -- (4.1,4.5) -- (1.2,4.6) -- cycle;
\fill[poporange!25, rounded corners] (1.2,4.7) -- (4.1,4.6) -- (3.3,5.5) -- (2.9,6.5) -- (3.3,7.4) -- (2.9,8.3) -- (2.2,7.9) -- (1.5,6.5) -- (1.1,5.5) -- cycle;
% Étiquettes des zones
\node[align=center, font=\footnotesize\bfseries, text=popgreen!60!black] at (2.6,1.7) {Grande forêt dense\\du bassin du Congo\\(bois, biodiversité)};
\node[align=center, font=\footnotesize\bfseries, text=popgold!60!black] at (2.6,3.8) {Hautes terres : sols\\fertiles, agriculture\\(café, cacao, maïs)};
\node[align=center, font=\footnotesize\bfseries, text=poporange!70!black] at (2.3,6.2) {Savanes : élevage,\\coton, mil};
% Lac Tchad
\fill[popblue!50] (3.0,8.6) ellipse (0.55 and 0.3);
\node[font=\footnotesize\bfseries, text=popblue!70!black] at (4.6,8.7) {Lac Tchad};
\draw[thin, popblue!70!black] (3.55,8.65) -- (4.05,8.7);
% Pétrole offshore
\fill[popdark] (0.0,1.2) circle (0.09);
\fill[popdark] (-0.15,1.9) circle (0.09);
\node[align=left, font=\footnotesize\bfseries, text=popdark] at (-2.2,1.6) {Pétrole offshore\\(Rio del Rey, Kribi)};
\draw[thin, popdark] (-0.15,1.9) -- (-0.9,1.8);
\draw[thin, popdark] (0.0,1.2) -- (-0.9,1.4);
% Pêche littorale
\node[align=left, font=\footnotesize\bfseries, text=popblue!70!black] at (-2.2,0.3) {Pêche maritime\\(estuaire du Wouri)};
\draw[thin, popblue!70!black] (0.2,0.7) -- (-0.9,0.45);
% Bauxite
\fill[poppink!80!black] (3.6,4.9) rectangle (3.8,5.1);
\node[align=left, font=\footnotesize\bfseries, text=poppink!80!black] at (5.9,5.0) {Bauxite\\(Minim-Martap)};
\draw[thin, poppink!80!black] (3.8,5.0) -- (4.9,5.0);
% Boussole
\draw[->, thick, popmuted] (5.6,7.6) -- (5.6,8.4) node[above, font=\footnotesize\bfseries]{N};
\end{tikzpicture}
\end{center}
\legende{Carte simplifiée du Cameroun localisant les grandes ressources naturelles : forêt dense du bassin du Congo au Sud (vert), sols agricoles des hautes terres (jaune), savanes et élevage au Nord (orange), lac Tchad à l'extrême Nord, pétrole offshore et pêche sur le littoral, bauxite de Minim-Martap sur l'Adamaoua.}
\end{popfigure}

\courssub{Complément pour le Probatoire : capital naturel et services écosystémiques}

\begin{definition}[Capital naturel et services écosystémiques]
Le \cle{capital naturel} désigne l'ensemble des ressources naturelles et des écosystèmes d'un territoire, considérés comme un « stock de richesse » au même titre qu'un capital économique. Les \cle{services écosystémiques} sont les bénéfices que les écosystèmes rendent gratuitement à l'Humanité. On distingue :
\begin{itemize}
\item les \textbf{services d'approvisionnement} : eau, bois, poissons, fibres, plantes médicinales ;
\item les \textbf{services de régulation} : stockage du carbone par les forêts, régulation du climat et des crues, purification de l'eau, pollinisation des cultures ;
\item les \textbf{services culturels} : tourisme (parcs de Waza et de Korup), valeurs spirituelles et récréatives ;
\item les \textbf{services de soutien} : formation des sols, cycle des nutriments, photosynthèse.
\end{itemize}
\end{definition}

\begin{exemplebox}[La mangrove de l'estuaire du Wouri : un capital naturel méconnu]
La mangrove au sud de Douala rend simultanément plusieurs services : elle fournit du bois et des poissons (approvisionnement), protège les côtes contre l'érosion et les tempêtes (régulation), sert de nurserie aux poissons du golfe de Guinée (soutien). Quand on détruit la mangrove pour construire ou produire du charbon, on ne perd pas seulement des arbres : on perd \emph{tous ces services gratuits}, qu'il faudrait alors remplacer par des infrastructures coûteuses (digues, élevages piscicoles).
\end{exemplebox}

\begin{aretenir}[L'essentiel de la section]
\begin{itemize}
\item Une \cle{ressource naturelle} est un élément de la nature prélevé et utilisé par l'Homme pour satisfaire ses besoins.
\item Les ressources \cle{renouvelables} (eau, sols, forêts, biodiversité) se régénèrent à l'échelle humaine ; les ressources \cle{non renouvelables} (pétrole, minerais) s'épuisent définitivement.
\item Une ressource renouvelable ne le reste que si le \cle{rythme d'exploitation} ne dépasse pas le \cle{rythme de régénération}.
\item Le Cameroun possède un capital naturel exceptionnel mais inégalement réparti : forêts au Sud, sols agricoles au Centre-Ouest, savanes et élevage au Nord, pêche et pétrole sur le littoral.
\end{itemize}
\end{aretenir}

Maintenant que nous savons ce que sont les ressources naturelles et comment les classer, la section suivante examinera les \emph{pressions} que les activités humaines exercent sur elles.

\coursec{Les pressions des activités humaines sur les ressources}

Dans la section précédente, nous avons défini les \cle{ressources naturelles} (eau, sols, forêts, biodiversité) et établi leur typologie. On distingue les ressources \textbf{renouvelables} (eau, forêts, espèces vivantes) qui peuvent se régénérer à l'échelle humaine \emph{si} leur exploitation respecte leur rythme de renouvellement, et les ressources \textbf{non renouvelables} (minerais, pétrole, gaz) dont le stock est fini. Le sol occupe une position intermédiaire : sa formation (pédogenèse) est si lente (siècles à millénaires pour quelques centimètres) qu'il est considéré comme \textbf{non renouvelable à l'échelle humaine}, même s'il peut être entretenu. Or, partout au Cameroun comme dans le monde, les \cle{activités humaines} exercent sur ces ressources des \cle{pressions} croissantes : prélèvements, rejets, destructions. C'est l'objet de cette section : identifier, ressource par ressource, les pressions exercées et comprendre leurs mécanismes, afin de pouvoir ensuite (sections 3 et 4) proposer des mesures de gestion durable.

\courssub{Notion de pression : de l'activité humaine à la dégradation de la ressource}

\begin{definition}[Pression sur une ressource naturelle]
On appelle \cle{pression} toute action humaine qui modifie l'état d'une ressource naturelle, soit en \textbf{quantité} (prélèvements, destructions), soit en \textbf{qualité} (pollutions, dégradations). Une pression devient problématique lorsqu'elle dépasse la \cle{capacité de régénération} de la ressource (pour les ressources renouvelables) ou la \cle{capacité d'autoépuration} du milieu (pour les pollutions).
\end{definition}

Prenons une situation concrète. À Bafoussam, un abattage anarchique de la forêt communautaire sur les versants est suivi, quelques années plus tard, de glissements de terrain et de tarissement des sources qui alimentaient les quartiers. Comment relier ces faits ? Il faut raisonner en \cle{chaîne de causalité} : une activité humaine provoque un impact direct, qui entraîne des impacts indirects, qui aboutissent à la dégradation de la ressource.

\begin{popfigure}
\begin{center}
\begin{tikzpicture}[
  node distance=0.55cm,
  act/.style={rectangle, rounded corners=3pt, draw=poporange, fill=poporangeL, text width=3.1cm, align=center, font=\small\bfseries},
  imp/.style={rectangle, rounded corners=3pt, draw=popblue, fill=popblueL, text width=3.1cm, align=center, font=\small},
  cons/.style={rectangle, rounded corners=3pt, draw=popink, fill=poppinkL, text width=3.1cm, align=center, font=\small\bfseries},
  fleche/.style={-{Stealth[length=3mm]}, thick, popdark}
]
\node[act, align=center] (a){Activité humaine\\ déforestation sur les versants};
\node[imp, below=of a, align=center] (b){Impact direct\\ disparition du couvert végétal};
\node[imp, below=of b, align=center] (c){Impacts indirects\\ érosion des sols, ruissellement, envasement des cours d'eau};
\node[cons, below=of c, align=center] (d){Conséquence sur les ressources\\ sols appauvris, eaux turbides, sources taries};
\draw[fleche] (a) -- (b);
\draw[fleche] (b) -- (c);
\draw[fleche] (c) -- (d);
\end{tikzpicture}
\end{center}
\legende{Chaîne de causalité reliant une activité humaine (déforestation) à la dégradation des ressources en sol et en eau. Cette démarche en chaîne est l'outil de base de l'analyse d'impact.}
\end{popfigure}

\begin{attention}[Pollution ou surexploitation ? Ne pas confondre]
Erreur très fréquente au Probatoire : confondre \cle{pollution} et \cle{surexploitation}.
\begin{itemize}
  \item La \textbf{pollution} est une dégradation de la \textbf{qualité} de la ressource (eau rendue impropre à la consommation par des rejets, sol contaminé par des pesticides).
  \item La \textbf{surexploitation} est une dégradation de la \textbf{quantité} de la ressource (nappe phréatique épuisée par les pompages, forêt réduite par les coupes, poissons raréfiés par la surpêche).
\end{itemize}
Une même activité peut combiner les deux : l'irrigation intensive prélève l'eau (surexploitation) et rejette des engrais (pollution). Dans une réponse, précise toujours de quel type d'impact il s'agit.
\end{attention}

\courssub{Les pressions sur la ressource en eau}

\textbf{Situation concrète.} En aval des zones industrielles de Douala-Bonabéri, les riverains constatent que l'eau des puits devient impropre à la boisson, tandis qu'au Nord, le lac Tchad, autrefois immense, a spectaculairement rétréci. Deux phénomènes différents, deux types de pressions.

\textbf{1. La pollution de l'eau.} Elle provient principalement de :
\begin{itemize}
  \item la \cle{pollution domestique} : eaux usées rejetées sans traitement, ordures déversées dans les rigoles et les cours d'eau, déjections humaines et animales. Elles apportent des matières organiques et des micro-organismes pathogènes responsables du choléra, de la typhoïde, des amibiases — maladies encore fréquentes lors des épidémies saisonnières au Cameroun ;
  \item la \cle{pollution industrielle} : rejets d'usines (huiles, métaux lourds, solvants, colorants), souvent déversés sans épuration dans les rivières comme la Wouri ou la Dibamba ;
  \item la \cle{pollution agricole} : engrais (nitrates \ce{NO3-}, phosphates \ce{PO4^3-}) et pesticides lessivés par les pluies vers les nappes et les rivières, provoquant l'eutrophisation des plans d'eau.
\end{itemize}

\textbf{2. Les prélèvements excessifs et l'assèchement.} Lorsque les pompages (irrigation, adduction d'eau des villes en croissance comme Yaoundé ou Maroua) dépassent la recharge des nappes par les pluies, le niveau des nappes phréatiques baisse, les puits se tarissent et les lacs régressent.

\begin{exemplebox}[Le cas du lac Tchad : un assèchement emblématique]
Le lac Tchad, partagé entre le Cameroun, le Tchad, le Niger et le Nigéria, a vu sa superficie d'eau libre permanente passer d'environ \num{25000} km\textsuperscript{2} dans les années 1960 à moins de \num{2000} km\textsuperscript{2} ces dernières décennies (avec de fortes variations interannuelles). Causes combinées : baisse des précipitations (sécheresses sahéliennes des années 1970-1980) et prélèvements massifs pour l'irrigation sur les fleuves qui l'alimentent (Chari, Logone). Conséquences : effondrement de la pêche, perte des pâturages inondables (yaérés), conflits entre agriculteurs, éleveurs et pêcheurs, déplacements de populations. C'est l'exemple-type à mobiliser pour illustrer la \textbf{surexploitation} d'une ressource en eau aggravée par le changement climatique.
\end{exemplebox}

\courssub{Les pressions sur les sols}

\textbf{Observation.} Dans les collines de l'Ouest (Bamiléké), après quelques années de culture sur des parcelles défrichées, les rendements de maïs chutent et les paysans défrichent de nouvelles terres. Dans l'Extrême-Nord, les dunes avancent sur les champs. Hypothèse : les pratiques agricoles et la disparition de la végétation dégradent le sol lui-même.

\textbf{Les mécanismes de dégradation :}
\begin{itemize}
  \item la \cle{déforestation} supprime la protection du sol : sans racines ni litière, l'eau de pluie ruisselle et emporte les particules fines — c'est l'\cle{érosion hydrique} (ravines, gullies) ; dans les zones sèches, le vent emporte les particules — c'est l'\cle{érosion éolienne} (dunes, dépôts éoliens) ;
  \item l'\cle{agriculture intensive} et les cultures répétées sans fumure organique ni rotation exportent les éléments minéraux (azote, phosphore, potassium) plus vite qu'ils ne sont restitués : c'est l'\cle{épuisement de la fertilité} (perte de matière organique, acidification) ;
  \item le \cle{surpâturage} tue la végétation herbacée et tasse le sol sous les sabots du bétail (compaction, perte de porosité) ;
  \item les \cle{feux de brousse} répétés détruisent la matière organique de surface, source d'humus, et tuent la microfaune du sol.
\end{itemize}

\begin{definition}[Désertification]
La \cle{désertification} est la dégradation des terres en \textbf{zone sèche} (aride, semi-aride ou sub-humide sèche) sous l'effet combiné des \textbf{variations climatiques} (sécheresses prolongées) et des \textbf{activités humaines} (déboisement, surpâturage, cultures inadaptées, feux). Elle aboutit à la perte durable de la productivité biologique du sol.
\end{definition}

\begin{attention}[Désertification $\neq$ avancée du désert]
La désertification ne signifie pas que le Sahara « avance » comme une marée de sable. Il s'agit d'une dégradation \emph{locale} des terres, en taches, autour des villages, des puits et des zones de pâturage surexploitées. Au Cameroun, elle menace surtout les départements de l'Extrême-Nord (Mayo-Sava, Mayo-Tsanaga, Diamaré, Logone-et-Chari) et le Nord (Bénoué, Faro, Mayo-Rey).
\end{attention}

\courssub{Les pressions sur les forêts}

Le Cameroun possède l'un des plus grands massifs forestiers d'Afrique centrale (environ 20 millions d'hectares, soit plus de 40 \% du territoire). Pourtant ce patrimoine recule.

\begin{exemplebox}[Un ordre de grandeur à retenir]
Le Cameroun perd chaque année environ \num{100000} à \num{200000} hectares de forêt (soit l'équivalent de plusieurs milliers de terrains de football par jour), selon les estimations des organismes de suivi forestier (Global Forest Watch, MINFOF). Les causes principales : l'agriculture itinérante sur brûlis, l'exploitation forestière (industrielle et artisanale) et l'extension des villes.
\end{exemplebox}

\begin{popfigure}
\begin{center}
\begin{tikzpicture}
\begin{axis}[
  ybar, bar width=0.55cm,
  width=0.85\linewidth, height=5.5cm,
  ymin=0, ymax=26,
  ylabel={Superficie forestière (millions d'ha)},
  symbolic x coords={1990,2000,2010,2020},
  xtick=data,
  nodes near coords,
  every node near coord/.append style={font=\small},
  axis lines=left,
  ymajorgrids=true, grid style={popline, dashed},
]
\addplot[fill=popgreen!70, draw=popgreen] coordinates {(1990,24.3) (2000,22.8) (2010,21.2) (2020,20.0)};
\end{axis}
\end{tikzpicture}
\end{center}
\legende{Évolution indicative de la superficie forestière du Cameroun par décennie (données simplifiées à des fins pédagogiques, sources FAO/MINFOF). La tendance est à une diminution régulière d'environ 1 à 1,5 million d'hectares par décennie.}
\end{popfigure}

\textbf{Les pressions en détail :}
\begin{itemize}
  \item l'\cle{exploitation forestière} (industrielle et artisanale) prélève les essences précieuses (ayous, sapelli, iroko, moabi) ; même « sélective », elle ouvre des pistes qui facilitent ensuite l'installation de champs, de chasseurs et de braconniers ;
  \item l'\cle{agriculture itinérante sur brûlis} : le paysan coupe et brûle une parcelle de forêt, la cultive 2 à 3 ans jusqu'à épuisement, puis en défriche une autre ; avec la croissance démographique, la durée de jachère se raccourcit (parfois < 5 ans) et la forêt n'a plus le temps de se régénérer ;
  \item les \cle{feux de brousse} répétés, surtout dans les savanes de l'Adamaoua et du Nord, qui empêchent la régénération naturelle des lisères forestières ;
  \item l'\cle{extension urbaine} et la demande en bois-énergie (bois de chauffe et charbon de bois, principale source d'énergie domestique de plus de 70 \% des ménages camerounais) ;
  \item l'expansion des \cle{plantations industrielles} (palmier à huile \emph{Elaeis guineensis}, hévéa, bananier) qui convertissent la forêt primaire ou secondaire en monocultures.
\end{itemize}

\begin{savaistu}[Les feux de brousse de l'Adamaoua : un cercle vicieux]
Chaque année, pendant la saison sèche (novembre-mars), d'immenses feux de brousse parcourent les savanes de l'Adamaoua, souvent allumés par des chasseurs pour débusquer le gibier ou par des éleveurs pour « rajeunir » les pâturages (pousse verte post-feu). Or ces feux détruisent la \cle{matière organique} des sols, tuent les jeunes arbres et les micro-organismes du sol, et appauvrissent durablement les pâturages : l'herbe repousse vite, mais sur un sol de moins en moins fertile, plus sensible à l'érosion. C'est un cercle vicieux typique d'une pratique ancienne devenue destructrice à cause de sa fréquence accrue et de la pression démographique.
\end{savaistu}

\courssub{Les pressions sur la biodiversité}

\textbf{Situation concrète.} Sur les marchés de certaines villes (Bertoua, Yokadouma, Douala), on trouve encore de la viande de brousse (pangolin, crocodile, céphalophes, primates) et des écailles de pangolin destinées à l'exportation illicite. Pendant ce temps, les éléphants de la réserve de Waza ou du parc de Bouba Ndjida se raréfient. La biodiversité camerounaise, parmi les plus riches d'Afrique (5\textsuperscript{e} rang continental), est sous pression.

\textbf{Les pressions principales :}
\begin{itemize}
  \item le \cle{braconnage} et le commerce illégal d'espèces sauvages : les \textbf{éléphants} (de forêt \emph{Loxodonta cyclotis} et de savane \emph{Loxodonta africana}) sont tués pour l'ivoire, les \textbf{pangolins} (3 espèces présentes) pour leurs écailles (trafic international vers l'Asie) et leur viande, les \textbf{grands singes} (gorilles de Cross River, gorilles de l'Ouest, chimpanzés) pour la viande de brousse et le trafic de bébés ;
  \item la \cle{destruction et la fragmentation des habitats} : déforestation, barrages (ex. Lom Pangar, Memve'ele), routes et plantations industrielles (palmier à huile, hévéa, bananier) morcellent les forêts ; les populations animales, isolées en petits groupes, perdent leur diversité génétique (consanguinité) et deviennent vulnérables aux épidémies et aux aléas démographiques ;
  \item la pollution des milieux aquatiques (métaux lourds, hydrocarbures, plastiques) et la \cle{surpêche} (filets maillants de petite taille, empoisonnement à la plante \emph{Tephrosia vogelii} ou aux produits chimiques), qui appauvrissent la faune des fleuves (Sanaga, Wouri, Logone) et des côtes (kribi, Limbé).
\end{itemize}

\begin{aretenir}[Espèces menacées emblématiques au Cameroun]
Parmi les espèces classées « En danger critique » (CR) ou « En danger » (EN) par l'UICN et présentes au Cameroun :
\begin{itemize}
  \item \textbf{Gorille de Cross River} (\emph{Gorilla gorilla diehli}) : l'un des plus rares au monde, quelques centaines d'individus dans les montagnes du Sud-Ouest (parc national de Takamanda, Mone, Kagwene).
  \item \textbf{Éléphant de forêt} et \textbf{Éléphant de savane} : braconnage intensif pour l'ivoire.
  \item \textbf{Pangolins} (géant, à ventre blanc, à queue longue) : mammifères les plus braconnés du monde.
  \item \textbf{Lamantin d'Afrique} (\emph{Trichechus senegalensis}) : piégé dans les filets, chassé pour la viande.
  \item \textbf{Perroquet gris du Gabon} (\emph{Psittacus erithacus}) : capture massive pour le commerce d'oiseaux de cage.
\end{itemize}
La disparition d'une espèce est \textbf{irréversible} : la biodiversité est une ressource renouvelable seulement si les populations conservent un effectif viable et une connectivité génétique suffisante pour se reproduire.
\end{aretenir}

\courssub{Croissance démographique et urbanisation : des facteurs aggravants}

Aucune des pressions précédentes ne peut être comprise sans deux facteurs de fond :
\begin{itemize}
  \item la \cle{croissance démographique} : la population camerounaise est passée d'environ 5,4 millions d'habitants en 1960 (indépendance) à plus de 28 millions aujourd'hui (estimations 2023-2024). Plus de bouches à nourrir signifie plus de terres défrichées, plus de bois-énergie prélevé, plus d'eau consommée, plus de déchets rejetés ;
  \item l'\cle{urbanisation} rapide (taux d'urbanisation > 55 \%) : Douala et Yaoundé concentrent chacun plus de 3 millions d'habitants ; l'étalement urbain détruit terres agricoles et forêts périurbaines (ex. forêt de Mbankolo, mangroves de l'estuaire Wouri), tandis que les infrastructures d'assainissement (stations d'épuration, décharges contrôlées) ne suivent pas, aggravant la pollution des eaux et des sols.
\end{itemize}
Ces deux facteurs n'agissent pas directement : ils \textbf{amplifient} toutes les pressions décrites ci-dessus. C'est pourquoi on les appelle des \cle{facteurs aggravants} (ou facteurs de pression sous-jacents / \emph{drivers}).

\courssub{Méthode : relier une activité humaine à ses impacts à partir d'un document}

\begin{methode}[Analyser l'impact d'une activité sur une ressource]
Face à un document (texte, tableau, carte, photographie, graphique) présentant une activité humaine :
\begin{enumerate}
  \item \textbf{Identifier l'activité} et la \textbf{ressource} concernée (eau, sol, forêt, biodiversité).
  \item \textbf{Repérer le prélèvement ou le rejet} : que prend-on à la ressource (coupes, pompages, captures) ou que lui rejette-on (eaux usées, engrais, fumées, boues) ?
  \item \textbf{Qualifier l'impact direct} : pollution (qualité) ou surexploitation (quantité) ?
  \item \textbf{Suivre la chaîne de conséquences} : impacts indirects sur d'autres ressources (ex. déforestation $\to$ érosion $\to$ envasement des rivières $\to$ baisse de la pêche $\to$ perte de biodiversité).
  \item \textbf{Conclure} en précisant si la ressource peut se régénérer et à quelle condition (arrêt de la pression, restauration, temps nécessaire).
\end{enumerate}
\end{methode}

\begin{exoresolu}[Analyse d'un document : une plantation industrielle de palmiers à huile]
\textbf{Énoncé.} Un document décrit l'installation d'une plantation industrielle de palmiers à huile (type SOCAPALM/HEVECAM) dans le Sud-Ouest : défrichement de \num{5000} ha de forêt secondaire, construction d'une usine rejetant ses effluents (effluents de trituration riches en matière organique) dans une rivière, emploi d'engrais (NPK) et d'herbicides (glyphosate). Identifie les pressions exercées sur chaque ressource naturelle et construis la chaîne de causalité correspondante.
\tcblower
\textbf{Solution détaillée.}
\begin{itemize}
  \item \textbf{Forêt} : défrichement de \num{5000} ha $\to$ impact direct : destruction quantitative de la forêt (surexploitation/perte de ressource et d'habitat).
  \item \textbf{Biodiversité} : impact indirect de la déforestation $\to$ destruction et fragmentation des habitats $\to$ fuite ou disparition des espèces forestières (singes, céphalophes, oiseaux, insectes) $\to$ perte de connectivité entre aires protégées.
  \item \textbf{Eau} : rejets d'effluents d'usine (DBO\textsubscript{5} élevée) + lessivage des engrais (nitrates, phosphates) et pesticides $\to$ pollution (qualité) de la rivière $\to$ eutrophisation, mortalité des poissons par anoxie, eau impropre à la consommation en aval.
  \item \textbf{Sol} : monoculture intensive sur sols nus entre les jeunes palmiers $\to$ érosion hydrique $\to$ perte de la couche arable et envasement de la rivière (impact indirect sur l'eau) ; à long terme, épuisement de la fertilité (exportation de K, Mg par les régimes).
\end{itemize}
\textbf{Conclusion} : une seule activité exerce des pressions croisées sur les quatre ressources ; la déforestation initiale est le maillon commun de plusieurs chaînes de causalité. La gestion durable impose des mesures d'atténuation (couloirs biologiques, bassins de décantation, couverts végétaux, zones tampons riveraines).
\end{exoresolu}

\begin{exercice}[À toi de jouer]
Un article de presse rapporte qu'au bord du Logone (Extrême-Nord), les pêcheurs capturent de moins en moins de poissons depuis la construction de retenues d'eau (barrages, seuils) pour l'irrigation du riz en amont (périmètres de Maga, Yagoua). En appliquant la méthode en cinq étapes, construis la chaîne de causalité complète et précise, pour chaque maillon, s'il s'agit de pollution ou de surexploitation.
\end{exercice}

\begin{corrige}[Exercice : pêche et irrigation sur le Logone]
\begin{enumerate}
  \item \textbf{Activité} : irrigation intensive du riz (périmètres aménagés) ; \textbf{Ressource} : l'eau du fleuve Logone et la ressource halieutique (poissons).
  \item \textbf{Prélèvement} : captage d'une grande partie du débit du fleuve par les retenues et canaux d'irrigation (surexploitation quantitative de l'eau).
  \item \textbf{Impact direct} : \textbf{surexploitation} de l'eau (baisse quantitative drastique du débit en aval, assèchement des bras morts et zones inondables) — ce n'est pas une pollution.
  \item \textbf{Chaîne de conséquences} : baisse du débit $\to$ réduction drastique des zones inondables (yaérés) qui servent de frayères et de zones d'alimentation pour les alevins $\to$ diminution de la reproduction et du renouvellement des stocks $\to$ effondrement des captures (surexploitation indirecte de la ressource « poissons » par destruction de son habitat de reproduction).
  \item \textbf{Conclusion} : la ressource halieutique, pourtant renouvelable, s'épuise car sa régénération dépend d'une ressource en eau elle-même surexploitée. La solution passe par une gestion concertée du débit réservé (débit d'étiage maintenu) et la restauration des plaines d'inondation.
\end{enumerate}
\end{corrige}

\begin{aretenir}[L'essentiel de la section]
\begin{itemize}
  \item Une \cle{pression} est une action humaine qui dégrade une ressource en \textbf{qualité} (pollution) ou en \textbf{quantité} (surexploitation).
  \item \textbf{Eau} : pollutions domestiques, industrielles, agricoles (eutrophisation) ; prélèvements excessifs (cas du lac Tchad, nappes de Yaoundé/Maroua).
  \item \textbf{Sols} : érosion hydrique/éolienne, épuisement de la fertilité (perte humus, minéraux), compaction, désertification dans le Nord (climat + activités humaines).
  \item \textbf{Forêts} : exploitation (bois d'œuvre, bois-énergie), agriculture sur brûlis (jachères trop courtes), feux, urbanisation, plantations industrielles ; perte de \num{100000} à \num{200000} ha/an.
  \item \textbf{Biodiversité} : braconnage (ivoire, écailles, viande), destruction/fragmentation des habitats (routes, barrages, plantations), pollution/surpêche.
  \item \textbf{Facteurs aggravants} : croissance démographique et urbanisation rapide amplifient toutes les pressions.
  \item \textbf{Méthode d'analyse} : activité $\to$ prélèvement/rejet $\to$ impact direct (pollution/surexploitation) $\to$ impacts indirects (chaîne) $\to$ conséquence sur la ressource (régénération possible ?).
\end{itemize}
\end{aretenir}

Ces pressions étant désormais identifiées et comprises, une question s'impose : comment concilier les besoins légitimes des populations actuelles avec la conservation des ressources pour les générations futures ? C'est précisément l'objet de la section suivante, consacrée à la notion de \cle{développement durable} et de \cle{gestion raisonnée}.

\coursec{Le développement durable : notion et principes}

Dans la section précédente, nous avons dressé un constat inquiétant : la déforestation du bassin du Congo, la pollution des eaux du Wouri, l'épuisement des sols dans l'Extrême-Nord, la surpêche à Kribi... Les activités humaines exercent des pressions croissantes sur les ressources naturelles, au point de menacer leur renouvellement. Faut-il alors cesser toute exploitation ? Évidemment non : les populations ont besoin de bois, de poisson, de terres cultivées et d'eau pour vivre. Toute la question est donc de savoir \emph{comment} exploiter sans détruire. C'est précisément à cette question que répond la notion de \cle{développement durable}, que nous allons maintenant définir et explorer.

\courssub{La notion de développement durable}

\textbf{Une situation concrète pour comprendre}\par

Imaginons un village camerounais installé au bord d'une forêt communautaire, par exemple dans la région de l'Est. Chaque année, les habitants coupent du bois pour se chauffer, construire leurs cases et vendre du charbon. Deux attitudes sont possibles :

\begin{itemize}
\item \textbf{Attitude 1 :} couper tous les arbres exploitables, sans replanter. Les revenus sont élevés pendant quelques années, puis la forêt disparaît : plus de bois, sols érodés, sources taries, gibier parti. Les enfants du village héritent d'une terre appauvrie.
\item \textbf{Attitude 2 :} ne couper chaque année qu'une partie des arbres matures, replanter systématiquement, protéger les zones de régénération. Les revenus sont un peu plus faibles chaque année, mais ils se maintiennent indéfiniment : la forêt se régénère, et les enfants du village pourront exploiter la même ressource que leurs parents.
\end{itemize}

L'attitude 2 illustre l'idée centrale du développement durable : \textbf{répondre aux besoins d'aujourd'hui sans empêcher les générations futures de répondre aux leurs}.

\begin{definition}[Développement durable]
Le \cle{développement durable} est un mode de développement qui \textbf{répond aux besoins des générations présentes sans compromettre la capacité des générations futures à répondre à leurs propres besoins}. Cette définition a été officialisée en 1987 dans le rapport « Notre avenir à tous » de la Commission mondiale sur l'environnement et le développement de l'ONU, dite \textbf{Commission Brundtland} (du nom de sa présidente, Gro Harlem Brundtland).
\end{definition}

Deux idées sont contenues dans cette définition et doivent être retenues séparément :

\begin{itemize}
\item l'idée de \textbf{besoins} : le développement doit améliorer les conditions de vie des populations actuelles (alimentation, santé, éducation, revenus), en particulier des plus pauvres ;
\item l'idée de \textbf{limites} : les ressources de la planète et la capacité des écosystèmes à absorber les pollutions ne sont pas infinies ; l'état de l'environnement que nous transmettons conditionne les possibilités des générations futures.
\end{itemize}

\begin{savaistu}[D'où vient l'expression « développement durable » ?]
L'expression « développement durable » (en anglais \emph{sustainable development}) a été officialisée en \textbf{1987} par la \textbf{Commission Brundtland} de l'Organisation des Nations Unies, dans le rapport « Notre avenir à tous » (\emph{Our Common Future}). Elle a ensuite été popularisée par le \textbf{Sommet de la Terre de Rio de Janeiro en 1992}, première grande conférence mondiale réunissant chefs d'État, scientifiques et associations autour de l'environnement et du développement. Depuis, le concept structure les politiques environnementales du monde entier, y compris au Cameroun (stratégies nationales de développement durable, loi-cadre sur la gestion de l'environnement de 1996).
\end{savaistu}

\courssub{Les trois piliers du développement durable}

\textbf{Observation et problème}\par

Reprenons notre village forestier. Un projet de « conservation » consiste à interdire totalement l'accès à la forêt. Résultat : la forêt est protégée, mais les habitants n'ont plus de bois ni de revenus ; la pauvreté augmente, les jeunes quittent le village, et certains braconnent la nuit. Le projet échoue. À l'inverse, un projet purement économique (exploitation intensive par une entreprise) détruit la forêt et ne profite pas aux habitants. Il échoue aussi.

\textbf{Conclusion de l'observation :} une solution durable ne peut pas viser un seul objectif. Elle doit concilier \textbf{trois dimensions} simultanément.

\begin{definition}[Les trois piliers du développement durable]
Le développement durable repose sur \textbf{trois piliers indissociables} :
\begin{itemize}
\item le pilier \cle{environnemental} : protéger les écosystèmes, préserver les ressources naturelles et la biodiversité, limiter les pollutions ;
\item le pilier \cle{économique} : assurer une activité économique viable, créatrice de revenus et d'emplois, sans gaspillage des ressources ;
\item le pilier \cle{social} : garantir l'équité, l'amélioration des conditions de vie (santé, éducation, alimentation) et la participation des populations.
\end{itemize}
Un développement n'est \textbf{durable} que s'il se situe à l'\textbf{intersection des trois piliers} : économiquement viable, socialement équitable et écologiquement soutenable.
\end{definition}

\begin{popfigure}
\begin{center}
\begin{tikzpicture}[scale=1.0]
% Cercle environnement (haut)
\fill[popgreen!35] (0,1.3) circle (2.2);
% Cercle économie (bas gauche)
\fill[popblue!30] (-1.6,-0.9) circle (2.2);
% Cercle social (bas droit)
\fill[poporange!35] (1.6,-0.9) circle (2.2);
% Contours
\draw[popgreen!80!black, thick] (0,1.3) circle (2.2);
\draw[popblue!80!black, thick] (-1.6,-0.9) circle (2.2);
\draw[poporange!80!black, thick] (1.6,-0.9) circle (2.2);
% Étiquettes des piliers
\node[font=\bfseries, text=popgreen!60!black] at (0,3.0) {Environnement};
\node[font=\bfseries, text=popblue!60!black] at (-3.0,-2.3) {Économie};
\node[font=\bfseries, text=poporange!70!black] at (3.0,-2.3) {Social};
% Intersections partielles
\node[font=\small, text=popdark] at (-2.0,0.9) {Viable};
\node[font=\small, text=popdark] at (2.0,0.9) {Équitable};
\node[font=\small, text=popdark] at (0,-1.7) {Soutenable};
% Intersection centrale
\node[font=\bfseries\small, text=popdark, align=center] at (0,-0.1) {Développement\\durable};
\end{tikzpicture}
\end{center}
\legende{Diagramme des trois piliers du développement durable. L'intersection économie--environnement correspond à un développement \emph{viable} ; social--environnement à un développement \emph{soutenable} ; économie--social à un développement \emph{équitable}. Seule la zone commune aux trois cercles correspond à un développement réellement \emph{durable}.}
\end{popfigure}

\textbf{Pourquoi l'équilibre des trois piliers est-il nécessaire ?}\par

Raisonnons sur chaque « pilier manquant » :

\begin{itemize}
\item \textbf{Sans le pilier environnemental} (économie + social seuls) : l'exploitation peut être rentable et créer des emplois, mais elle épuise la ressource ; l'activité s'effondre à terme, entraînant avec elle emplois et revenus. Exemple : l'effondrement d'un stock de poissons surexploité ruine les pêcheurs eux-mêmes.
\item \textbf{Sans le pilier économique} (environnement + social seuls) : la protection est idéale sur le papier, mais sans ressources financières ni activité viable, elle n'est pas financée et ne dure pas.
\item \textbf{Sans le pilier social} (économie + environnement seuls) : l'exploitation peut être rentable et écologiquement correcte, mais si les populations locales sont exclues des bénéfices, les inégalités et les conflits grandissent, et le projet perd sa légitimité (braconnage, sabotage, contestation).
\end{itemize}

\begin{aretenir}[L'essentiel sur les trois piliers]
Le développement durable n'est pas synonyme de « protection de la nature » seule. C'est la recherche d'un \textbf{équilibre entre trois objectifs} : efficacité économique, équité sociale et préservation de l'environnement. Négliger l'un des trois condamne le projet à l'échec à moyen ou long terme.
\end{aretenir}

\begin{attention}[Erreur fréquente au Probatoire]
Beaucoup d'élèves écrivent : « le développement durable, c'est arrêter d'exploiter les ressources pour protéger la nature ». C'est \textbf{faux}. Le développement durable ne signifie \textbf{pas} l'arrêt de toute exploitation : il signifie une exploitation \textbf{raisonnée}, qui maintient la ressource tout en satisfaisant les besoins humains. Dans une copie, associez toujours développement durable aux trois piliers et à l'idée d'équilibre, jamais à une interdiction totale.
\end{attention}

\courssub{La gestion raisonnée d'une ressource : exploiter sans épuiser}

\textbf{Démarche d'investigation : à quelle condition une pêche peut-elle durer ?}\par

Considérons le lac de pêche d'un village (situation inspirée des pêcheries du Logone ou du lac Tchad). Menons la démarche scientifique complète.

\textbf{Observation.} Les pêcheurs constatent que lorsque le nombre de pirogues augmente trop, les captures par pirogue diminuent d'année en année, et les poissons pêchés sont de plus en plus petits.

\textbf{Hypothèse.} Le rythme de prélèvement (captures annuelles) dépasse le rythme de renouvellement naturel du stock de poissons (reproduction et croissance).

\textbf{Données.} On estime que le lac contient un stock de poissons qui se régénère naturellement d'environ 100 tonnes par an (naissances et croissance compensant les morts naturelles). On compare trois scénarios de pêche sur plusieurs années :

\begin{center}
\begin{tabularx}{\linewidth}{|l|X|X|X|}
\hline
\rowcolor{popblueL} \textbf{Scénario} & \textbf{Prélèvement annuel} & \textbf{Régénération annuelle} & \textbf{Évolution du stock} \\
\hline
A & 80 tonnes & 100 tonnes & Stock stable ou en légère augmentation : pêche durable \\
\hline
B & 100 tonnes & 100 tonnes & Stock stable, à la limite : pêche au seuil de durabilité \\
\hline
C & 130 tonnes & 100 tonnes & Stock en diminution régulière : épuisement à terme \\
\hline
\end{tabularx}
\end{center}

\textbf{Interprétation.} Tant que le prélèvement reste inférieur ou égal à la régénération, le stock se maintient : la pêche peut durer indéfiniment. Dès que le prélèvement dépasse durablement la régénération, le stock décroît, la régénération elle-même diminue (moins de reproducteurs), et l'épuisement s'accélère.

\textbf{Conclusion.} Une exploitation est durable si, et seulement si, le rythme de prélèvement ne dépasse pas le rythme de régénération de la ressource.

\begin{propriete}[Condition de durabilité d'une exploitation — admise]
Une exploitation d'une ressource renouvelable est \cle{durable} si le \textbf{taux de prélèvement reste inférieur ou égal au taux de régénération} de la ressource :
\[ \text{prélèvement annuel} \leq \text{régénération annuelle} \]
Cette propriété est \textbf{admise} (sa justification complète relève de modèles mathématiques de dynamique des populations non exigibles au Probatoire), mais tu dois savoir l'\textbf{énoncer} et l'\textbf{appliquer} à des données chiffrées simples.
\end{propriete}

\begin{popfigure}
\begin{center}
\begin{tikzpicture}
\begin{axis}[
width=0.85\linewidth, height=6.5cm,
xlabel={Temps (années)},
ylabel={Stock de la ressource (tonnes)},
xmin=0, xmax=20, ymin=0, ymax=1400,
xtick={0,5,10,15,20},
ytick={0,400,800,1000,1200},
legend style={at={(0.03,0.03)}, anchor=south west, font=\small},
grid=both, grid style={popline!40},
axis line style={popdark},
]
% Régénération seule : stock stable
\addplot[popgreen, very thick, domain=0:20, samples=50] {1000};
\addlegendentry{Prélèvement = régénération (durable)}
% Prélèvement < régénération : stock augmente
\addplot[popblue, very thick, domain=0:20, samples=50] {800 + 20*x};
\addlegendentry{Prélèvement $<$ régénération (durable)}
% Prélèvement > régénération : déclin
\addplot[poporange, very thick, domain=0:20, samples=50] {1000 - 55*x + 1.2*x*x};
\addlegendentry{Prélèvement $>$ régénération (épuisement)}
\end{axis}
\end{tikzpicture}
\end{center}
\legende{Évolution du stock d'une ressource renouvelable selon le rythme d'exploitation. Lorsque le prélèvement annuel est inférieur ou égal à la régénération annuelle, le stock se maintient ou augmente (courbes bleue et verte) : l'exploitation est durable. Lorsque le prélèvement dépasse la régénération (courbe orange), le stock s'effondre : la ressource s'épuise.}
\end{popfigure}

\begin{definition}[Gestion raisonnée (ou gestion durable) d'une ressource]
La \cle{gestion raisonnée} d'une ressource naturelle consiste à organiser son exploitation de telle sorte que le \textbf{prélèvement ne dépasse jamais durablement la capacité de régénération} de la ressource, tout en satisfaisant les besoins humains. On dit : \textbf{exploiter sans épuiser}.
\end{definition}

\textbf{Les outils concrets de la gestion raisonnée}\par

Pour maintenir le prélèvement sous le seuil de durabilité, on utilise plusieurs instruments :

\begin{itemize}
\item le \cle{quota} : on fixe une quantité maximale prélevable par an (tonnage de pêche, nombre de mètres cubes de bois, nombre d'animaux chassés). Exemple : quotas de pêche imposés aux flottilles industrielles au large des côtes camerounaises ;
\item la \cle{rotation} : on divise l'espace exploité en parcelles exploitées à tour de rôle, le temps que les autres se régénèrent. En foresterie, on parle de \textbf{rotation des coupes} : une parcelle coupée n'est réexploitée qu'après le temps nécessaire à la reconstitution du peuplement (souvent 25 à 30 ans pour les essences exploitées en forêt dense camerounaise) ;
\item le \cle{renouvellement} (reboisement, restauration) : on compense le prélèvement par une régénération active. Exemple de la \textbf{coupe forestière avec reboisement} : pour chaque arbre exploité, on plante un ou plusieurs jeunes plants d'essences locales, et l'on protège la régénération naturelle ;
\item les \textbf{périodes de fermeture} : interdiction de pêche ou de chasse pendant la saison de reproduction (par exemple la fermeture de la pêche pendant la période de frai), pour laisser les populations se reproduire ;
\item la \textbf{taille minimale de capture} : ne prélever que les individus ayant atteint la maturité sexuelle, afin que chacun ait eu le temps de se reproduire au moins une fois.
\end{itemize}

\begin{methode}[Juger de la durabilité d'une exploitation à partir de données]
Face à un document chiffré (tableau, graphique) décrivant l'exploitation d'une ressource :
\begin{enumerate}
\item \textbf{Repérer} dans le document le prélèvement annuel (quantité exploitée par an) et la régénération annuelle (capacité de renouvellement par an).
\item \textbf{Comparer} quantitativement les deux valeurs : calculer éventuellement le rapport prélèvement/régénération.
\item \textbf{Conclure} : si prélèvement $\leq$ régénération, l'exploitation est durable ; si prélèvement $>$ régénération de façon répétée, elle conduit à l'épuisement.
\item \textbf{Proposer} le cas échéant des mesures correctives adaptées : réduction du quota, rotation, reboisement, période de fermeture, taille minimale de capture.
\end{enumerate}
\end{methode}

\begin{exoresolu}[La forêt communautaire de la région de l'Est]
Une forêt communautaire de la région de l'Est du Cameroun couvre \SI{5000}{ha}. Les inventaires forestiers montrent que la forêt produit naturellement, par croissance des arbres, environ \SI{1,2}{m^3} de bois exploitable par hectare et par an. Le plan de gestion autorise une coupe annuelle de \SI{4000}{m^3} avec reboisement des parcelles coupées.

\textbf{1.} Calculer la régénération annuelle totale de la forêt.

\textbf{2.} L'exploitation est-elle durable ? Justifier quantitativement.

\textbf{3.} Une entreprise propose de porter la coupe annuelle à \SI{9000}{m^3}. Que peut-on prévoir ?

\tcblower

\textbf{1.} Régénération annuelle totale :
\[ 5000 \times 1,2 = 6000 \ \si{m^3/an} \]
La forêt se régénère de \SI{6000}{m^3} de bois exploitable chaque année.

\textbf{2.} On compare prélèvement et régénération :
\[ \text{prélèvement} = 4000 \ \si{m^3/an} < \text{régénération} = 6000 \ \si{m^3/an} \]
Le taux de prélèvement est inférieur au taux de régénération : le stock de bois se maintient (il augmente même légèrement). De plus, le reboisement accélère la reconstitution des parcelles coupées. L'exploitation est donc \textbf{durable}.

\textbf{3.} Avec une coupe de \SI{9000}{m^3/an} :
\[ 9000 > 6000 \ \si{m^3/an} \]
Le prélèvement dépasse la régénération de \SI{3000}{m^3} chaque année : le stock forestier diminue d'année en année. À terme, la forêt s'appauvrit, la régénération elle-même chute (moins d'arbres producteurs et de semenciers), et l'exploitation s'effondre. Cette proposition n'est \textbf{pas durable} ; il faudrait la rejeter ou la ramener sous le seuil de \SI{6000}{m^3/an}, idéalement avec une marge de sécurité.
\end{exoresolu}

\courssub{Principe de précaution et responsabilité intergénérationnelle}

\textbf{Le principe de précaution}\par

Certaines activités humaines présentent des risques graves mais mal connus : rejets de produits chimiques nouveaux, introduction d'espèces exotiques, exploitation en eaux profondes... Faut-il attendre la preuve scientifique complète du danger pour agir ? L'histoire montre que cette attente peut être catastrophique : lorsque le danger est enfin prouvé, les dégâts sont souvent irréversibles (espèces disparues, sols contaminés pour des siècles).

\begin{definition}[Principe de précaution]
Le \cle{principe de précaution} affirme que l'\textbf{absence de certitude scientifique complète ne doit pas servir de prétexte pour retarder l'adoption de mesures} visant à prévenir un risque de dommage grave ou irréversible à l'environnement. En cas de doute sérieux, on s'abstient ou on encadre strictement l'activité.
\end{definition}

\begin{exemplebox}[Appliquer le principe de précaution]
Avant d'autoriser l'épandage massif d'un nouvel herbicide dans les plantations du littoral camerounais, on exige des études d'impact sur les cours d'eau et la santé des populations. Tant que la toxicité à long terme n'est pas établie, on limite les doses autorisées et on surveille les nappes phréatiques. On agit \emph{avant} la catastrophe, pas après.
\end{exemplebox}

\textbf{La responsabilité intergénérationnelle}\par

Le développement durable repose enfin sur une exigence éthique : la \cle{responsabilité intergénérationnelle}. Les générations actuelles ne sont pas propriétaires des ressources de la planète ; elles en sont les \textbf{dépositaires} pour les générations futures. Chaque décision d'exploitation doit donc intégrer ses conséquences à long terme : un gisement minier épuisé, une espèce éteinte ou un sol stérilisé par l'érosion sont des pertes \textbf{définitives} que subiront nos enfants et petits-enfants sans en être responsables.

\begin{aretenir}[Les deux principes complémentaires]
\begin{itemize}
\item \textbf{Principe de précaution} : en cas de risque grave ou irréversible, l'incertitude scientifique n'empêche pas d'agir pour protéger.
\item \textbf{Responsabilité intergénérationnelle} : transmettre aux générations futures un patrimoine naturel au moins aussi riche que celui reçu.
\end{itemize}
Ces deux principes fondent juridiquement et moralement les politiques de gestion durable des ressources.
\end{aretenir}

\courssub{Complément pour le Probatoire : l'empreinte écologique}

\begin{definition}[Empreinte écologique — complément utile au Probatoire]
L'\cle{empreinte écologique} est un \textbf{indicateur de la pression exercée par une population humaine sur la nature}. Elle mesure la surface de terres et de mers biologiquement productives nécessaire pour produire les ressources consommées par cette population et absorber ses déchets (notamment le \ce{CO2} émis). Elle s'exprime en \textbf{hectares globaux (hag) par habitant}.
\end{definition}

L'intérêt de cet indicateur est de permettre une comparaison simple : on la confronte à la \textbf{biocapacité} de la planète (la surface productive réellement disponible par habitant, environ 1,6 hag). Si l'empreinte écologique moyenne dépasse la biocapacité, l'humanité consomme plus que ce que la Terre régénère : elle vit « à crédit », en liquidant le capital naturel. C'est le cas actuellement à l'échelle mondiale : l'empreinte de l'humanité correspond à plus d'une planète et demie. Les pays riches ont des empreintes très élevées (souvent 4 à 8 hag par habitant), tandis que les pays comme le Cameroun ont des empreintes plus faibles, mais une biodiversité et des forêts dont la préservation est un enjeu mondial.

\begin{attention}[Ne pas confondre les notions]
Ne confonds pas \textbf{empreinte écologique} (surface nécessaire pour soutenir un mode de vie) et \textbf{gestion raisonnée} (mode d'exploitation d'une ressource précise). L'empreinte écologique est un \emph{indicateur global de pression} ; la gestion raisonnée est un \emph{outil d'action} sur une ressource donnée. Les deux sont liées : généraliser la gestion raisonnée fait diminuer l'empreinte écologique.
\end{attention}

\courssub{Synthèse de la section}

Le développement durable, défini par le rapport Brundtland (1987), est un développement qui répond aux besoins présents sans compromettre ceux des générations futures. Il exige l'équilibre de trois piliers — environnemental, économique et social — et se traduit concrètement par la \textbf{gestion raisonnée} des ressources : maintenir le prélèvement annuel au niveau ou en dessous de la régénération annuelle, grâce à des outils comme les quotas, la rotation des coupes, le reboisement ou les périodes de fermeture. Guidée par le principe de précaution et la responsabilité intergénérationnelle, cette approche sera mise en pratique dans la section suivante, à travers les politiques concrètes de gestion durable des ressources au Cameroun.

\begin{exercice}[Pour t'entraîner]
Un stock de crevettes au large de Kribi se régénère de 500 tonnes par an. Les captures annuelles sont passées de 300 à 700 tonnes en dix ans.
\begin{enumerate}
\item À partir de quelles captures annuelles l'exploitation cesse-t-elle d'être durable ?
\item Prévois l'évolution du stock si les captures restent à 700 tonnes par an.
\item Propose deux mesures de gestion raisonnée adaptées à cette pêcherie, en justifiant chacune.
\end{enumerate}
\end{exercice}

\begin{corrige}[Exercice : la pêche crevettière de Kribi]
\begin{enumerate}
\item L'exploitation cesse d'être durable lorsque le prélèvement dépasse la régénération, c'est-à-dire au-delà de 500 tonnes par an.
\item À 700 tonnes prélevées pour 500 tonnes régénérées, le stock diminue de 200 tonnes par an. Il s'épuise progressivement ; la régénération elle-même finira par chuter (moins de reproducteurs), accélérant l'effondrement de la pêcherie.
\item Deux mesures possibles : (a) fixer un \textbf{quota} annuel de captures inférieur à 500 tonnes, avec marge de sécurité (par exemple 400 tonnes) ; (b) instaurer une \textbf{fermeture de la pêche pendant la période de reproduction} et une taille minimale de capture, pour protéger les reproducteurs et les juvéniles. Ces mesures ramènent le prélèvement sous le seuil de régénération et permettent la reconstitution du stock.
\end{enumerate}
\end{corrige}

\coursec{La gestion durable des ressources au Cameroun}

Dans la section précédente, tu as découvert le \cle{développement durable} : un développement qui répond aux besoins du présent sans compromettre la capacité des générations futures à répondre aux leurs. Mais une notion ne vaut que si elle se traduit en actes. Comment un pays comme le Cameroun, riche en eau, en forêts et en biodiversité, passe-t-il du principe à la pratique ? C'est ce que nous allons voir : pour chaque grande ressource naturelle, des mesures concrètes de gestion durable existent, portées par l'État, les communes et les populations locales.

\courssub{La gestion durable de l'eau}

\textbf{Situation concrète.} Dans certains quartiers de Yaoundé, les puits sont contaminés après les pluies ; dans l'Extrême-Nord, le lac Tchad a perdu plus de 90 \% de sa superficie depuis les années 1960. L'eau est abondante au Cameroun, mais sa qualité et sa disponibilité sont menacées.

\begin{attention}[L'eau n'est pas inépuisable]
Une ressource renouvelable n'est renouvelable que si le prélèvement reste inférieur à la capacité de renouvellement. Pomper une nappe plus vite qu'elle ne se recharge, ou polluer une rivière plus vite qu'elle ne s'épure, c'est transformer une ressource renouvelable en ressource épuisée.
\end{attention}

Les principales mesures de gestion durable de l'eau au Cameroun sont :

\begin{itemize}
\item \textbf{La protection des \cle{bassins versants}} : un bassin versant est l'ensemble des terres dont les eaux de pluie alimentent un même cours d'eau. Reboiser les versants, interdire les cultures sur les berges et protéger les sources limitent l'érosion et la pollution des rivières.
\item \textbf{Le traitement des \cle{eaux usées}} : les eaux domestiques et industrielles doivent être épurées (décantation, traitement biologique par des bactéries) avant d'être rejetées, afin de ne pas contaminer les nappes et les cours d'eau.
\item \textbf{Les forages raisonnés} : un forage capte l'eau d'une nappe souterraine. Il doit être implanté loin des sources de pollution (latrines, dépôts d'ordures) et exploité sans dépasser le débit de recharge de la nappe.
\item \textbf{Le projet de restauration du \cle{lac Tchad}} : le Cameroun participe, avec les pays riverains réunis dans la Commission du Bassin du Lac Tchad (CBLT), à des projets de transfert d'eau depuis le bassin de l'Oubangui et de gestion concertée des prélèvements, afin de restaurer le lac et les activités (pêche, agriculture) qui en dépendent.
\end{itemize}

\courssub{La gestion durable des sols}

\textbf{Situation concrète.} Dans l'Ouest camerounais, des collines cultivées de haut en bas perdent chaque année leur terre fertile, emportée par les pluies : les rendements chutent et les rivières s'ensablent. Le sol se forme en des siècles, mais se détruit en quelques saisons.

Les techniques de gestion durable des sols sont :

\begin{itemize}
\item \textbf{L'\cle{agroforesterie}} : association, sur une même parcelle, d'arbres et de cultures (par exemple caféiers sous ombrage, cacaoyers et arbres fruitiers). Les racines des arbres fixent le sol, leurs feuilles fournissent de la matière organique et limitent l'érosion.
\item \textbf{Les cultures en \cle{courbes de niveau}} : sur les reliefs, on cultive en bandes horizontales suivant les courbes de niveau, et non dans le sens de la pente. L'eau de pluie est ainsi freinée et s'infiltre au lieu de ruisseler et d'emporter la terre.
\item \textbf{La \cle{jachère}} : mise en repos temporaire d'une parcelle, qui permet au sol de reconstituer sa fertilité grâce à la décomposition de la végétation spontanée. Une jachère efficace dure plusieurs années ; sa durée doit être respectée.
\item \textbf{La lutte contre les \cle{feux de brousse}} : les feux tardifs détruisent la matière organique, la faune et les jeunes arbres. Pare-feux, feux précoces contrôlés et sensibilisation des populations limitent les dégâts.
\item \textbf{La \cle{Grande Muraille Verte}} : vaste initiative panafricaine (à laquelle le Cameroun participe) visant à planter une bande d'arbres à travers le Sahel, du Sénégal à Djibouti, pour stopper l'avancée du désert et restaurer les terres dégradées.
\end{itemize}

\courssub{La gestion durable des forêts}

\textbf{Situation concrète.} Le Sud camerounais abrite la deuxième forêt tropicale du monde (bassin du Congo). L'exploitation forestière rapporte des devises, mais une coupe anarchique détruit la forêt, fait disparaître la faune et appauvrit les villages. Comment exploiter sans détruire ?

\begin{definition}[Forêt communautaire]
Une \cle{forêt communautaire} est une portion de forêt (au maximum \SI{5000}{hectares}) attribuée à une communauté villageoise, qui la gère elle-même selon un \cle{plan de gestion} (ou plan d'aménagement) approuvé par l'administration forestière. Les revenus de l'exploitation reviennent à la communauté, qui doit en contrepartie respecter les règles de coupe et de reboisement.
\end{definition}

Les outils de gestion durable des forêts au Cameroun sont :

\begin{itemize}
\item \textbf{Les forêts communautaires} : elles responsabilisent les villages, qui deviennent les premiers gardiens de leur forêt.
\item \textbf{Les concessions avec \cle{plan d'aménagement}} : une entreprise forestière ne peut exploiter une concession que selon un plan qui découpe la forêt en parcelles exploitées à tour de rôle (rotation de 25 à 30 ans), avec des \cle{coupes sélectives} : on ne prélève que les arbres matures d'essences autorisées, au-dessus d'un diamètre minimal, en laissant les jeunes arbres assurer la régénération.
\item \textbf{Le \cle{reboisement}} : plantation d'arbres (essences locales ou à croissance rapide) sur les surfaces exploitées ou dégradées.
\item \textbf{La \cle{certification}} (par exemple le label FSC) : elle garantit, par un contrôle indépendant, que le bois provient d'une forêt gérée durablement, ce qui valorise le bois camerounais sur les marchés internationaux.
\end{itemize}

\begin{popfigure}
\begin{center}
\begin{tikzpicture}[scale=1.05, every node/.style={font=\small}]
\node[draw=popblue, fill=popblueL, rounded corners, minimum width=3.4cm, minimum height=1cm] (inv) {1. Inventaire forestier};
\node[draw=popblue, fill=popblueL, rounded corners, minimum width=3.4cm, minimum height=1cm, right=1.1cm of inv] (plan) {2. Plan d'aménagement};
\node[draw=poporange, fill=poporangeL, rounded corners, minimum width=3.4cm, minimum height=1cm, below=1.1cm of plan] (coupe) {3. Coupe sélective};
\node[draw=popgreen, fill=popgreenL, rounded corners, minimum width=3.4cm, minimum height=1cm, left=1.1cm of coupe] (rebois) {4. Reboisement};
\node[draw=poppurple, fill=poppurpleL, rounded corners, minimum width=3.4cm, minimum height=1cm, below=1.1cm of rebois] (suivi) {5. Suivi et contrôle};
\draw[-{Stealth}, thick, popdark] (inv) -- (plan);
\draw[-{Stealth}, thick, popdark] (plan) -- (coupe);
\draw[-{Stealth}, thick, popdark] (coupe) -- (rebois);
\draw[-{Stealth}, thick, popdark] (rebois) -- (suivi);
\draw[-{Stealth}, thick, popdark] (suivi.west) -- ++(-1.2,0) |- (inv.west);
\node[font=\small\itshape, popmuted] at ($(inv.west)+(-1.2,-3.3)$) {retour à l'inventaire};
\end{tikzpicture}
\end{center}
\legende{Cycle de gestion raisonnée d'une forêt communautaire : après l'inventaire des ressources, la communauté élabore un plan d'aménagement approuvé par l'État, pratique des coupes sélectives, reboise les parcelles exploitées, puis assure le suivi avant un nouvel inventaire.}
\end{popfigure}

\begin{savaistu}[Une innovation camerounaise]
Le Cameroun a créé en 1994 (loi forestière n° 94/01) les premières \cle{forêts communautaires} d'Afrique centrale. Ce modèle, qui permet aux villages de gérer eux-mêmes leurs forêts et d'en tirer des revenus, a depuis été repris par de nombreux pays du bassin du Congo.
\end{savaistu}

\courssub{La conservation de la biodiversité}

\textbf{Situation concrète.} Au parc national de Waza, les éléphants, girafes et lions attirent des visiteurs du monde entier. Mais le braconnage et l'avancée des cultures menacent ces espèces. Protéger la biodiversité, c'est protéger un patrimoine irremplaçable : une espèce disparue ne revient jamais.

\begin{definition}[Aire protégée]
Une \cle{aire protégée} est un espace géographique délimité, reconnu et géré, par des moyens légaux ou d'autres moyens efficaces, en vue d'assurer la conservation durable de la nature. Selon le niveau de protection, on distingue notamment les \cle{parcs nationaux} (protection intégrale, aucune exploitation) et les \cle{réserves} (de faune, de forêt), où certaines activités encadrées peuvent être tolérées.
\end{definition}

Le Cameroun possède un réseau remarquable d'aires protégées :

\begin{itemize}
\item le \textbf{parc national de Waza} (Extrême-Nord) : faune de savane (éléphants, girafes, lions, antilopes) ;
\item le \textbf{parc national de la Bénoué} (Nord) : savane soudanienne, hippopotames, buffles, lions ;
\item le \textbf{parc national de Korup} (Sud-Ouest) : l'une des plus anciennes forêts tropicales d'Afrique, d'une richesse botanique exceptionnelle ;
\item le \textbf{parc national de Campo-Ma'an} (Sud) : gorilles, chimpanzés, éléphants de forêt ;
\item la \textbf{réserve de faune du Dja} (Sud-Est), classée au patrimoine mondial de l'UNESCO.
\end{itemize}

\begin{exemplebox}[La réserve de faune du Dja]
Classée au patrimoine mondial de l'UNESCO depuis 1987, la réserve de faune du \cle{Dja} couvre \SI{526000}{hectares}, soit l'une des plus vastes forêts tropicales protégées d'Afrique. Presque entièrement entourée par la rivière Dja, elle abrite des gorilles, des chimpanzés, des éléphants de forêt et plus de 300 espèces d'oiseaux.
\end{exemplebox}

\begin{popfigure}
\begin{center}
\begin{tikzpicture}[scale=0.9, font=\small]
% Contour simplifié du Cameroun
\draw[thick, popdark, fill=popcream] (2.2,5.6) -- (3.0,6.4) -- (3.6,5.8) -- (4.6,5.4) -- (5.2,4.6) -- (5.0,3.6) -- (5.6,2.8) -- (5.2,1.6) -- (4.4,0.8) -- (3.4,0.4) -- (2.4,0.6) -- (1.6,1.4) -- (0.8,1.2) -- (0.4,2.0) -- (1.0,2.8) -- (0.8,3.8) -- (1.4,4.6) -- (2.0,4.8) -- cycle;
% Aires protégées
\fill[popgreen!70] (3.9,5.3) circle (0.28);
\fill[popgreen!70] (3.6,4.3) circle (0.28);
\fill[popgreen!70] (1.1,2.2) circle (0.28);
\fill[popgreen!70] (4.6,1.4) circle (0.28);
\fill[popgreen!70] (2.0,0.9) circle (0.28);
% Légendes pointées
\draw[thin, popmuted] (3.9,5.3) -- (5.6,5.9) node[right, popdark]{Waza (Extrême-Nord)};
\draw[thin, popmuted] (3.6,4.3) -- (5.9,4.3) node[right, popdark]{Bénoué (Nord)};
\draw[thin, popmuted] (1.1,2.2) -- (-0.4,2.9) node[left, popdark]{Korup (Sud-Ouest)};
\draw[thin, popmuted] (4.6,1.4) -- (6.0,1.7) node[right, popdark]{Dja (Sud-Est)};
\draw[thin, popmuted] (2.0,0.9) -- (0.2,0.1) node[left, popdark]{Campo-Ma'an (Sud)};
\node[popmuted, font=\small\itshape] at (3.0,6.9) {Cameroun};
\end{tikzpicture}
\end{center}
\legende{Localisation schématique des principales aires protégées du Cameroun : parcs nationaux de Waza, de la Bénoué, de Korup et de Campo-Ma'an, et réserve de faune du Dja. Le contour du pays est volontairement simplifié.}
\end{popfigure}

La \textbf{lutte contre le \cle{braconnage}} complète ce dispositif : écogardes armés, patrouilles, répression du commerce illégal d'ivoire et de viande de brousse, et implication des populations riveraines, qui tirent des revenus de l'écotourisme et ont donc intérêt à protéger la faune.

\courssub{Le cadre institutionnel camerounais}

La gestion durable des ressources repose sur des institutions et des lois :

\begin{itemize}
\item le \textbf{MINEPDED} (Ministère de l'Environnement, de la Protection de la Nature et du Développement Durable) élabore et applique la politique environnementale ;
\item le \textbf{MINFOF} (Ministère des Forêts et de la Faune) gère les forêts, la faune et les aires protégées ;
\item la \textbf{loi n° 96/12 du 5 août 1996} portant loi-cadre relative à la gestion de l'environnement fixe les grands principes : prévention des pollutions, études d'impact obligatoires avant les grands projets, droit de chacun à un environnement sain ;
\item les \textbf{aires protégées} couvrent aujourd'hui environ \textbf{20 \% du territoire national}, un des taux les plus élevés d'Afrique.
\end{itemize}

\courssub{Le rôle des populations locales et de l'éducation environnementale}

Aucune politique ne réussit sans les populations. Les communautés villageoises gèrent les forêts communautaires, participent aux comités de surveillance des parcs, pratiquent l'agroforesterie et signalent les braconniers. L'\cle{éducation environnementale} — à l'école, dans les clubs de jeunes, par les médias — forme des citoyens capables de comprendre les enjeux et d'agir : trier ses déchets, éviter les feux tardifs, planter des arbres, refuser la viande d'espèces protégées. Tu es toi-même un acteur de la gestion durable.

\begin{methode}[Proposer une mesure de gestion durable]
Face à une situation d'exploitation d'une ressource, procède en trois étapes :
\begin{enumerate}
\item \textbf{Identifie la pression} exercée sur la ressource (surpêche, déforestation, pollution d'une nappe, braconnage...) ;
\item \textbf{Choisis une mesure adaptée} qui agit sur cette pression selon l'une des trois stratégies : \emph{réduire le prélèvement} (quotas, coupe sélective, forage raisonné), \emph{restaurer la ressource} (reboisement, jachère, épuration des eaux), ou \emph{protéger le milieu} (aire protégée, pare-feux, loi) ;
\item \textbf{Justifie} ta proposition en montrant qu'elle concilie conservation de la ressource et besoins des populations (revenus, alimentation, emplois).
\end{enumerate}
\end{methode}

\begin{exoresolu}[Analyser une situation locale]
Énoncé. Dans un village de la région de l'Est, une entreprise exploite une concession forestière en abattant tous les arbres des parcelles, sans replanter. Les habitants constatent la disparition du gibier, l'ensablement de la rivière et la baisse des pluies locales. Propose trois mesures de gestion durable adaptées, en les justifiant.
\tcblower
Solution détaillée.
\begin{enumerate}
\item \textbf{Imposer un plan d'aménagement avec coupes sélectives} : la pression identifiée est une coupe rase qui détruit la forêt. La coupe sélective (arbres matures uniquement, rotation de 25 à 30 ans) réduit le prélèvement et laisse la forêt se régénérer ; la faune conserve son habitat.
\item \textbf{Reboiser les parcelles exploitées} : mesure de restauration. Planter des essences locales reconstitue le couvert forestier, ce qui limite l'érosion responsable de l'ensablement de la rivière et restaure le cycle local de l'eau.
\item \textbf{Associer les populations} (forêt communautaire voisine, emplois dans le reboisement, comité de suivi) : les habitants, tirant des revenus de la forêt, deviennent acteurs de sa protection au lieu d'en subir la destruction.
\end{enumerate}
Ces trois mesures illustrent les trois stratégies de la méthode : réduire le prélèvement, restaurer la ressource, protéger le milieu en impliquant les populations.
\end{exoresolu}

\begin{attention}[Pièges fréquents au Probatoire]
\begin{itemize}
\item Ne confonds pas \textbf{parc national} (protection intégrale) et \textbf{réserve} (activités encadrées possibles) : ce sont deux catégories d'aires protégées.
\item Une forêt communautaire n'est pas une forêt « libre d'accès » : elle est gérée selon un plan approuvé par l'État.
\item Ne cite pas la Grande Muraille Verte comme un projet camerounais isolé : c'est une initiative panafricaine à laquelle le Cameroun participe.
\item Dans une proposition de mesure, justifie toujours le lien entre la pression et la mesure : une mesure non justifiée vaut peu de points.
\end{itemize}
\end{attention}

\begin{aretenir}[L'essentiel de la section]
\begin{itemize}
\item Eau : protection des bassins versants, traitement des eaux usées, forages raisonnés, restauration du lac Tchad.
\item Sols : agroforesterie, courbes de niveau, jachères, lutte contre les feux de brousse, Grande Muraille Verte.
\item Forêts : forêts communautaires, plans d'aménagement avec coupes sélectives, reboisement, certification.
\item Biodiversité : aires protégées (Waza, Bénoué, Korup, Campo-Ma'an, Dja — \SI{526000}{ha}, patrimoine UNESCO), lutte contre le braconnage.
\item Cadre institutionnel : MINEPDED, MINFOF, loi-cadre de 1996, environ 20 \% du territoire en aires protégées.
\item La gestion durable ne réussit qu'avec la participation des populations et l'éducation environnementale.
\end{itemize}
\end{aretenir}

\coursec{Analyse de dossiers documentaires et synthèse}

Après avoir exploré les principes de la gestion durable et les mesures concrètes mises en œuvre au Cameroun, nous abordons maintenant une compétence clé du programme : \textbf{l'analyse critique de dossiers documentaires}. Au Probatoire, vous serez souvent confronté à un ensemble de documents (textes, tableaux, cartes, graphiques) décrivant une situation d'exploitation d'une \cle{ressource naturelle}. Il ne s'agit pas de les recopier, mais d'en extraire l'information pertinente pour \textbf{identifier les pressions, expliquer les conséquences et proposer des mesures de gestion durable argumentées}. Cette section vous donne la méthode, l'applique sur trois grands cas camerounais et synthétise toute la leçon.

\courssub{Méthodologie d'analyse d'un dossier documentaire}

\begin{methode}[Analyser un dossier documentaire en 4 étapes]
\begin{enumerate}
    \item \textbf{Lire la consigne et repérer la question centrale.} Souligne les mots-clés : « \cle{pression} », « \cle{conséquence} », « \cle{mesure de gestion} », « \cle{acteur} », « \cle{développement durable} ». Identifie la ressource concernée (eau, sol, forêt, biodiversité).
    \item \textbf{Extraire les données chiffrées et qualitatives de chaque document.} Pour chaque document (texte, tableau, carte, photo), note : les chiffres (superficies, volumes, populations, dates), les acteurs (État, entreprises, ONG, populations locales, braconniers), les conflits d'usage, les lois citées.
    \item \textbf{Relier l'activité humaine à l'impact sur la ressource.} Construis la chaîne causale : \emph{Activité} $\rightarrow$ \emph{Pression} (prélèvement, pollution, \cle{fragmentation}) $\rightarrow$ \emph{Conséquence} (érosion, perte de biodiversité, assèchement, appauvrissement des populations). Utilise les flèches pour visualiser les liens.
    \item \textbf{Proposer une mesure de gestion durable justifiée.} La mesure doit répondre aux trois piliers : \textbf{écologique} (préserver la ressource), \textbf{économique} (maintenir une activité rentable), \textbf{social} (respecter les droits des populations locales). Cite une loi, une convention ou un exemple local (\cle{Plan d'aménagement}, \cle{UFA}, \cle{AP}, \cle{reboisement}, \cle{quota}, \cle{écotourisme}).
\end{enumerate}
\end{methode}

\begin{attention}[Piège fréquent : recopier au lieu d'analyser]
Ne recopie \textbf{jamais} des paragraphes entiers du dossier. L'examinateur attend une \textbf{synthèse personnelle} : « Le document 2 montre que le volume de bois prélevé est passé de 2,5 millions de m³ en 2000 à 3,8 millions de m³ en 2015, ce qui dépasse le taux de renouvellement naturel estimé à 1,5 m³/ha/an (Doc 3). Cette \cle{surexploitation} entraîne une dégradation de la strate dominante… » C'est cette mise en relation de chiffres issus de documents différents qui vaut des points.
\end{attention}

\courssub{Étude de cas 1 : L'exploitation forestière dans le Sud-Est Cameroun}

\begin{experience}[Dossier : Forêt dense humide du Sud-Est (région du Dja-et-Lobo)]
\textbf{Contexte.} La forêt du bassin du Congo couvre environ 22 millions d'hectares au Cameroun. Le Sud-Est abrite des \cle{Unités Forestières d'Aménagement (UFA)} attribuées à des entreprises (souvent étrangères) et des \cle{Aires Protégées (AP)} (Parc du Dja, Réserve du Dja). 

\textbf{Documents types (simulés pour l'exercice).}
\begin{itemize}
    \item \textbf{Doc 1 (Carte) :} Localisation des UFA (zones d'exploitation) et des Aires Protégées. On observe une \cle{fragmentation} : les UFA découpent la forêt, isolant les AP.
    \item \textbf{Doc 2 (Tableau) :} Exportations de grumes et sciages (2010–2020). Volume grumes : 2,1 M m³ (2010) $\rightarrow$ 3,4 M m³ (2020). Recettes : 150 milliards FCFA/an. Emplois directs : 12 000.
    \item \textbf{Doc 3 (Texte scientifique) :} « Le \cle{diamètre minimum d'exploitation (DME)} pour l'\cle{ayous} est de 60 cm. Le \cle{cycle de rotation} de 30 ans est insuffisant pour reconstituer le \cle{stock commercial}. L'ouverture de pistes forestières facilite le \cle{braconnage} et l'agriculture sur brûlis. »
    \item \textbf{Doc 4 (Témoignage ONG) :} « Les populations \cle{Baka} et \cle{Bantu} perdent l'accès aux \cle{produits forestiers non ligneux (PFNL)} : chenilles, miel, écorces médicamenteuses et aux sites sacrés. Les \cle{redevances forestières annuelles (RFA)} versées aux communes sont souvent détournées. »
\end{itemize}

\textbf{Analyse guidée.}
\begin{enumerate}
    \item \textbf{Ressource :} Forêt dense humide (bois d'œuvre, biodiversité, PFNL, \cle{services écosystémiques} : carbone, climat).
    \item \textbf{Activité / Pression :} Exploitation industrielle sélective (mais intensive) + pistes d'évacuation + chasse commerciale facilitée par les pistes.
    \item \textbf{Conséquences écologiques :} Fragmentation de l'habitat (isolement des populations de grands mammifères : gorille, éléphant de forêt), perte de diversité génétique, perturbation du cycle de l'eau, érosion sur pistes.
    \item \textbf{Conséquences socio-économiques :} Revenus d'exportation et emplois (pilier économique) MAIS inéquité du partage des bénéfices (RFA détournées), perte des moyens de subsistance pour autochtones (pilier social), risque d'épuisement du stock commercial à 30 ans (pilier écologique non respecté).
    \item \textbf{Mesures de gestion durable proposées :}
    \begin{itemize}
        \item Respect strict du \cle{Plan d'Aménagement (PA)} : augmentation du DME, allongement de la rotation à 60 ans, réduction de l'\cle{assiette de coupe annuelle}.
        \item Certification \cle{FSC} (Forest Stewardship Council) obligatoire pour l'export vers l'UE (réglementation \cle{FLEGT}).
        \item Création de \cle{corridors écologiques} entre AP (ex. corridor Dja-Lobéké) pour contrer la fragmentation.
        \item \cle{Gestion communautaire} des forêts riveraines (Forêts Communautaires) pour les PFNL et écotourisme, avec contrats de partage de recettes transparents.
    \end{itemize}
\end{enumerate}
\end{experience}

\begin{exoresolu}[Question type Probatoire : Exploitation forestière]
\textbf{Énoncé.} À partir des documents 1 à 4, montrez que l'exploitation forestière actuelle dans le Sud-Est n'est pas durable et proposez deux mesures concrètes pour la rendre compatible avec le développement durable.

\textbf{Solution détaillée.}
\textbf{1. Non-durabilité (argumentation croisée documents) :}
\begin{itemize}
    \item \textbf{Écologique :} Doc 3 indique un cycle de 30 ans insuffisant pour la reconstitution du stock (diamètre 60 cm non atteint). Doc 1 montre la fragmentation par les pistes/UFA $\rightarrow$ perte de connectivité pour la mégafaune.
    \item \textbf{Économique :} Doc 2 montre une hausse des volumes exportés (+60\% en 10 ans) non compensée par une transformation locale (peu de sciage/usinage sur place) $\rightarrow$ export de \cle{valeur brute}, vulnérabilité aux cours mondiaux.
    \item \textbf{Social :} Doc 4 : populations locales exclues des bénéfices (RFA détournées) et privées de ressources (PFNL, sites culturels).
\end{itemize}
\textbf{2. Mesures proposées (justifiées) :}
\begin{enumerate}
    \item \textbf{Allonger la rotation à 60 ans et augmenter le DME} (ex. 80 cm pour l'ayous, 100 cm pour le \cle{sapelli}) \emph{→} permet la reconstitution du capital bois (pilier écologique) et garantit l'approvisionnement futur de l'industrie (pilier économique). Base légale : Loi 94/01 régime des forêts, décret d'application sur les PA.
    \item \textbf{Instaurer une transformation locale obligatoire (quota de sciage sur place 70\%) et un fonds de développement communautaire alimenté par 50\% des RFA} \emph{→} crée de l'emploi qualifié local (économique), réduit l'export de grumes (écologique : moins de pistes/lourdes infrastructures), redistribue la rente aux populations (social). Exemple : modèle de la forêt communautaire de Nkolbisson (région Centre) ou conventions UFA-riverains.
\end{enumerate}
\end{exoresolu}

\courssub{Étude de cas 2 : L'assèchement du lac Tchad et ses conséquences humaines}

\begin{experience}[Dossier : Bassin du lac Tchad (Extrême-Nord Cameroun)]
\textbf{Contexte.} Le lac Tchad, partagé entre Cameroun, Tchad, Niger, Nigeria, a perdu \textbf{90\% de sa surface} depuis les années 1960 (25 000 km² $\rightarrow$ ~2 500 km² aujourd'hui, avec fluctuations saisonnières).

\textbf{Documents types.}
\begin{itemize}
    \item \textbf{Doc 1 (Série chronologique d'images satellites) :} Régression nette de l'eau libre, apparition de zones sableuses et de végétation herbacée (nouvelles terres agricoles).
    \item \textbf{Doc 2 (Graphique climatique) :} Précipitations moyennes annuelles à N'Djamena : 320 mm (années 1950) $\rightarrow$ 220 mm (années 1980) $\rightarrow$ 300 mm (années 2010). Variabilité interannuelle forte.
    \item \textbf{Doc 3 (Tableau usages de l'eau) :} Prélèvements pour irrigation (riz, maïs) : 2,5 km³/an (1980) $\rightarrow$ 5,8 km³/an (2020). Population du bassin : 17 millions (1990) $\rightarrow$ 45 millions (2020).
    \item \textbf{Doc 4 (Rapport OMS/UNICEF) :} Taux de malnutrition aiguë > 15\% dans la zone (seuil d'urgence). Conflits éleveurs/agriculteurs/pêcheurs pour l'accès à l'eau et aux terres émergées. Déplacements de populations (réfugiés climatiques + insécurité Boko Haram).
\end{itemize}

\textbf{Analyse guidée.}
\begin{enumerate}
    \item \textbf{Ressource :} Eau douce du lac Tchad (pêche, irrigation, eau de boisson, pâturages, biodiversité : oiseaux migrateurs, hippopotames, lamantins).
    \item \textbf{Pressions combinées :}
    \begin{itemize}
        \item \textbf{Climatique :} Sécheresses récurrentes (variabilité naturelle + \cle{changement climatique anthropique}) $\rightarrow$ baisse du bilan hydrique (P < \cle{ETP}).
        \item \textbf{Anthropique directe :} Surprélèvement pour irrigation amont (barrages sur Chari/Logone : Maga, Yagoua) $\rightarrow$ réduction de l'apport au lac. Démographie galopante $\rightarrow$ pression sur terres émergées (agriculture de décrue non régulée).
    \end{itemize}
    \item \textbf{Conséquences :} Effondrement de la pêche (production divisée par 10), \cle{salinisation} des sols par évaporation, perte de biodiversité (disparition de l'hippopotame sur la partie camerounaise), conflits violents pour l'eau, insécurité alimentaire majeure, radicalisation facilitée par la misère.
    \item \textbf{Mesures de gestion durable (échelle locale à régionale) :}
    \begin{itemize}
        \item \textbf{Local (Cameroun, Extrême-Nord) :} Régulation des prélèvements irrigation (\cle{quotas}, \cle{goutte-à-goutte}), restauration des berges (végétation fixatrice : \emph{Acacia nilotica}, \emph{Balanites aegyptiaca}), plans de gestion concertée des pêcheries (périodes de repos biologique, mailles réglementaires), développement de l'\cle{aquaculture} hors lac.
        \item \textbf{Régional (Commission du Bassin du Lac Tchad - CBLT) :} Projet de transfert d'eau du bassin de l'Oubangui (Congo) vers le Chari (projet \cle{Transaqua}) — \emph{controversé : coût énorme, impact écologique amont, géopolitique}. Alternative prioritaire : \cle{Gestion Intégrée des Ressources en Eau (GIRE)} au niveau du bassin, accord de partage de l'eau entre États riverains, \cle{système d'alerte précoce} sécheresses.
    \end{itemize}
\end{enumerate}
\end{experience}

\begin{aretenir}[Lac Tchad : le piège du « tout climatique »]
Au Probatoire, ne dites pas « c'est à cause du changement climatique, on ne peut rien faire ». Les documents montrent \textbf{deux leviers d'action locaux} : (1) la gestion de la demande en eau (irrigation efficiente, cultures moins gourmandes), (2) la gouvernance locale des ressources (comités de gestion de l'eau, résolution des conflits). C'est ce que le jury attend : \textbf{des solutions à l'échelle de l'acteur local}, même si le problème est global.
\end{aretenir}

\courssub{Étude de cas 3 : Le braconnage des éléphants dans le parc national de la Bénoué}

\begin{experience}[Dossier : Parc National de la Bénoué (Région Nord, 180 000 ha)]
\textbf{Contexte.} Zone soudano-sahélienne, écosystème de savane boisée. Espèces phares : éléphant de savane (\emph{Loxodonta africana}), girafe, hippotrague, buffle, lion.

\textbf{Documents types.}
\begin{itemize}
    \item \textbf{Doc 1 (Recensements aériens MINFOF/WWF) :} Population éléphants : $\sim$ 5 000 (années 1970) $\rightarrow$ 1 200 (2000) $\rightarrow$ \textbf{200–300} (2020). Taux de braconnage estimé > 8\%/an (seuil critique).
    \item \textbf{Doc 2 (Saisies douanières / ONG TRAFFIC) :} Ivoire saisi au port de Douala : 1,5 tonnes (2015), 800 kg (2018). Origine : Nord Cameroun, transit par Garoua/Ngaoundéré. Prix ivoire brut : 150 000–200 000 FCFA/kg (marché noir local) $\rightarrow$ 1 000 000 FCFA/kg (marché asiatique).
    \item \textbf{Doc 3 (Témoignages éco-gardes) :} Effectifs : 40 éco-gardes pour 180 000 ha (1/4500 ha, norme UICN : 1/500–1000 ha). Équipement : véhicules vieux, pas de GPS, pas de drones, armes légères vs braconniers armés AK-47, motos. Corruption : certains gardes complices.
    \item \textbf{Doc 4 (Enquête socio-économique riverains) :} 60\% des ménages dépendent de la chasse villageoise (protéines). Conflits homme-éléphant : destruction champs (mil, maïs, coton) $\rightarrow$ tueries de représailles. Revenus tourisme : quasi nuls (insécurité, pistes impraticables).
\end{itemize}

\textbf{Analyse guidée.}
\begin{enumerate}
    \item \textbf{Ressource :} Biodiversité (éléphant = \cle{espèce parapluie}, \cle{ingénieur d'écosystème} : dispersion graines, ouverture milieux, points d'eau).
    \item \textbf{Pressions :}
    \begin{itemize}
        \item \textbf{Braconnage commercial organisé :} Réseaux transfrontaliers (Tchad, RCA, Soudan) $\rightarrow$ ivoire $\rightarrow$ Asie. Facteur déclencheur : prix élevé, impunité, corruption.
        \item \textbf{Braconnage de subsistance / conflits H-E :} Chasse villageoise pour viande de brousse + tueries protection cultures.
        \item \textbf{Insuffisance des moyens de protection :} Sous-effectif, sous-équipement, manque de formation, gouvernance faible.
    \end{itemize}
    \item \textbf{Conséquences :} Effondrement démographique (risque d'extinction locale), perte de la fonction écologique (régénération forestière compromise), perte économique (tourisme de vue impossible), insécurité (bandes armées), appauvrissement protéique populations.
    \item \textbf{Mesures de gestion durable (approche intégrée) :}
    \begin{itemize}
        \item \textbf{Renforcement application loi :} Recrutement/formation 100 éco-gardes, équipements (drones, GPS, VHF, véhicules 4x4), patrouilles mixtes (MINFOF + BIR + douanes), lutte anti-corruption (rotation postes, primes de résultat).
        \item \textbf{Gestion des conflits H-E :} Clôtures électriques légères (essai parc Waza), cultures répulsives (piment, sésame, citronnelle) en bordure, indemnisation rapide dégâts (fonds assurance indiciel), couloirs de migration sécurisés vers parc Faro et Bouba Ndjidda.
        \item \textbf{Valorisation économique pour riverains :} Partage recettes tourisme (redevances entrées parc $\rightarrow$ micro-projets : forages, écoles, santé), emploi éco-gardes locaux, développement filières PFNL (karité, baobab, miel) + label « produit ami de l'éléphant ».
        \item \textbf{Coopération transfrontalière :} Accord tripartite Cameroun-Tchad-RCA pour patrouilles conjointes, harmonisation législative (peines prison ferme ivoire), échange renseignements.
    \end{itemize}
\end{enumerate}
\end{experience}

\begin{exemplebox}[Chiffres clés : L'effondrement des éléphants de la Bénoué]
\begin{center}
\begin{tabular}{|c|c|c|}\hline
\textbf{Année} & \textbf{Effectif estimé} & \textbf{Tendance} \\ \hline
1970 (premiers recensements) & $\sim$ 5 000 & Référence \\ \hline
2000 & $\sim$ 1 200 & -76\% en 30 ans \\ \hline
2010 & $\sim$ 600 & -50\% en 10 ans \\ \hline
2020 (dernier recensement aérien) & \textbf{200–300} & \textbf{-95\% en 50 ans} \\ \hline
\end{tabular}
\end{center}
\legende{Source : synthèse MINFOF / WWF / UICN. Le taux de braconnage (>8\%/an) dépasse largement le taux d'accroissement naturel de l'espèce (~5–6\%/an), conduisant à l'extinction quasi certaine sans action drastique.}
\end{exemplebox}

\courssub{Construction d'un tableau bilan : Ressource / Pression / Conséquence / Mesure de gestion durable}

La synthèse des trois études de cas (et des autres ressources vues en section 4) permet de construire un tableau récapitulatif, outil indispensable pour réviser et répondre aux questions de synthèse du Probatoire.

\begin{popfigure}
\begin{center}
\begin{tikzpicture}[
    cell/.style={rectangle, draw=popline, minimum height=1.1cm, align=center, text width=3.2cm, font=\small},
    header/.style={cell, fill=popblue, text=white, font=\small\bfseries},
    rowA/.style={cell, fill=popblueL!30},
    rowB/.style={cell, fill=popgreenL!30},
    rowC/.style={cell, fill=poporangeL!30},
    rowD/.style={cell, fill=poppurpleL!30},
]
% En-têtes
\node[header] (h1) at (0,0) {Ressource naturelle};
\node[header] (h2) at (3.4,0) {Pression(s) anthropique(s) principale(s)};
\node[header] (h3) at (6.8,0) {Conséquence(s) majeure(s)};
\node[header] (h4) at (10.2,0) {Mesure(s) de gestion durable (ex. Cameroun)};

% Ligne 1 : Forêt
\node[rowA] (r11) at (0,-1.1) {Forêt dense humide (Sud-Est) : bois d'œuvre, biodiversité, PFNL, carbone};
\node[rowA] (r12) at (3.4,-1.1) {Exploitation industrielle intensive (UFA), pistes, chasse commerciale facilitée, agriculture sur brûlis};
\node[rowA] (r13) at (6.8,-1.1) {Fragmentation habitat, perte stock commercial, érosion, appauvrissement populations autochtones (Baka/Bantu)};
\node[rowA] (r14) at (10.2,-1.1) {PA respectés (rotation 60 ans, DME majoré), certification FSC/FLEGT, corridors écologiques, forêts communautaires, transparence RFA};

% Ligne 2 : Eau / Lac Tchad
\node[rowB] (r21) at (0,-2.2) {Eau douce Lac Tchad (Extrême-Nord) : pêche, irrigation, pâture, biodiversité};
\node[rowB] (r22) at (3.4,-2.2) {Sécheresses climatiques + surprélèvements irrigation amont (barrages Maga/Yagoua) + démographie};
\node[rowB] (r23) at (6.8,-2.2) {Assèchement 90\%, effondrement pêche, salinisation, conflits violents, insécurité alimentaire, réfugiés};
\node[rowB] (r24) at (10.2,-2.2) {GIRE bassin (CBLT), quotas irrigation + goutte-à-goutte, cultures peu gourmandes, restauration berges, aquaculture, accord partage eau};

% Ligne 3 : Biodiversité / Éléphant
\node[rowC] (r31) at (0,-3.3) {Biodiversité savane (Parc Bénoué) : éléphant (espèce parapluie), grands mammifères};
\node[rowC] (r32) at (3.4,-3.3) {Braconnage ivoire organisé (réseaux transfrontaliers), conflits H-E (destruction cultures), sous-effectif éco-gardes, corruption};
\node[rowC] (r33) at (6.8,-3.3) {Chute population -95\% (5000 $\rightarrow$ 200-300), perte fonction écologique, insécurité, perte revenus tourisme, appauvrissement};
\node[rowC] (r34) at (10.2,-3.3) {Renforcement effectifs/équipements éco-gardes, patrouilles mixtes, clôtures/cultures répulsives, couloirs migration, partage recettes tourisme, coopération transfrontalière};

% Ligne 4 : Sols (rappel section 4)
\node[rowD] (r41) at (0,-4.4) {Sols (zones cotonnières Nord, périmètres irrigués) : fertilité, structure};
\node[rowD] (r42) at (3.4,-4.4) {Monoculture coton (intrants chimiques), surpâturage, déforestation, irrigation mal drainée (salinisation)};
\node[rowD] (r43) at (6.8,-4.4) {Érosion hydrique/éolienne, perte matière organique, salinisation/sodification, baisse rendements, désertification};
\node[rowD] (r44) at (10.2,-4.4) {Agroforesterie (Faidherbia), rotation cultures, jachère améliorée, compost/fumure, drainage périmètres, CES/DRS, agriculture de conservation};
\end{tikzpicture}
\end{center}
\legende{Tableau bilan des quatre ressources majeures étudiées : identification de la pression dominante, de la conséquence critique et de mesures de gestion durable concrètes et contextualisées au Cameroun. Ce tableau est une « fiche de révision » idéale pour le Probatoire.}
\end{popfigure}

\courssub{Synthèse générale : De l'identification des impacts à la proposition de solutions locales}

\begin{popfigure}
\begin{center}
\begin{tikzpicture}[
    mindmap,
    concept color=popblue,
    level 1 concept/.append style={level distance=4.5cm, sibling angle=90, font=\small\bfseries},
    level 2 concept/.append style={level distance=3.2cm, sibling angle=45, font=\small},
    level 3 concept/.append style={level distance=2.5cm, sibling angle=30, font=\footnotesize},
    every node/.style={concept, execute at begin node=\hskip0pt},
    text=white,
    grow cyclic,
]
\node[concept, minimum size=2.2cm, align=center] (center){RESSOURCES NATURELLES\\CAMEROUNAISES\\Gestion Durable}
    child[concept color=popgreen] {
        node[concept, minimum size=1.8cm] {EAU}
        child { node {Lac Tchad : GIRE, quotas, restauration berges} }
        child { node {Barrages : sédimentation, migration poissons} }
        child { node {Pollution urbaine/industrielle} }
    }
    child[concept color=poporange] {
        node[concept, minimum size=1.8cm] {SOLS}
        child { node {Agroforesterie, CES/DRS, rotation} }
        child { node {Lutte salinisation (drainage, lessivage)} }
        child { node {Régénération naturelle assistée (RNA)} }
    }
    child[concept color=poppurple] {
        node[concept, minimum size=1.8cm] {FORÊTS}
        child { node {UFA : PA respectés, certification FSC} }
        child { node {Corridors écologiques, AP connectées} }
        child { node {Forêts communautaires, PFNL, écotourisme} }
    }
    child[concept color=poppink] {
        node[concept, minimum size=1.8cm] {BIODIVERSITÉ}
        child { node {Anti-braconnage : effectifs, tech, loi} }
        child { node {Conflits H-E : clôtures, couloirs, indemnisation} }
        child { node {Valorisation économique locale (tourisme, labels)} }
    };
% Liens transversaux (pressions communes)
\node[draw=popdark, thick, dashed, fit=(center), label={[font=\tiny\bfseries, text=popdark]above:PRESSIONS COMMUNES : Démographie • Changement climatique • Gouvernance faible • Pauvreté • Marchés illicites}] {};
\end{tikzpicture}
\end{center}
\legende{Carte mentale de synthèse de la leçon ND07. Au centre, les quatre ressources. Sur chaque branche, les mesures de gestion durable spécifiques (niveau 2) et les leviers d'action (niveau 3). Le cadre pointillé rappelle les pressions transversales qui exigent des réponses systémiques (politiques publiques, éducation, coopération internationale).}
\end{popfigure}

\begin{aretenir}[Synthèse pour le Probatoire : 3 idées forces]
\begin{enumerate}
    \item \textbf{Diagnostic partagé :} Toute ressource (eau, sol, forêt, biodiversité) subit des \textbf{pressions multiples} (climatiques + anthropiques directes) qui s'additionnent. L'analyse documentaire doit savoir les \textbf{dissocier} (ex. : lac Tchad = climat + irrigation) pour proposer des mesures ciblées.
    \item \textbf{Gestion durable = triangle équilibré :} Une mesure n'est durable que si elle est \textbf{écologiquement efficace} (ressource préservée), \textbf{économiquement viable} (activités maintenues/transformées) et \textbf{socialement équitable} (populations associées, bénéfices partagés). Citez toujours les trois piliers.
    \item \textbf{L'échelle locale est le levier d'action :} Même pour des problèmes globaux (climat, trafic ivoire), le Probatoire attend des \textbf{mesures concrètes à l'échelle communale, départementale ou d'aire protégée} : plan d'aménagement, comité de gestion, quota, reboisement, culture répulsive, écogarde local, forêt communautaire.
\end{enumerate}
\end{aretenir}

\begin{savaistu}[Le Probatoire : ce qui tombe (presque) tous les ans]
\begin{itemize}
    \item \textbf{Question 1 (Analyse de document) :} Un dossier 3–4 documents (carte évolution forêt, tableau volumes bois, texte loi forestière, témoignage ONG). \textbf{Consigne type :} « Montrez que l'exploitation n'est pas durable. Proposez deux mesures. » $\rightarrow$ Appliquez la \textbf{Méthode 4 étapes} ci-dessus.
    \item \textbf{Question 2 (Mise en œuvre) :} « Vous êtes chef de poste MINFOF dans une zone de conflit homme-éléphant. Décrivez la démarche participative pour réduire les dégâts aux cultures. » $\rightarrow$ Réponse : Diagnostic participatif $\rightarrow$ Choix mesures (clôture piment, couloirs) $\rightarrow$ Convention signée $\rightarrow$ Suivi-évaluation $\rightarrow$ Indemnisation.
    \item \textbf{Question 3 (Synthèse) :} « En vous appuyant sur un exemple camerounais, expliquez comment la gestion communautaire d'une ressource naturelle contribue au développement durable. » $\rightarrow$ Exemples types : Forêt communautaire (bois/PFNL), Plan de gestion pêche lac Lagdo, Comité de gestion eau villageois. Structurez : Acteurs $\rightarrow$ Règles $\rightarrow$ Bénéfices 3 piliers.
\end{itemize}
\end{savaistu}

\begin{exercice}[Entraînement Probatoire : Dossier « Forêt communautaire de Ngoyla »]
\textbf{Documents fournis (imaginés) :}
\begin{itemize}
    \item \textbf{Doc A :} Décret 95/531/PM fixant les modalités d'attribution des forêts communautaires.
    \item \textbf{Doc B :} Carte : Forêt communautaire de 5 000 ha, limitrophe d'une UFA et d'un parc national. Zonage : zone de production (3 000 ha), zone de conservation (1 500 ha), zone d'agroforesterie (500 ha).
    \item \textbf{Doc C :} Tableau : Revenus annuels communauté (2018–2022). Vente bois (grumes) : 15 M FCFA/an. Vente PFNL (miel, Njansang (Djansang), moabi) : 3 M FCFA/an. Écotourisme (guide, hébergement) : 0,5 M FCFA/an. Parts reversées aux ménages : 40\%. Frais de gestion (inventaire, garde) : 30\%. Dette envers entreprise d'exploitation : 20 M FCFA.
    \item \textbf{Doc D :} Extrait rapport ONG : « L'inventaire de 2021 montre un taux de prélèvement supérieur au taux de reconstitution pour l'essence Moabi. Les pistes d'évacuation traversent la zone de conservation. Les jeunes émigrent faute d'emplois qualifiés. »
\end{itemize}
\textbf{Questions :}
\begin{enumerate}
    \item À partir des documents B et C, identifiez deux pressions exercées sur la forêt communautaire et leurs conséquences écologiques et socio-économiques.
    \item Le document D révèle des dysfonctionnements dans la gestion. Citez-en deux et proposez une mesure corrective pour chacun.
    \item En vous appuyant sur l'ensemble du dossier, arguez pour ou contre le renouvellement de l'accord de gestion communautaire à l'échéance 2025. Structurez votre réponse autour des trois piliers du développement durable.
\end{enumerate}
\end{exercice}

\begin{corrige}[Exercice : Forêt communautaire de Ngoyla]
\textbf{1. Pressions et conséquences :}
\begin{itemize}
    \item \textbf{Pression 1 (Doc B + C) :} Exploitation bois grumes (zone de production 3 000 ha) $\rightarrow$ \textbf{Conséquence écologique :} prélèvement > reconstitution pour Moabi (Doc D) $\rightarrow$ risque épuisement stock, perte diversité génétique. \textbf{Socio-éco :} revenus bois (15 M/an) mais dette 20 M $\rightarrow$ dépendance vis-à-vis exploitant, faible valeur ajoutée locale (pas de transformation).
    \item \textbf{Pression 2 (Doc B + D) :} Pistes d'évacuation traversant zone de conservation $\rightarrow$ \textbf{Conséquence écologique :} fragmentation habitat, perturbation faune, facilitation braconnage. \textbf{Socio-éco :} zone conservation (1 500 ha) non respectée $\rightarrow$ perte services écosystémiques (eau, carbone, tourisme potentiel).
\end{itemize}
\textbf{2. Dysfonctionnements et mesures correctives :}
\begin{itemize}
    \item \textbf{Dysfonctionnement 1 :} Taux de prélèvement Moabi non durable (Doc D). \textbf{Mesure :} Réviser le \cle{Plan d'Aménagement Simple (PAS)} : baisser l'assiette de coupe annuelle Moabi, augmenter DME, introduire rotation 60 ans, enrichissement en Moabi (plantations d'enrichissement).
    \item \textbf{Dysfonctionnement 2 :} Non-respect zonage (pistes en zone conservation) + endettement + pas d'emplois qualifiés jeunes. \textbf{Mesure :} Renégocier contrat avec exploitant : clause « transformation sur place (scierie mobile) = emplois jeunes », clause « respect strict zonage = pistes hors zone conservation », plan d'apurement dette sur 5 ans prélevé sur recettes, formation jeunes (garde, guide, transformation PFNL).
\end{itemize}
\textbf{3. Avis sur renouvellement 2025 (argumenté 3 piliers) :}
\textbf{POUR le renouvellement (SI conditions) :}
\begin{itemize}
    \item \textbf{Écologique :} Zonage existant (conservation 30\%) protège cœur de forêt. Mesures correctives (ci-dessus) permettent durabilité Moabi. Connectivité avec parc national maintenue.
    \item \textbf{Économique :} Revenus réels (18,5 M/an) mais mal répartis. Transformation locale + écotourisme (potentiel parc limitrophe) + PFNL (Njansang, miel) = diversification revenus, réduction dépendance grumes, création emplois qualifiés jeunes.
    \item \textbf{Social :} 40\% ménages = base. Il faut : assemblée générale transparente (comptes), fonds de développement communautaire (école, santé, eau) alimenté par recettes, inclusion femmes/jeunes dans instances décision.
\end{itemize}
\textbf{CONCLUSION :} Renouvellement \textbf{conditionnel} à la signature d'un nouveau PAS intégrant les mesures correctives, un plan d'affaires (business plan) communautaire 5 ans, et un contrat de partenariat équilibré avec l'exploitant (ou reprise en régie communautaire avec appui technique MINFOF/ONG). C'est cela, la gestion durable : \textbf{adapter, corriger, partager}.
\end{corrige}

\courssub{Bilan de la leçon ND07 : Compétences validées}

\begin{aretenir}[Check-list de révision — Leçon ND07]
\begin{itemize}
    \item \textbf{Savoirs :} Je sais définir \cle{ressource naturelle}, \cle{développement durable}, \cle{gestion raisonnée}. Je cite 4 ressources camerounaises (eau, sol, forêt, biodiversité) et leurs pressions spécifiques.
    \item \textbf{Savoir-faire :} J'applique la \textbf{Méthode 4 étapes} pour analyser n'importe quel dossier (forêt, lac, parc, sol). Je construis un \textbf{tableau bilan Ressource/Pression/Conséquence/Mesure}. Je propose des mesures \textbf{locales, concrètes, argumentées 3 piliers}.
    \item \textbf{Savoir-être :} Je comprends l'urgence (lac Tchad, éléphants Bénoué) mais je refuse le fatalisme. J'adopte une posture de \textbf{citoyen responsable} : je questionne ma propre consommation (bois, eau, viande de brousse), je soutiens les initiatives locales (forêts communautaires, agroforesterie), j'exige la transparence (RFA, contrats).
\end{itemize}
\end{aretenir}

\begin{methode}[Dernier conseil pour l'épreuve écrite]
Face au dossier du Probatoire : \textbf{1.} Lisez la question AVANT les documents. \textbf{2.} Annotez chaque document (chiffres, acteurs, flèches causales). \textbf{3.} Rédigez en phrases courtes, structurées (Pression $\rightarrow$ Conséquence $\rightarrow$ Mesure). \textbf{4.} Utilisez le vocabulaire précis : « taux de prélèvement », « taux de reconstitution », « fragmentation », « corridor écologique », « GIRE », « parties prenantes », « trois piliers ». \textbf{5.} Concluez toujours par une phrase d'ouverture : « Cette gestion durable nécessite un suivi-évaluation participatif et un financement pérenne (ex. fonds environnement, paiement pour services écosystémiques). »
\end{methode}

\textbf{Félicitations !} Vous avez maintenant tous les outils pour maîtriser le Module 1 « Le Monde Vivant » et briller au Probatoire. La gestion durable des ressources naturelles n'est pas un chapitre de plus : c'est \textbf{la compétence citoyenne majeure} pour l'avenir du Cameroun et de la planète. À vous de jouer !

\newpage

\coursec{Activité d'intégration}

\coursec{Leçon ND07 : Réduction des conséquences néfastes des activités humaines sur les ressources naturelles}

\courssub{Objectifs de l'activité}

Cette activité d'intégration te place dans la peau d'un conseiller municipal chargé de proposer un plan de gestion durable des ressources naturelles d'un village camerounais. À l'issue de cette activité, tu devras être capable de :

\begin{itemize}
\item \cle{identifier} les impacts d'activités humaines données (extension urbaine, agriculture sur brûlis, exploitation d'une carrière) sur les ressources naturelles (eau, sol, forêt, biodiversité) à partir d'un dossier documentaire ;
\item \cle{exploiter} des données chiffrées : calculer un taux de déboisement annuel, le comparer à un rythme de régénération naturelle et en tirer une conclusion sur la durabilité de l'exploitation ;
\item \cle{proposer} des mesures de gestion durable adaptées à une situation locale réelle, en justifiant chaque mesure par les savoirs de la leçon ;
\item \cle{communiquer} une argumentation scientifique devant un public (le « conseil municipal »), en distinguant faits, interprétations et propositions.
\end{itemize}

\courssub{Situation}

\textbf{Conseil municipal de gestion durable : le cas du village de Nkolbisson}

Nkolbisson est un village situé à la périphérie ouest de Yaoundé, dans la région du Centre. Traditionnellement, les habitants y vivaient de l'agriculture vivrière (manioc, macabo, plantain, arachide) et de la cueillette des produits forestiers. Le village est traversé par une source pérenne, la source \emph{Mefou}, qui alimente environ \num{1200} habitants en eau de boisson et d'usage domestique.

Depuis une dizaine d'années, la croissance rapide de la ville de Yaoundé a profondément transformé le village :

\begin{itemize}
\item l'\textbf{extension des habitations} : de nouveaux quartiers se construisent sur les anciennes terres agricoles et forestières, sans plan d'assainissement ; les eaux usées et les ordures sont rejetées à l'air libre, en amont de la source ;
\item l'\textbf{agriculture sur brûlis} : pour nourrir une population croissante, les paysans défrichent chaque année de nouvelles parcelles de forêt par le feu, cultivent deux ou trois ans, puis abandonnent les terres épuisées pour en défricher d'autres ;
\item l'\textbf{exploitation d'une carrière de sable} : une entreprise prélève du sable dans le lit du ruisseau Mefou, en aval de la source, sans autorisation ni remise en état du site.
\end{itemize}

Le chef du village, inquiet, a fait réaliser un suivi sur dix ans (2014--2024). Voici les données recueillies.

\begin{center}
\rowcolors{1}{popblueL}{white}
\begin{tabularx}{\linewidth}{|l|X|X|X|}
\hline
\rowcolor{popblue!40} \textbf{Année} & \textbf{Débit de la source Mefou (\SI{}{L/min})} & \textbf{Surface forestière du village (\SI{}{ha})} & \textbf{Rendement moyen du manioc (\SI{}{t/ha})} \\
\hline
2014 & 45 & 320 & 12 \\
2016 & 40 & 296 & 11 \\
2018 & 34 & 268 & 9 \\
2020 & 27 & 240 & 7 \\
2022 & 21 & 216 & 6 \\
2024 & 16 & 192 & 5 \\
\hline
\end{tabularx}
\end{center}

Autres informations du dossier :

\begin{itemize}
\item Les analyses d'eau réalisées en 2024 par le centre de santé montrent une \textbf{contamination bactériologique} de la source (présence de coliformes fécaux) et une \textbf{turbidité élevée} (eau trouble), surtout en saison des pluies. Les cas de diarrhées et de typhoïde ont augmenté.
\item Les pentes déboisées présentent des \textbf{ravines d'érosion} ; la couche de terre végétale est emportée vers le ruisseau.
\item La forêt communautaire abritait autrefois des essences précieuses (iroko, moabi), des petits mammifères et de nombreux oiseaux ; ces espèces sont devenues rares.
\item On estime qu'une forêt secondaire laissée en jachère longue \textbf{régénère environ \SI{2}{ha/an}} sur les terres abandonnées du village, si elle n'est pas perturbée.
\end{itemize}

\textbf{Témoignage de Mme Abena, cultivatrice} : « Avant, la source coulait fort toute l'année et l'eau était claire. Maintenant, en saison sèche, on fait la queue dès quatre heures du matin. Mes champs de manioc ne donnent plus rien : la terre est devenue dure et rouge, l'eau l'emporte quand il pleut. Les jeunes coupent le bois pour vendre des planches aux gens de Yaoundé. Si ça continue, nos enfants n'auront ni eau, ni terre, ni forêt. »

Tu es membre d'un groupe d'experts invité par le conseil municipal. Ta mission : analyser la situation et proposer un \textbf{plan communal de gestion durable} des ressources de Nkolbisson.

\courssub{Consignes}

\textbf{Travail 1 — Identification des impacts.} À partir du dossier, complète un tableau à trois colonnes : \emph{activité humaine} / \emph{ressource affectée} / \emph{impact observé (avec la donnée du dossier qui le prouve)}. Chaque ligne du tableau de données doit être exploitée au moins une fois.

\textbf{Travail 2 — Analyse quantitative de la durabilité.}
\begin{enumerate}
\item Calcule la surface forestière perdue entre 2014 et 2024, puis le \textbf{taux de déboisement annuel moyen} (en \SI{}{ha/an}).
\item Compare ce taux au rythme de régénération naturelle (\SI{2}{ha/an}). La forêt de Nkolbisson est-elle exploitée de manière durable ? Justifie.
\item Si rien ne change, en combien d'années la forêt communale aura-t-elle totalement disparu à partir de 2024 ?
\item Calcule la baisse du débit de la source entre 2014 et 2024 en pourcentage. Quel lien peux-tu établir avec le déboisement ?
\end{enumerate}

\textbf{Travail 3 — Rédaction du plan communal de gestion durable.} Rédige un plan comportant \textbf{au moins une mesure concrète par ressource} (eau, sol, forêt, biodiversité). Pour chaque mesure, indique : la ressource visée, l'impact qu'elle corrige, et l'acteur responsable (commune, villageois, entreprise de la carrière, services de l'État). Intègre au moins une mesure de \textbf{protection de la source}, une mesure d'\textbf{agroforesterie}, une mesure de \textbf{reboisement} et une mesure de \textbf{réglementation de la carrière}.

\textbf{Travail 4 — Restitution en plénière.} Chaque groupe présente son plan en cinq minutes, comme devant le conseil municipal. Les autres groupes jouent le rôle des conseillers : ils posent des questions et vérifient que chaque mesure proposée est justifiée par une donnée du dossier.

\courssub{Ressources à mobiliser}

\begin{itemize}
\item \textbf{Savoirs} : typologie des ressources naturelles (eau, sols, forêts, biodiversité) ; pressions des activités humaines (déforestation, pollution, érosion, surexploitation) ; notion de \cle{développement durable} (« répondre aux besoins du présent sans compromettre la capacité des générations futures à répondre aux leurs ») et ses trois piliers (environnemental, économique, social) ; mesures de conservation (reboisement, agroforesterie, zones tampons, réglementation des prélèvements).
\item \textbf{Savoir-faire} : identifier un impact à partir d'un document ; analyser un dossier documentaire ; proposer des mesures adaptées à une situation locale.
\item \textbf{Outils de calcul} :
\[ \text{Taux annuel moyen} = \frac{\text{variation totale}}{\text{nombre d'années}} \qquad ; \qquad \text{Baisse en \%} = \frac{\text{valeur initiale} - \text{valeur finale}}{\text{valeur initiale}} \times 100 \]
\item \textbf{Critère de durabilité d'un prélèvement} : une ressource renouvelable est exploitée durablement si le \cle{rythme de prélèvement} est inférieur ou égal au \cle{rythme de régénération}.
\end{itemize}

\courssub{Démarche et corrigé commenté}

\begin{exoresolu}[Corrigé commenté de l'activité]
\textbf{Énoncé.} Identifier les impacts des activités humaines à Nkolbisson, évaluer la durabilité de l'exploitation forestière et proposer un plan communal de gestion durable.

\tcblower

\textbf{Travail 1 — Tableau des impacts.}

\begin{center}
\rowcolors{1}{popgreenL}{white}
\begin{tabularx}{\linewidth}{|X|l|X|}
\hline
\rowcolor{popgreen!40} \textbf{Activité humaine} & \textbf{Ressource affectée} & \textbf{Impact (preuve dans le dossier)} \\
\hline
Extension des habitations sans assainissement & Eau & Contamination bactériologique de la source (coliformes fécaux) ; hausse des diarrhées et de la typhoïde \\
\hline
Agriculture sur brûlis & Forêt & Recul de la forêt : de \SI{320}{ha} (2014) à \SI{192}{ha} (2024) \\
\hline
Agriculture sur brûlis + déboisement des pentes & Sol & Érosion en ravines, terre végétale emportée ; rendement du manioc divisé par plus de 2 (\SI{12}{t/ha} à \SI{5}{t/ha}) \\
\hline
Carrière de sable dans le lit du ruisseau & Eau et sol & Turbidité élevée de l'eau, surtout en saison des pluies ; déstabilisation des berges \\
\hline
Déboisement global (coupe de bois, brûlis) & Biodiversité & Raréfaction de l'iroko, du moabi, des petits mammifères et des oiseaux \\
\hline
Déboisement global & Eau & Chute du débit de la source : de \SI{45}{L/min} à \SI{16}{L/min} \\
\hline
\end{tabularx}
\end{center}

\emph{Commentaire} : chaque impact est relié à une donnée précise du dossier — c'est la différence entre une affirmation et une preuve scientifique.

\textbf{Travail 2 — Analyse quantitative.}
\begin{enumerate}
\item Surface perdue : $320 - 192 = \SI{128}{ha}$ en 10 ans, soit un taux de déboisement annuel moyen de :
\[ \frac{128}{10} = \num{12,8} \ \text{ha/an} \]
\item Comparaison : \num{12,8} \SI{}{ha/an} de perte contre \num{2} \SI{}{ha/an} de régénération. Le prélèvement est \textbf{6,4 fois supérieur} à la régénération ($12,8 / 2 = 6,4$). L'exploitation n'est donc \textbf{pas durable} : le capital forestier diminue chaque année.
\item Durée avant disparition totale à partir de 2024 : en tenant compte de la régénération, la perte nette est de $12,8 - 2 = \num{10,8} \ \text{ha/an}$, d'où :
\[ \frac{192}{10,8} \approx 17,8 \ \text{ans} \]
Sans changement, la forêt communale aura disparu vers \textbf{2042}.
\item Baisse du débit de la source :
\[ \frac{45 - 16}{45} \times 100 \approx 64\ \% \]
La source a perdu environ \textbf{64 \%} de son débit en dix ans. Le lien avec le déboisement est double : la forêt retient l'eau de pluie et favorise son infiltration dans le sol (recharge de la nappe qui alimente la source) ; une fois la forêt coupée, l'eau ruisselle, érode le sol et s'écoule directement vers le ruisseau au lieu d'alimenter la nappe.
\end{enumerate}

\textbf{Travail 3 — Exemple de plan communal de gestion durable.}

\begin{center}
\rowcolors{1}{popblueL}{white}
\begin{tabularx}{\linewidth}{|l|X|X|l|}
\hline
\rowcolor{popblue!40} \textbf{Ressource} & \textbf{Mesure proposée} & \textbf{Impact corrigé} & \textbf{Acteur} \\
\hline
Eau & Créer une zone de protection de 50 m autour de la source (interdiction de construire et de cultiver) ; construire des latrines et un tri des ordures en amont & Contamination bactériologique, maladies hydriques & Commune + villageois \\
\hline
Sol & Introduire l'agroforesterie (association manioc--arbres fruitiers et fixateurs d'azote comme le calliandra) ; cultiver en courbes de niveau sur les pentes ; allonger les jachères & Érosion, baisse des rendements & Villageois + encadrement agricole \\
\hline
Forêt & Reboiser \SI{15}{ha/an} avec des essences locales (iroko, moabi) ; créer un comité villageois de gestion des coupes de bois & Déboisement (\num{12,8} \SI{}{ha/an}) supérieur à la régénération & Commune + services des Eaux et Forêts \\
\hline
Eau et sol & Réglementer la carrière : autorisation préalable, prélèvement limité hors lit du ruisseau, remise en état obligatoire du site & Turbidité, déstabilisation des berges & Entreprise + État (contrôle) \\
\hline
Biodiversité & Protéger un reliquat forestier comme réserve communautaire ; interdire la chasse en période de reproduction & Raréfaction des espèces & Village + ONG \\
\hline
\end{tabularx}
\end{center}

\emph{Commentaire} : le plan respecte les trois piliers du développement durable — \textbf{environnemental} (reboisement, protection de la source), \textbf{économique} (l'agroforesterie maintient les revenus paysans, la carrière continue mais encadrée), \textbf{social} (eau potable, santé, implication des habitants).

\textbf{Travail 4 — Restitution.} En plénière, chaque groupe doit pouvoir répondre à la question du conseiller : « quelle donnée du dossier justifie cette mesure ? ». Par exemple, le reboisement de \SI{15}{ha/an} se justifie ainsi : il compense la perte nette de \num{10,8} \SI{}{ha/an} et permet de reconstituer progressivement le couvert forestier.
\end{exoresolu}

\courssub{Bilan de l'activité}

Cette activité t'a permis de mettre en évidence plusieurs idées essentielles de la leçon :

\begin{itemize}
\item Les ressources naturelles (eau, sol, forêt, biodiversité) sont \textbf{interdépendantes} : à Nkolbisson, le déboisement a entraîné l'érosion des sols, la baisse des rendements, la pollution et le tarissement progressif de la source. Dégrader une ressource, c'est souvent en dégrader plusieurs.
\item Une ressource renouvelable n'est renouvelable que si le \textbf{rythme de prélèvement reste inférieur au rythme de régénération} : le calcul du taux de déboisement (\num{12,8} \SI{}{ha/an} contre \num{2} \SI{}{ha/an} de régénération) transforme une impression (« la forêt recule ») en un diagnostic chiffré (« la forêt aura disparu vers 2042 »).
\item Le \textbf{développement durable} n'est pas l'arrêt des activités humaines, mais leur \textbf{gestion raisonnée} : la carrière peut continuer si elle est réglementée, l'agriculture peut nourrir le village si l'agroforesterie remplace le brûlis itinérant.
\item Les solutions efficaces sont \textbf{locales et concertées} : elles associent les villageois, la commune, l'entreprise et l'État, et s'appuient sur des données mesurées — exactement la démarche que tu devras reproduire au Probatoire face à un dossier documentaire sur l'exploitation d'une ressource au Cameroun.
\end{itemize}

\begin{attention}[Pièges fréquents]
Ne confonds pas \emph{constat} et \emph{mesure} : « la source est polluée » est un constat, « créer une zone de protection de 50 m » est une mesure. Dans les calculs, n'oublie pas de soustraire la régénération pour obtenir la perte nette, et exprime toujours tes résultats avec une unité (\SI{}{ha/an}, \%, années). Enfin, chaque mesure proposée doit être reliée à une donnée du dossier : une proposition sans justification chiffrée ne vaut pas de point en analyse de dossier.
\end{attention}

\newpage

\coursec{Méthodes et exercices résolus}

\coursec{ND07 : Réduction des conséquences néfastes des activités humaines sur les ressources naturelles}

\courssub{Rappels de cours : Ressources naturelles et développement durable}

\begin{definition}[Ressource naturelle]
Une \cle{ressource naturelle} est un élément de l'environnement (milieu physique ou vivant) que les sociétés humaines utilisent pour satisfaire leurs besoins (alimentation, énergie, habitat, santé, culture). On distingue classiquement :
\begin{itemize}
\item les \cle{ressources renouvelables} : elles se régénèrent naturellement à une échelle de temps humaine \textbf{si leur rythme d'exploitation ne dépasse pas leur rythme de renouvellement} (eau douce, sols, forêts, biodiversité, air, ressources halieutiques) ;
\item les \cle{ressources non renouvelables} (ou exhaustibles) : leur formation géologique prend des millions d'années, tout prélèvement est définitif à l'échelle humaine (minerais, pétrole, gaz naturel, charbon).
\end{itemize}
\textbf{Attention :} une ressource renouvelable devient \emph{épuisable} si elle est surexploitée (ex. : nappe phréatique, stock forestier, population de poissons).
\end{definition}

\begin{aretenir}[Principales pressions anthropiques sur les ressources au Cameroun]
\begin{center}
\begin{tabularx}{\linewidth}{|l|X|}\hline
\textbf{Ressource} & \textbf{Pressions principales (activités humaines)} \\ \hline
\cle{Eau} & Pollution (industrielle, agricole, domestique), prélèvements excessifs (irrigation, barrage), envasement des cours d'eau, destruction des zones humides. \\ \hline
\cle{Sols} & Érosion hydrique et éolienne (déforestation, surpâturage, cultures sur pente), appauvrissement en matière organique (raccourcissement des jachères), salinisation (mauvaise irrigation), pollution (pesticides, déchets). \\ \hline
\cle{Forêts} & Exploitation forestière illégale ou non durable, déforestation pour l'agriculture (palmier à huile, cacao, cultures vivrières), feux de brousse, prélèvement de bois-énergie, urbanisation. \\ \hline
\cle{Biodiversité} & Braconnage (viande de brousse, ivoire, écailles de pangolin), destruction/fragmentation des habitats, espèces invasives, pollution, changement climatique. \\ \hline
\end{tabularx}
\end{center}
\end{aretenir}

\begin{definition}[Développement durable]
Le \cle{développement durable} est « un développement qui répond aux besoins des générations présentes sans compromettre la capacité des générations futures à répondre aux leurs » (\emph{Rapport Brundtland}, 1987). Il repose sur l'articulation de \textbf{trois piliers indissociables} :
\begin{popfigure}
\begin{tikzpicture}[scale=0.9, every node/.style={font=\small}]
  % Couleurs
  \colorlet{env}{popgreen}
  \colorlet{eco}{popblue}
  \colorlet{soc}{poporange}
  % Cercles
  \node[circle, draw=env, thick, fill=env!20, minimum width=3.5cm, align=center] (E) at (0,1.5){\textbf{Environnemental}\\\emph{Préservation des ressources, biodiversité, climat}};
  \node[circle, draw=eco, thick, fill=eco!20, minimum width=3.5cm, align=center] (C) at (-1.7,-1){\textbf{Économique}\\\emph{Efficacité, rentabilité long terme, emplois verts}};
  \node[circle, draw=soc, thick, fill=soc!20, minimum width=3.5cm, align=center] (S) at (1.7,-1){\textbf{Social}\\\emph{Équité, participation, éducation, santé, culture}};
  % Intersections
  \node[align=center, font=\footnotesize\bfseries] at (0,0) {\textcolor{popdark}{Développement Durable}};
  \node[align=center, font=\footnotesize] at (-0.9,0.3) {Viable};
  \node[align=center, font=\footnotesize] at (0.9,0.3) {Équitable};
  \node[align=center, font=\footnotesize] at (0,-1.3) {Supportable};
\end{tikzpicture}
\legende{Les trois piliers du développement durable et leurs intersections (Viable = Environnement + Économie ; Équitable = Économie + Social ; Supportable = Social + Environnement).}
\end{popfigure}
\end{definition}

\begin{propriete}[Principes de la gestion raisonnée des ressources]
Pour qu'une exploitation soit \cle{durable}, elle doit respecter trois règles opérationnelles (inspirées de H. Daly, 1990) :
\begin{enumerate}
\item \textbf{Règle des ressources renouvelables :} le taux de prélèvement $\le$ taux de régénération naturel.
\item \textbf{Règle des ressources non renouvelables :} le taux d'épuisement $\le$ taux de création de substituts renouvelables (investissement dans les alternatives).
\item \textbf{Règle des déchets/polluants :} le taux d'émission $\le$ capacité d'assimilation naturelle des écosystèmes.
\end{enumerate}
\end{propriete}

\courssub{Méthode 1 : Identifier les impacts d'une activité humaine sur une ressource naturelle}

\begin{methode}[Analyser l'impact d'une activité humaine à partir d'un document]
Pour identifier rigoureusement les impacts d'une activité humaine (déforestation, agriculture intensive, urbanisation, exploitation minière...) sur une ressource naturelle, suis toujours la même démarche en 5 étapes :
\begin{enumerate}
\item \textbf{Identifier la ressource concernée} : eau, sol, forêt, biodiversité (faune et flore), air, ressources minérales. Précise s'il s'agit d'une ressource \cle{renouvelable} (eau, forêt, sol — à condition de respecter leur rythme de renouvellement) ou \cle{non renouvelable} (minerais, pétrole).
\item \textbf{Décrire l'activité humaine} en la citant précisément à partir du document (ex. : « abattage de 500 ha de forêt pour l'extension de plantations de palmiers à huile »).
\item \textbf{Relever les données chiffrées} du document : surfaces, tonnages, pourcentages, dates. Un impact non chiffré est un impact mal argumenté.
\item \textbf{Établir la chaîne de causalité} : activité $\rightarrow$ mécanisme $\rightarrow$ conséquence sur la ressource.
\item \textbf{Classer l'impact} : direct/indirect, à court terme/à long terme, réversible/irréversible (l'extinction d'une espèce est toujours irréversible).
\end{enumerate}
\textbf{Réflexe rédaction} : rédige en phrases complètes, en citant le document (« le document 1 montre que... »), et conclus par un bilan : la ressource est-elle menacée dans sa quantité, sa qualité, ou les deux ?
\end{methode}

\begin{popfigure}
\begin{tikzpicture}[>=Stealth, node distance=1.2cm and 1.5cm, every node/.style={font=\small}]
  \node[draw, rounded corners, fill=popblue!10, thick] (act) {Activité humaine\normalsize (ex. exploitation forestière)};
  \node[draw, rounded corners, fill=poporange!10, thick, right=of act] (mec) {Mécanisme physique/biologique\normalsize (ex. suppression du couvert, pistes d'accès)};
  \node[draw, rounded corners, fill=popred!10, thick, right=of mec] (cons) {Conséquence sur la ressource\normalsize (ex. érosion, braconnage, turbidité)};
  \draw[->, thick] (act) -- (mec);
  \draw[->, thick] (mec) -- (cons);
\end{tikzpicture}
\legende{Schéma générique d'une chaîne de causalité : Activité $\rightarrow$ Mécanisme $\rightarrow$ Conséquence.}
\end{popfigure}

\begin{exoresolu}[Impact de l'exploitation forestière dans le Sud-Est Cameroun]
\textbf{Énoncé.} Document : « Dans la région de l'Est du Cameroun, une société forestière exploite une concession de \SI{80000}{ha}. Entre 2010 et 2020, le couvert forestier de la zone est passé de 92 \% à 68 \% de la surface totale. Les pistes d'exploitation ont favorisé l'installation de chasseurs ; les captures de gorilles et de chimpanzés destinés à la viande de brousse ont été multipliées par trois. Les cours d'eau traversant la zone présentent une turbidité (trouble de l'eau) en augmentation constante. »
\begin{enumerate}
\item Identifie les ressources naturelles affectées par cette activité.
\item Pour chaque ressource, établis la chaîne de causalité reliant l'activité à l'impact observé.
\item Précise le caractère réversible ou irréversible de chaque impact.
\end{enumerate}
\tcblower
\textbf{Solution détaillée.}

\textbf{1. Ressources affectées.} Trois ressources sont touchées :
\begin{itemize}
\item la \cle{forêt} (ressource renouvelable à long terme) : le couvert forestier passe de 92 \% à 68 \%, soit une perte de 24 points en 10 ans ;
\item la \cle{biodiversité} (faune : gorilles, chimpanzés) ;
\item l'\cle{eau} des cours d'eau (dégradation de la qualité : turbidité en augmentation).
\end{itemize}

\textbf{2. Chaînes de causalité.}
\begin{itemize}
\item \emph{Forêt} : exploitation forestière (abattage des arbres) $\rightarrow$ diminution du couvert forestier de 92 \% à 68 \% $\rightarrow$ fragmentation de l'habitat et réduction du stock de bois.
\item \emph{Biodiversité} : ouverture des pistes d'exploitation $\rightarrow$ accès facilité aux zones autrefois inaccessibles $\rightarrow$ installation de chasseurs $\rightarrow$ braconnage $\rightarrow$ captures de grands singes multipliées par trois $\rightarrow$ déclin des populations de gorilles et chimpanzés.
\item \emph{Eau} : abattage des arbres et passage d'engins $\rightarrow$ sol nu, non retenu par les racines $\rightarrow$ érosion par les pluies tropicales $\rightarrow$ transport de particules de sol vers les rivières $\rightarrow$ augmentation de la turbidité (dégradation de la qualité de l'eau).
\end{itemize}

\textbf{3. Réversibilité.}
\begin{itemize}
\item La perte de couvert forestier est \textbf{partiellement réversible} à très long terme (reforestation, régénération naturelle : plusieurs décennies à siècles pour une forêt primaire).
\item Le déclin des grands singes peut devenir \textbf{irréversible} : si la population passe sous un seuil critique, l'espèce s'éteint localement — une extinction est définitive.
\item La turbidité de l'eau est \textbf{réversible} à moyen terme si l'érosion cesse (restauration du couvert végétal).
\end{itemize}

\textbf{Conclusion.} L'exploitation forestière non maîtrisée dégrade simultanément trois ressources — forêt (quantité), faune (biodiversité) et eau (qualité) — ce qui illustre l'interdépendance des ressources naturelles et justifie une gestion durable des concessions forestières.
\end{exoresolu}

\courssub{Méthode 2 : Proposer des mesures de gestion durable adaptées à une situation locale}

\begin{methode}[Construire une proposition de gestion durable]
Une bonne proposition de mesure de gestion durable doit être \textbf{adaptée, argumentée et réaliste}. Procède en 4 étapes :
\begin{enumerate}
\item \textbf{Diagnostiquer} : nomme précisément le problème (quelle ressource ? quelle pression ? quelle échelle : village, bassin versant, région ?).
\item \textbf{Proposer des mesures} en couvrant idéalement les trois piliers du \cle{développement durable} :
\begin{itemize}
\item \emph{environnemental} : reboisement, création d'aires protégées (parcs nationaux, réserves), rotation des cultures, jachère, agroforesterie, traitement des eaux usées ;
\item \emph{économique} : écotourisme, certification forestière (FSC), agriculture raisonnée rentable, valorisation des produits forestiers non ligneux (miel, fruits, plantes médicinales) ;
\item \emph{social} : éducation environnementale, implication des populations locales (forêts communautaires), création d'emplois verts, respect des droits des communautés autochtones (Baka, Bagyéli).
\end{itemize}
\item \textbf{Justifier chaque mesure} : explique le mécanisme par lequel elle réduit la pression (ex. : « la jachère permet la reconstitution de la matière organique du sol, donc le maintien de sa fertilité »).
\item \textbf{Vérifier la faisabilité locale} : coût, acceptabilité par les populations, existence d'un cadre légal (loi forestière camerounaise de 1994, forêts communautaires, plans d'aménagement).
\end{enumerate}
\textbf{À retenir} : une mesure de gestion durable vise à exploiter la ressource à un rythme \textbf{inférieur ou égal à son rythme de renouvellement}, tout en assurant les besoins des générations présentes sans compromettre ceux des générations futures (définition du développement durable, rapport Brundtland, 1987).
\end{methode}

\begin{exoresolu}[Dégradation des sols agricoles dans l'Ouest Cameroun]
\textbf{Énoncé.} Dans les Hautes Terres de l'Ouest Cameroun, la densité de population élevée a conduit au raccourcissement des jachères (de 10 ans à moins de 2 ans), à la culture sur des pentes raides et à l'usage croissant d'engrais chimiques. Les rendements du maïs sont passés de 2,5 à 1,2 tonne par hectare en 15 ans, et les ravines d'érosion se multiplient.
\begin{enumerate}
\item Explique l'origine de la baisse des rendements.
\item Propose quatre mesures de gestion durable des sols adaptées à cette région, en justifiant chacune.
\end{enumerate}
\tcblower
\textbf{Solution détaillée.}

\textbf{1. Origine de la baisse des rendements.} La jachère est une période de repos du sol pendant laquelle la végétation naturelle restitue de la matière organique et des éléments minéraux (azote, phosphore, potassium). Son raccourcissement (de 10 ans à moins de 2 ans) ne laisse plus au sol le temps de reconstituer sa fertilité : les exportations minérales par les récoltes successives épuisent le stock nutritif. De plus, la culture sur pentes raides sans protection expose le sol nu aux pluies : l'érosion hydrique emporte l'horizon supérieur, le plus riche en humus (d'où les ravines). Les rendements chutent donc de 2,5 à 1,2 t/ha, soit une baisse de :
\[ \frac{2{,}5 - 1{,}2}{2{,}5} \times 100 = 52 \,\% \]
Les engrais chimiques compensent partiellement mais n'apportent pas de matière organique et peuvent acidifier le sol à long terme.

\textbf{2. Mesures de gestion durable.}
\begin{itemize}
\item \emph{Agroforesterie} (association d'arbres et de cultures, ex. caféiers sous ombrage, \emph{Acacia} ou \emph{Leucaena} en haies) : les racines des arbres fixent le sol, leur litière apporte de la matière organique — lutte contre l'érosion et restaure la fertilité.
\item \emph{Culture en courbes de niveau et construction de bandes végétalisées (haies) sur les pentes} : elles ralentissent le ruissellement, donc réduisent l'érosion hydrique et la formation de ravines.
\item \emph{Allongement des jachères ou rotation culturale incluant des légumineuses} (haricot, arachide, soja, mucuna) : les légumineuses fixent l'azote atmosphérique grâce aux bactéries symbiotiques de leurs nodosités racinaires (\emph{Rhizobium} spp.), ce qui enrichit naturellement le sol en azote et réduit le recours aux engrais chimiques.
\item \emph{Compostage et fumure organique} (fumier, déchets végétaux, résidus de cuisine) : apport de matière organique qui améliore la structure du sol, sa rétention en eau et sa fertilité durable, à faible coût pour les paysans.
\end{itemize}

\textbf{Conclusion.} Ces mesures, accessibles aux paysans des Hautes Terres, concilient les trois piliers du développement durable : elles restaurent le sol (environnement), maintiennent les rendements (économie) et ne demandent pas d'investissement inaccessible (social).
\end{exoresolu}

\courssub{Méthode 3 : Analyser un dossier documentaire sur l'exploitation d'une ressource au Cameroun}

\begin{methode}[Analyser un dossier documentaire (méthode complète)]
L'analyse de dossier est l'exercice le plus fréquent au Probatoire. Adopte la méthode « LIRE — TRIER — CROISER — CONCLURE » :
\begin{enumerate}
\item \textbf{LIRE} : lis d'abord les questions, puis les documents (texte, tableau, carte, graphique) en notant pour chacun : nature du document, source, date, information principale.
\item \textbf{TRIER} : pour chaque question, repère le(s) document(s) utile(s) et extrais uniquement les informations pertinentes (chiffres, faits).
\item \textbf{CROISER} : confronte les documents entre eux (un tableau de chiffres + un texte de loi + une carte permettent souvent de répondre à une question de synthèse). Cite toujours tes sources : « d'après le document 2... ».
\item \textbf{CONCLURE} : rédige une synthèse structurée : rappel du problème $\rightarrow$ faits établis $\rightarrow$ enjeux (environnementaux, économiques, sociaux) $\rightarrow$ solutions possibles.
\end{enumerate}
\textbf{Réflexes} : ne recopie jamais un document intégralement ; distingue toujours le \emph{fait} (donnée du document) de l'\emph{interprétation} (ce que tu en déduis) ; sois attentif aux unités et aux échelles des graphiques.
\end{methode}

\begin{exoresolu}[Exploitation pétrolière et ressources halieutiques dans le Golfe de Guinée]
\textbf{Énoncé.} Document 1 : « Le Cameroun exploite du pétrole offshore au large de Kribi et Limbé depuis 1977. La production nationale a culminé à environ 180 000 barils/jour en 1985, puis a décliné pour atteindre environ 70 000 barils/jour en 2020. » Document 2 : « En 2019, une fuite sur un oléoduc a provoqué une nappe d'hydrocarbures de plusieurs kilomètres carrés. Les pêcheurs artisanaux de la côte ont signalé une chute de 40 \% de leurs captures dans les mois suivants. » Document 3 : « Le pétrole représente encore une part importante des recettes d'exportation du Cameroun et finance des infrastructures (routes, hôpitaux). »
\begin{enumerate}
\item À quel type de ressource le pétrole appartient-il ? Justifie.
\item Dégage le conflit d'usage révélé par les documents.
\item Rédige une synthèse de 8 à 10 lignes présentant les enjeux et les solutions possibles.
\end{enumerate}
\tcblower
\textbf{Solution détaillée.}

\textbf{1. Type de ressource.} Le pétrole est une ressource naturelle \cle{non renouvelable} : sa formation (décomposition de matière organique enfouie pendant des millions d'années) est infiniment plus lente que sa consommation. Le déclin de la production camerounaise (de 180 000 à 70 000 barils/jour entre 1985 et 2020, soit une baisse de $\dfrac{180000-70000}{180000}\times 100 \approx 61\,\%$) illustre l'épuisement progressif des gisements.

\textbf{2. Conflit d'usage.} Les documents 2 et 3 révèlent un conflit entre :
\begin{itemize}
\item l'exploitation pétrolière, source de devises et de financements publics (doc. 3) ;
\item la pêche artisanale, ressource renouvelable dont dépendent les communautés côtières, mise en danger par les pollutions accidentelles (chute de 40 \% des captures après la fuite, doc. 2).
\end{itemize}
La pollution par hydrocarbures tue poissons, crustacés et larves, et contamine les zones de reproduction (mangroves, estuaires) : c'est la ressource \emph{biodiversité marine} et la ressource \emph{halieutique} (halieutique = relatif à la pêche) qui paient le coût de l'exploitation pétrolière.

\textbf{3. Synthèse (exemple de rédaction).} Le pétrole, ressource non renouvelable, procure au Cameroun des recettes d'exportation essentielles (doc. 3), mais son exploitation offshore menace les ressources renouvelables de la mer : la fuite de 2019 a entraîné une chute de 40 \% des captures artisanales (doc. 2), frappant des populations déjà vulnérables. De plus, l'épuisement des gisements (déclin de 61 \% de la production entre 1985 et 2020, doc. 1) impose de préparer l'après-pétrole. Une gestion durable suppose : le renforcement des normes de sécurité et des contrôles des oléoducs, la création de fonds d'indemnisation des pêcheurs, la protection des zones de reproduction (mangroves), et surtout l'investissement des revenus pétroliers dans les énergies renouvelables (hydroélectricité, solaire) et la diversification économique, afin de concilier développement actuel et préservation des ressources pour les générations futures.

\textbf{Conclusion.} Ce dossier illustre le principe central du développement durable : une ressource non renouvelable ne doit pas être exploitée au détriment des ressources renouvelables ni de l'avenir des populations.
\end{exoresolu}

\courssub{Méthode 4 : Exploiter des données chiffrées (taux de déforestation, évolutions en pourcentage)}

\begin{methode}[Calculer et interpréter un taux d'évolution]
Les sujets de Probatoire demandent souvent de calculer puis d'interpréter un taux d'évolution d'une ressource. Retiens :
\begin{enumerate}
\item \textbf{Taux d'évolution} (en \%) :
\[ t = \frac{V_{\text{finale}} - V_{\text{initiale}}}{V_{\text{initiale}}} \times 100 \]
Un résultat négatif traduit une diminution.
\item \textbf{Taux annuel moyen} (si l'évolution s'étale sur $n$ années) : $t_{\text{moyen}} = \dfrac{t}{n}$ (approximation linéaire acceptée au lycée).
\item \textbf{Interprétation obligatoire} : un calcul non commenté ne rapporte pas tous les points. Traduis toujours le résultat en une phrase : « la forêt a perdu X \% de sa surface en n ans, soit Y ha par an ».
\item \textbf{Vérification} : contrôle l'ordre de grandeur (un taux de perte ne peut pas dépasser 100 \% ; une surface restante ne peut pas être négative).
\end{enumerate}
\end{methode}

\begin{exoresolu}[Déforestation du bassin du Congo]
\textbf{Énoncé.} La couverture forestière d'une province camerounaise était de \SI{1200000}{ha} en 2000. En 2020, elle n'était plus que de \SI{1020000}{ha}.
\begin{enumerate}
\item Calcule la surface de forêt perdue et le taux de déforestation sur la période.
\item Calcule le rythme annuel moyen de perte en hectares.
\item Si ce rythme se maintient, en quelle année la forêt aura-t-elle perdu la moitié de sa surface de l'an 2000 ? Commente.
\end{enumerate}
\tcblower
\textbf{Solution détaillée.}

\textbf{1. Surface perdue et taux de déforestation.}
\[ \text{Surface perdue} = 1\,200\,000 - 1\,020\,000 = \SI{180000}{ha} \]
\[ t = \frac{1\,020\,000 - 1\,200\,000}{1\,200\,000} \times 100 = \frac{-180\,000}{1\,200\,000} \times 100 = -15\,\% \]
La forêt a perdu \SI{180000}{ha}, soit \textbf{15 \% de sa surface en 20 ans}.

\textbf{2. Rythme annuel moyen.}
\[ \frac{180\,000}{20} = \SI{9000}{ha/an} \]
Chaque année, en moyenne, \SI{9000}{ha} de forêt disparaissent (l'équivalent d'environ 12 600 terrains de football de 0,714 ha chacun).

\textbf{3. Prévision.} La moitié de la surface de 2000 vaut :
\[ \frac{1\,200\,000}{2} = \SI{600000}{ha} \]
Il faut donc perdre \SI{600000}{ha} au total. À \SI{9000}{ha/an}, la durée nécessaire est :
\[ \frac{600\,000}{9\,000} \approx 66{,}7 \text{ ans} \]
soit vers l'année $2000 + 67 = 2067$.

\textbf{Commentaire.} Ce modèle linéaire est une simplification : en réalité, la déforestation s'accélère souvent avec la croissance démographique et l'extension des fronts agricoles, donc l'échéance pourrait être plus proche. Inversement, des politiques de reboisement et d'aménagement forestier pourraient la repousser. Ce calcul montre néanmoins qu'un rythme apparemment « faible » (15 \% en 20 ans) conduit, s'il se poursuit, à la disparition de la moitié de la forêt en moins d'un siècle : d'où l'urgence d'une gestion durable.

\textbf{Conclusion.} La province a perdu 15 \% de sa forêt entre 2000 et 2020, au rythme moyen de \SI{9000}{ha/an} ; au même rythme, la moitié de la forêt de 2000 aurait disparu vers 2067.
\end{exoresolu}

\courssub{Méthode 5 : Définir et mobiliser la notion de développement durable dans une argumentation}

\begin{methode}[Utiliser correctement la notion de développement durable dans une copie]
\begin{enumerate}
\item \textbf{Cite la définition de référence} : le développement durable est « un développement qui répond aux besoins des générations présentes sans compromettre la capacité des générations futures à répondre aux leurs » (rapport Brundtland, 1987).
\item \textbf{Mobilise les trois piliers} : environnemental (préserver les ressources et la biodiversité), économique (assurer une activité rentable à long terme), social (équité, implication des populations, éducation). Une solution qui néglige un pilier est incomplète.
\item \textbf{Applique les principes de gestion raisonnée} : ne pas prélever une ressource renouvelable plus vite qu'elle ne se régénère ; ne pas émettre de déchets plus vite que la nature ne les recycle ; substituer des ressources renouvelables aux ressources non renouvelables.
\item \textbf{Illustre par un exemple camerounais} : parcs nationaux (Waza, Bénoué, Korup, Campo-Ma'an), forêts communautaires et plans d'aménagement (loi de 1994), reboisement (opération « Sahel vert », Grande Muraille Verte dans l'Extrême-Nord), traitement des eaux usées, interdiction des sacs plastiques non biodégradables.
\item \textbf{Conclus toujours} en reliant la mesure proposée à la préservation de la ressource pour les générations futures.
\end{enumerate}
\end{methode}

\begin{popfigure}
\begin{tikzpicture}[scale=0.85, every node/.style={font=\small}]
  % Trois piliers comme colonnes
  \node[draw, rounded corners, fill=popgreen!15, thick, minimum width=3cm, minimum height=2.5cm, align=center] (E) at (0,0){\textbf{Pilier Environnemental}\\\textit{Préserver le capital naturel}\\\textbullet\ Biodiversité\\ \textbullet\ Qualité air/eau/sol\\ \textbullet\ Climat\\ \textbullet\ Non-prélèvement $>$ Renouvellement};
  \node[draw, rounded corners, fill=popblue!15, thick, minimum width=3cm, minimum height=2.5cm, align=center] (C) at (4,0){\textbf{Pilier Économique}\\\textit{Efficacité \& Viabilité}\\\textbullet\ Rentabilité long terme\\ \textbullet\ Emplois verts\\ \textbullet\ Certification (FSC)\\ \textbullet\ Substitution ressources};
  \node[draw, rounded corners, fill=poporange!15, thick, minimum width=3cm, minimum height=2.5cm, align=center] (S) at (8,0){\textbf{Pilier Social}\\\textit{Équité \& Participation}\\\textbullet\ Droits communautés\\ \textbullet\ Éducation env.\textbullet\ Santé\\ \textbullet\ Partage des bénéfices};
  % Flèches convergence
  \draw[->, thick, popdark] (E.east) -- node[above, font=\tiny] {Viable} (C.west);
  \draw[->, thick, popdark] (C.east) -- node[above, font=\tiny] {Équitable} (S.west);
  \draw[->, thick, popdark] (S.north west) .. controls +(0,1) and +(0,1) .. node[above, font=\tiny] {Supportable} (E.north east);
  % Centre
  \node[circle, draw=popdark, thick, fill=popgold!30, minimum width=2.5cm, align=center] at (4,-2){\textbf{Développement}\\\textbf{Durable}\\\textit{Brundtland 1987}};
  \draw[->, dashed, popdark] (E.south) -- (4,-1.5);
  \draw[->, dashed, popdark] (C.south) -- (4,-1.5);
  \draw[->, dashed, popdark] (S.south) -- (4,-1.5);
\end{tikzpicture}
\legende{Articulation des trois piliers du développement durable : chaque pilier a des objectifs propres, leur intersection définit la viabilité, l'équité et la supportabilité, et leur convergence centrale est le développement durable.}
\end{popfigure}

\begin{exoresolu}[Le lac Tchad et la gestion partagée d'une ressource en eau]
\textbf{Énoncé.} Le lac Tchad, partagé entre le Cameroun, le Tchad, le Niger et le Nigéria, est passé d'environ \SI{25000}{km^2} en 1960 à moins de \SI{2500}{km^2} certaines années récentes, sous l'effet combiné de la sécheresse, des prélèvements pour l'irrigation et de la croissance démographique. Environ 30 millions de personnes vivent de ses ressources (pêche, agriculture, élevage).
\begin{enumerate}
\item Calcule le taux de rétraction de la surface du lac.
\item Explique pourquoi la gestion de cette ressource relève du développement durable, en mobilisant les trois piliers.
\item Propose deux mesures de gestion concertée entre les pays riverains.
\end{enumerate}
\tcblower
\textbf{Solution détaillée.}

\textbf{1. Taux de rétraction.}
\[ t = \frac{2\,500 - 25\,000}{25\,000} \times 100 = \frac{-22\,500}{25\,000} \times 100 = -90\,\% \]
Le lac a perdu environ \textbf{90 \% de sa surface} en un peu plus d'un demi-siècle.

\textbf{2. Développement durable et trois piliers.}
\begin{itemize}
\item \emph{Pilier environnemental} : l'eau du lac est une ressource renouvelable mais \emph{limitée} ; des prélèvements supérieurs au rythme de recharge (pluies, apports du fleuve Chari) conduisent à son assèchement, à la disparition des zones humides et de leur biodiversité (poissons, oiseaux migrateurs).
\item \emph{Pilier économique} : 30 millions de personnes tirent leurs revenus de la pêche, de l'agriculture et de l'élevage autour du lac ; sa disparition détruirait ces activités et provoquerait paupérisation et migrations.
\item \emph{Pilier social} : le partage de la ressource entre quatre pays et entre usages concurrents (irrigation, pêche, abreuvement du bétail) exige équité, coopération et prévention des conflits.
\end{itemize}
La gestion du lac illustre donc parfaitement la définition de Brundtland : satisfaire les besoins actuels de 30 millions de personnes sans hypothéquer la ressource pour les générations futures.

\textbf{3. Mesures de gestion concertée.}
\begin{itemize}
\item Renforcer la \emph{Commission du Bassin du Lac Tchad} (CBLT, créée en 1964) : quotas de prélèvement d'eau négociés entre les États, suivi scientifique commun du niveau du lac, techniques d'irrigation économes (goutte-à-goutte) plutôt que l'irrigation par submersion.
\item Restaurer les écosystèmes du bassin : reboisement des berges et des bassins versants (réduction de l'érosion et de l'envasement), protection des zones de reproduction des poissons, et sensibilisation des populations riveraines à une pêche et une agriculture durables.
\end{itemize}

\textbf{Conclusion.} La rétraction de 90 \% du lac Tchad est un cas d'école : seule une gestion concertée, équitable et scientifiquement fondée — c'est-à-dire durable — peut sauver une ressource vitale partagée par quatre nations.
\end{exoresolu}

\courssub{Méthode 6 : Analyser un exemple camerounais de conservation (aire protégée)}

\begin{methode}[Évaluer une mesure de conservation]
Pour analyser un exemple de conservation (parc national, réserve, reboisement), structure ta réponse ainsi :
\begin{enumerate}
\item \textbf{Présente la mesure} : nature (aire protégée, plan d'aménagement, campagne de reboisement), localisation, date de création, superficie.
\item \textbf{Identifie la ressource protégée} : biodiversité (espèces emblématiques), forêt, sols, bassin versant.
\item \textbf{Évalue l'efficacité} à partir des données du document : évolution des populations animales, taux de braconnage, surface reboisée. Distingue les succès des limites (braconnage persistant, conflits homme-faune, moyens insuffisants des éco-gardes).
\item \textbf{Montre la dimension « développement durable »} : retombées pour les populations riveraines (écotourisme, emplois d'éco-gardes, forêts communautaires reversant des redevances aux villages).
\end{enumerate}
\textbf{Exemples camerounais à connaître} : parc national de Waza (Nord, créé en 1934 comme réserve, parc en 1968 ; lions, éléphants, girafes de Kordofan), parc de la Bénoué (réserve de biosphère UNESCO), parc de Korup (Sud-Ouest, forêt primaire parmi les plus anciennes d'Afrique), réserve de faune du Dja (patrimoine mondial UNESCO, 1987), sanctuaires de grands singes (Mefou, Limbe Wildlife Centre).
\end{methode}

\begin{exoresolu}[Le parc national de Waza face aux pressions humaines]
\textbf{Énoncé.} Document : « Le parc national de Waza (Extrême-Nord, environ \SI{170000}{ha}) abritait dans les années 1970 plusieurs centaines de lions et de vastes troupeaux d'éléphants. Le braconnage, l'empiètement des troupeaux domestiques, le prélèvement de bois et l'insécurité ont fait chuter ces populations : en 2010, on comptait moins de 30 lions. Un programme de réhabilitation (renforcement des éco-gardes, aménagement de points d'eau, coopération avec les villages riverains) a été lancé. »
\begin{enumerate}
\item Identifie les pressions exercées sur la biodiversité du parc.
\item Calcule le taux de déclin de la population de lions entre les années 1970 (on retiendra 300 lions) et 2010.
\item Explique en quoi la coopération avec les villages riverains est une condition de succès de la conservation.
\end{enumerate}
\tcblower
\textbf{Solution détaillée.}

\textbf{1. Pressions identifiées.} Le document mentionne quatre pressions d'origine humaine :
\begin{itemize}
\item le \emph{braconnage} (prélèvement direct de la faune) ;
\item l'\emph{empiètement des troupeaux domestiques} (compétition pour les pâturages et l'eau, transmission de maladies, dégradation de l'habitat) ;
\item le \emph{prélèvement de bois} (dégradation du couvert végétal, donc de l'habitat des espèces) ;
\item l'\emph{insécurité} (qui empêche la surveillance et l'écotourisme, source de revenus).
\end{itemize}

\textbf{2. Taux de déclin des lions.}
\[ t = \frac{30 - 300}{300} \times 100 = \frac{-270}{300} \times 100 = -90\,\% \]
La population de lions a chuté d'environ \textbf{90 \%} en quatre décennies, passant sous un seuil où la viabilité de la population devient critique (consanguinité, difficulté de reproduction).

\textbf{3. Rôle des villages riverains.} Une aire protégée entourée de populations pauvres et hostiles ne peut pas survivre : si la conservation n'apporte rien aux riverains, ceux-ci continueront à prélever bois et gibier pour vivre. La coopération transforme les villageois en acteurs de la protection : emplois d'éco-gardes et de guides, revenus de l'écotourisme reversés aux communautés, forêts communautaires périphériques servant de zone tampon, éducation environnementale. C'est l'application concrète du pilier social du développement durable : la conservation de la biodiversité (pilier environnemental) n'est durable que si elle améliore aussi la vie des populations (piliers social et économique).

\textbf{Conclusion.} Le cas de Waza montre qu'une aire protégée ne se gère pas seulement par l'interdiction, mais par une gestion intégrée associant surveillance, aménagement et développement local — condition pour que la faune emblématique du Nord Cameroun soit transmise aux générations futures.
\end{exoresolu}

\begin{attention}[Les 5 erreurs qui coûtent des points au Probatoire]
\begin{enumerate}
\item \textbf{Confondre ressource renouvelable et inépuisable.} L'eau, le sol ou la forêt sont renouvelables \emph{à condition} que le rythme de prélèvement respecte le rythme de renouvellement : une ressource renouvelable surexploitée peut disparaître (lac Tchad, sols érodés). Ne jamais écrire « l'eau est inépuisable ».
\item \textbf{Donner un calcul sans unité et sans interprétation.} Un taux de déforestation de « 15 » n'a aucun sens : il faut écrire « 15 \% » et ajouter une phrase d'interprétation (« la forêt a perdu 15 \% de sa surface en 20 ans »). Le calcul seul ne vaut que la moitié des points.
\item \textbf{Citer le développement durable sans le définir ni l'appliquer.} Écrire « il faut un développement durable » sans la définition de Brundtland ni les trois piliers (environnemental, économique, social) est une réponse creuse. Exige-toi toujours : définition + piliers + exemple concret.
\item \textbf{Proposer des mesures génériques non adaptées au contexte.} « Interdire la pêche » ou « planter des arbres » sans lien avec le document est hors sujet : chaque mesure doit répondre au problème précis du dossier (pression identifiée, échelle locale, faisabilité pour les populations).
\item \textbf{Recopier les documents au lieu de les exploiter.} Au Probatoire, il faut \emph{extraire} l'information pertinente, la \emph{mettre en relation} avec tes connaissances (chaîne de causalité : activité $\rightarrow$ mécanisme $\rightarrow$ impact) et \emph{citer la source} (« d'après le document 2... »). La paraphrase intégrale d'un document est sanctionnée.
\end{enumerate}
\end{attention}

\newpage

\coursec{Exercices}

\courssub{Je vérifie mes connaissances}

\begin{exercice}[QCM : ressources et développement durable]
Pour chaque question, entoure la (ou les) bonne(s) réponse(s).

\textbf{1.} Une ressource naturelle renouvelable est une ressource :
\begin{itemize}
\item[a)] qui ne s'épuise jamais, quelle que soit la manière dont on l'utilise ;
\item[b)] qui se régénère à l'échelle d'une vie humaine si le prélèvement ne dépasse pas sa capacité de renouvellement ;
\item[c)] qui se forme sur des millions d'années ;
\item[d)] qui est toujours gratuite.
\end{itemize}

\textbf{2.} Parmi les ressources suivantes, laquelle est \textbf{non renouvelable} ?
\begin{itemize}
\item[a)] l'eau des nappes phréatiques profondes fossiles ;
\item[b)] le pétrole du bassin du Rio del Rey ;
\item[c)] le bois d'une forêt exploitée avec reboisement ;
\item[d)] le minerai de fer de Mbalam.
\end{itemize}

\textbf{3.} Le développement durable est un développement qui :
\begin{itemize}
\item[a)] interdit toute exploitation des ressources naturelles ;
\item[b)] répond aux besoins du présent sans compromettre la capacité des générations futures à répondre aux leurs ;
\item[c)] concerne uniquement la protection de l'environnement ;
\item[d)] repose sur trois piliers : économique, social et environnemental.
\end{itemize}

\textbf{4.} Le reboisement est une mesure de conservation adaptée à la ressource :
\begin{itemize}
\item[a)] eau ; \quad b) sol ; \quad c) forêt ; \quad d) biodiversité.
\end{itemize}
\end{exercice}

\begin{exercice}[Vrai ou faux justifié]
Indique si chaque affirmation est vraie ou fausse, puis justifie ta réponse en une ou deux phrases.
\begin{enumerate}
\item « La déforestation ne concerne que la disparition des arbres ; elle n'a aucune conséquence sur le climat. »
\item « Une ressource renouvelable peut disparaître si elle est surexploitée. »
\item « Les aires protégées du Cameroun (parcs nationaux, réserves) servent uniquement au tourisme. »
\item « Les feux de brousse répétés dans l'Extrême-Nord favorisent la désertification. »
\end{enumerate}
\end{exercice}

\begin{exercice}[Définitions]
Définis en deux ou trois phrases chacun des termes suivants, en donnant à chaque fois un exemple camerounais :
\begin{enumerate}
\item \cle{ressource naturelle} ;
\item \cle{gestion raisonnée} (ou \cle{gestion durable}) d'une ressource ;
\item \cle{aire protégée} ;
\item \cle{désertification}.
\end{enumerate}
\end{exercice}

\begin{exercice}[Calcul direct : taux de déforestation]
La superficie forestière d'une région du Cameroun est passée de \num{500000} hectares en 1990 à \num{410000} hectares en 2020.
\begin{enumerate}
\item Calcule la superficie de forêt perdue sur cette période.
\item Calcule le pourcentage de forêt perdu par rapport à la superficie de 1990.
\item Calcule le taux de perte annuel moyen en hectares par an.
\item À ce rythme constant, en combien d'années la forêt de cette région aurait-elle totalement disparu à partir de 2020 ? Conclus sur la durabilité de cette exploitation.
\end{enumerate}
\end{exercice}

\courssub{Je m'entraîne}

\begin{exercice}[Le lac Tchad : un écosystème en danger]
\textbf{Document.} Le lac Tchad, situé à l'extrême nord du Cameroun, couvrait environ \num{25000} km$^2$ en 1960. Il n'en couvre plus qu'environ \num{2500} km$^2$ aujourd'hui. Ce recul s'explique par la baisse des précipitations dans la région sahélienne, mais aussi par le pompage intensif de l'eau pour l'irrigation des cultures et par la construction de barrages sur les rivières qui alimentent le lac. Des milliers de pêcheurs et d'agriculteurs camerounais, tchadiens, nigérians et nigériens dépendent directement de ce lac.
\begin{enumerate}
\item À partir du document, distingue les causes \textbf{naturelles} et les causes \textbf{humaines} du recul du lac Tchad.
\item Cite deux conséquences de cet assèchement sur les populations riveraines.
\item Propose trois mesures concrètes de gestion durable de la ressource en eau dans ce bassin, en précisant pour chacune le pilier du développement durable concerné (économique, social ou environnemental).
\end{enumerate}
\end{exercice}

\begin{exercice}[Associer mesures et ressources]
Recopie et complète le tableau suivant en associant à chaque ressource naturelle une pression humaine majeure et une mesure de conservation adaptée, choisies dans les listes proposées.

\textbf{Pressions :} déforestation ; pompage excessif ; érosion et épuisement des sols ; braconnage ; pollution industrielle.

\textbf{Mesures :} reboisement ; création d'aires protégées ; techniques d'irrigation économes ; rotation des cultures et agroforesterie ; traitement des rejets avant déversement.

\begin{center}
\begin{tabularx}{\linewidth}{|l|X|X|}\hline
\rowcolor{popblueL} \textbf{Ressource} & \textbf{Pression humaine majeure} & \textbf{Mesure de conservation} \\ \hline
Eau & & \\ \hline
Sols & & \\ \hline
Forêts & & \\ \hline
Biodiversité & & \\ \hline
\end{tabularx}
\end{center}
\end{exercice}

\begin{exercice}[Carte des aires protégées du Cameroun]
Sur une carte mure du Cameroun (ou un fond de carte reproduit), place les quatre aires protégées suivantes : parc national de Waza (Extrême-Nord), parc national de la Bénoué (Nord), parc national de Korup (Sud-Ouest), réserve de faune du Dja (Sud-Est).
\begin{enumerate}
\item Pour chaque aire protégée, indique la région et le type de milieu dominé (savane sèche, savane arborée, forêt dense humide).
\item Explique en quelques lignes le rôle de ces aires protégées dans la conservation de la biodiversité camerounaise.
\item Cite deux menaces qui pèsent encore sur ces espaces malgré leur statut de protection.
\end{enumerate}
\end{exercice}

\begin{exercice}[Exploitation forestière au Sud-Est]
Une société forestière obtient une concession de \num{60000} hectares dans la région de l'Est. Elle prévoit d'abattre \num{3000} hectares par an, sans programme de reboisement, et emploie 400 personnes de la région.
\begin{enumerate}
\item Calcule en combien d'années la concession serait totalement coupée à ce rythme.
\item Identifie deux impacts économiques positifs et trois impacts écologiques négatifs de cette exploitation.
\item Propose deux conditions qui rendraient cette exploitation durable (pense notamment au taux de prélèvement et au renouvellement de la ressource).
\end{enumerate}
\end{exercice}

\courssub{Je me prépare au Probatoire}

\begin{exercice}[Étude de cas : la forêt camerounaise en chiffres]
\textbf{Document 1.} Évolution de la superficie forestière du Cameroun :

\begin{center}
\begin{tabularx}{\linewidth}{|l|X|X|X|X|}\hline
\rowcolor{popblueL} \textbf{Année} & 1990 & 2000 & 2010 & 2020 \\ \hline
Superficie forestière (millions d'ha) & 24,5 & 23,2 & 21,9 & 20,6 \\ \hline
\end{tabularx}
\end{center}

\textbf{Document 2.} La forêt camerounaise abrite plus de \num{8000} espèces végétales et constitue le deuxième massif forestier d'Afrique. Elle fournit du bois d'œuvre, des produits forestiers non ligneux (nourriture, pharmacopée), régule le régime des pluies et protège les sols de l'érosion.

\begin{enumerate}
\item Calcule la perte totale de superficie forestière entre 1990 et 2020, puis le taux de déforestation annuel moyen en millions d'hectares par an.
\item Représente graphiquement l'évolution de la superficie forestière en fonction du temps (axe des abscisses : années ; axe des ordonnées : superficie en millions d'hectares).
\item À partir des documents, cite trois fonctions de la forêt camerounaise qui sont menacées par cette déforestation.
\item L'exploitation forestière rapporte des devises et des emplois au Cameroun. Rédige un texte argumenté de dix à quinze lignes montrant comment concilier exploitation économique de la forêt et conservation de la ressource, en t'appuyant sur la notion de développement durable.
\end{enumerate}
\end{exercice}

\begin{exercice}[Situation-problème : un village de l'Extrême-Nord]
\textbf{Document.} Dans un village de la région de l'Extrême-Nord, les habitants pratiquent chaque année des feux de brousse pour défricher de nouveaux champs et coupent massivement les arbres pour le bois de chauffe, seule source d'énergie domestique. En vingt ans, la couverture végétale a fortement régressé : les sols, nus pendant la saison sèche, sont emportés par les vents et les premières pluies ; les rendements agricoles chutent ; les points d'eau s'assèchent de plus en plus tôt dans l'année. Certains jeunes quittent le village.

\begin{enumerate}
\item Construis, sous forme de chaîne de causes à effets (schéma fléché), les étapes qui mènent des pratiques du village à la désertification. Tu feras apparaître au minimum : feux de brousse et coupe de bois, disparition du couvert végétal, érosion et appauvrissement du sol, baisse des rendements, désertification.
\item Explique pourquoi la disparition du couvert végétal accélère l'érosion des sols.
\item Propose un plan d'action local en quatre mesures concrètes et réalistes (une mesure par domaine : énergie, agriculture, reboisement, sensibilisation), en justifiant chaque mesure par son effet attendu sur la ressource concernée.
\item Montre en quelques lignes que ce plan d'action correspond aux trois piliers du développement durable.
\end{enumerate}
\end{exercice}

\begin{exercice}[Analyse de dossier : l'eau potable à Yaoundé]
\textbf{Document 1.} La ville de Yaoundé compte plus de 4 millions d'habitants. La demande en eau potable croît d'environ 5 \% par an, tandis que les sources et rivières périurbaines (Mfoundi, Mingoa) reçoivent des rejets domestiques et industriels non traités. Le traitement de l'eau polluée coûte de plus en plus cher à la société de distribution.

\textbf{Document 2.} On estime qu'un habitant de Yaoundé utilise en moyenne 80 litres d'eau par jour, mais qu'environ 30 \% de l'eau produite est perdue par des fuites du réseau de distribution.

\begin{enumerate}
\item À partir du document 1, identifie les deux pressions majeures exercées par les activités humaines sur la ressource en eau à Yaoundé.
\item Explique le lien entre la pollution des rivières périurbaines et le coût de production de l'eau potable.
\item Un quartier de \num{50000} habitants consomme en moyenne 80 litres d'eau par habitant et par jour. Calcule la consommation journalière totale du quartier, puis le volume d'eau perdu chaque jour par les fuites du réseau (30 \%).
\item Propose trois mesures de gestion durable de l'eau adaptées à la situation de Yaoundé, en distinguant celles qui relèvent des pouvoirs publics, des entreprises et des habitants.
\end{enumerate}
\end{exercice}

\newpage

\coursec{Corrigés des exercices}

\begin{corrige}[Exercice 1 — QCM : ressources et développement durable]
\textbf{1. Réponse b).} Une ressource renouvelable (forêt, eau de surface, poissons) se régénère à l'échelle d'une vie humaine, \emph{à condition} que le rythme de prélèvement ne dépasse pas sa capacité de renouvellement. L'affirmation a) est fausse : une ressource renouvelable surexploitée peut s'épuiser (exemple : effondrement des stocks de poissons du lac Tchad). L'affirmation c) décrit une ressource \emph{non renouvelable} (pétrole, minerais). L'affirmation d) est fausse : le caractère renouvelable est une propriété \emph{biologique ou physico-chimique} de la ressource, sans aucun rapport avec son prix de marché.

\textbf{2. Réponses b) et d).} Le pétrole du bassin du Rio del Rey (région du Littoral) et le minerai de fer de Mbalam (région de l'Est) se sont formés sur des dizaines à des centaines de millions d'années : leur stock est fixe à l'échelle humaine, ce sont des ressources \textbf{non renouvelables}. L'eau des nappes souterraines se recharge par infiltration des eaux de pluie (lentement pour les nappes profondes) : elle relève des ressources hydriques, renouvelables si les pompages restent inférieurs à la recharge. Le bois d'une forêt exploitée avec reboisement est renouvelable.

\textbf{3. Réponses b) et d).} La définition de référence est celle du rapport Brundtland (1987, Commission mondiale sur l'environnement et le développement de l'ONU) : un développement qui répond aux besoins du présent sans compromettre la capacité des générations futures à répondre aux leurs. Il repose sur trois piliers indissociables : \cle{économique}, \cle{social} et \cle{environnemental}. L'affirmation a) est fausse : le développement durable encadre et régule l'exploitation, il ne l'interdit pas. L'affirmation c) est fausse : l'environnement n'est qu'un des trois piliers ; un projet qui protège la nature mais ruine les populations locales n'est pas durable.

\textbf{4. Réponses b), c) et d).} Le reboisement restaure directement la ressource \textbf{forêt} ; en reconstituant le couvert végétal, il protège aussi les \textbf{sols} de l'érosion (les racines fixent les particules, le feuillage amortit les gouttes de pluie) et abrite la \textbf{biodiversité} (reconstitution des habitats). Il n'est pas, à lui seul, une mesure directe de conservation de la ressource \textbf{eau}, même s'il améliore l'infiltration et la recharge des nappes.

\textbf{Points de vigilance :} au Probatoire, un QCM peut comporter plusieurs bonnes réponses ; lis toutes les propositions avant de répondre et justifie chaque choix si la consigne le demande.
\end{corrige}

\begin{corrige}[Exercice 2 — Vrai ou faux justifié]
\begin{enumerate}
\item \textbf{Faux.} La déforestation supprime des arbres qui absorbaient le \ce{CO2} atmosphérique par photosynthèse ; de plus, la combustion ou la décomposition de la biomasse abattue libère ce carbone sous forme de \ce{CO2}. Elle contribue donc à l'augmentation de l'effet de serre et aux changements climatiques, en plus de réduire localement l'évapotranspiration, donc les précipitations.
\item \textbf{Vrai.} Une ressource renouvelable ne se régénère que si le prélèvement reste inférieur à sa capacité de renouvellement : une forêt coupée plus vite qu'elle ne repousse, ou un stock de poissons prélevé au-delà de sa reproduction, peut disparaître durablement. Le caractère renouvelable est donc \emph{conditionnel}.
\item \textbf{Faux.} Les aires protégées (parc national de Waza, réserve de faune du Dja\ldots) ont d'abord pour rôle de \cle{conserver la biodiversité} : protéger les espèces menacées, préserver les milieux naturels, servir de zones de reproduction et de réservoirs génétiques. L'écotourisme n'est qu'une retombée économique secondaire, qui peut d'ailleurs financer la protection.
\item \textbf{Vrai.} Les feux de brousse répétés détruisent le couvert végétal ; les sols mis à nu sont alors érodés par le vent de la saison sèche et par les pluies, s'appauvrissent en humus et en éléments nutritifs, ce qui, dans un climat sahélien déjà sec, accélère l'avancée de la désertification.
\end{enumerate}
\textbf{Points de vigilance :} une réponse « vrai » ou « faux » sans justification ne rapporte aucun point ; la justification doit mobiliser un mécanisme précis (photosynthèse, érosion, capacité de renouvellement) et non une simple reformulation de l'affirmation.
\end{corrige}

\begin{corrige}[Exercice 3 — Définitions]
\begin{enumerate}
\item \textbf{Ressource naturelle :} élément du milieu naturel (eau, sol, forêt, minerais, espèces vivantes) que l'être humain prélève et utilise pour satisfaire ses besoins. On distingue les ressources \emph{renouvelables} (qui se régénèrent à l'échelle d'une vie humaine) et \emph{non renouvelables} (stock fixe). \emph{Exemple :} le bois de la forêt de l'Est du Cameroun, l'eau du fleuve Sanaga.
\item \textbf{Gestion raisonnée (durable) :} mode d'exploitation d'une ressource qui organise les prélèvements de façon qu'ils ne dépassent pas la capacité de renouvellement de la ressource, afin qu'elle reste disponible pour les générations futures. \emph{Exemple :} les plans d'aménagement des concessions forestières du Sud-Est, qui imposent des quotas de coupe, un diamètre minimal d'exploitabilité et des reboisements.
\item \textbf{Aire protégée :} espace terrestre ou aquatique délimité par la loi, où les activités humaines sont réglementées afin de conserver les milieux et les espèces. \emph{Exemple :} le parc national de Waza (Extrême-Nord), la réserve de faune du Dja (classée au patrimoine mondial de l'UNESCO depuis 1987).
\item \textbf{Désertification :} dégradation durable des terres en zone sèche, sous l'effet conjugué du climat (sécheresses récurrentes) et des activités humaines (déboisement, feux de brousse, surpâturage), aboutissant à des sols stériles et improductifs. \emph{Exemple :} l'avancée de la désertification dans la région de l'Extrême-Nord, autour du bassin du lac Tchad.
\end{enumerate}
\textbf{Points de vigilance :} une définition attendue au Probatoire comporte l'idée essentielle (nature de l'objet + propriété caractéristique) et, ici, un exemple camerounais explicitement demandé par la consigne.
\end{corrige}

\begin{corrige}[Exercice 4 — Calcul direct : taux de déforestation]
\begin{enumerate}
\item Superficie perdue entre 1990 et 2020 :
\[ 500000 - 410000 = \num{90000}\ \text{hectares}. \]
\item Pourcentage perdu par rapport à la superficie initiale de 1990 :
\[ \frac{90000}{500000}\times 100 = 18\ \%. \]
\item Durée écoulée : $2020 - 1990 = 30$ ans. Taux de perte annuel moyen :
\[ \frac{90000}{30} = \num{3000}\ \text{ha/an}. \]
\item En 2020, il reste \num{410000} ha. Au rythme constant de \num{3000} ha/an :
\[ \frac{410000}{3000} \approx 136,7\ \text{ans}. \]
La forêt aurait totalement disparu en environ \textbf{137 ans}, soit vers l'an 2157. \textbf{Conclusion :} ce rythme d'exploitation n'est pas durable, car la ressource est prélevée sans être renouvelée ; elle finira par s'épuiser, privant les générations futures de la ressource forestière et de ses services (bois, régulation du climat, protection des sols).
\end{enumerate}
\textbf{Points de vigilance :} bien distinguer le pourcentage perdu (rapport à la valeur \emph{initiale} de 1990) du taux annuel (rapport à la \emph{durée}) ; ne pas oublier l'unité (ha/an) et rédiger une conclusion argumentée sur la durabilité, comme l'exige la dernière question.
\end{corrige}

\begin{corrige}[Exercice 5 — Le lac Tchad : un écosystème en danger]
\begin{enumerate}
\item \textbf{Cause naturelle :} la baisse durable des précipitations dans la région sahélienne (sécheresses récurrentes depuis les années 1970). \textbf{Causes humaines :} le pompage intensif de l'eau pour l'irrigation et la construction de barrages sur les rivières tributaires (Chari, Logone), qui réduisent les apports au lac.
\item Conséquences sur les populations riveraines (deux parmi) : disparition de la pêche et perte de revenus pour les pêcheurs ; réduction des terres irriguées et insécurité alimentaire ; conflits entre usagers (agriculteurs, éleveurs, pêcheurs) pour l'eau et les terres restantes ; migrations forcées des populations.
\item Trois mesures de gestion durable :
\begin{itemize}
\item \textbf{Irrigation économe en eau} (goutte-à-goutte, canaux étanches) : réduit les prélèvements tout en maintenant les rendements agricoles (piliers \emph{environnemental} et \emph{économique}).
\item \textbf{Gestion concertée transfrontalière des prélèvements} via la Commission du bassin du lac Tchad (CBLT), avec quotas de pompage négociés entre les États riverains (pilier \emph{social} : garantir un accès équitable à l'eau pour tous les usagers).
\item \textbf{Reboisement des berges et restauration des zones humides} : limiter l'évaporation, protéger les rives de l'érosion et préserver les habitats de reproduction des poissons (pilier \emph{environnemental}).
\end{itemize}
\end{enumerate}
\textbf{Points de vigilance :} bien \emph{distinguer} causes naturelles et causes humaines à partir du document (ne rien inventer) ; chaque mesure proposée doit être explicitement rattachée à un pilier du développement durable, comme le demande la consigne.
\end{corrige}

\begin{corrige}[Exercice 6 — Associer mesures et ressources]
\begin{center}
\begin{tabularx}{\linewidth}{|l|X|X|}\hline
\rowcolor{popblueL} \textbf{Ressource} & \textbf{Pression humaine majeure} & \textbf{Mesure de conservation} \\ \hline
Eau & Pompage excessif & Techniques d'irrigation économes \\ \hline
Sols & Érosion et épuisement des sols & Rotation des cultures et agroforesterie \\ \hline
Forêts & Déforestation & Reboisement \\ \hline
Biodiversité & Braconnage & Création d'aires protégées \\ \hline
\end{tabularx}
\end{center}
\textbf{Remarque :} la pression « pollution industrielle » et la mesure « traitement des rejets avant déversement » forment un couple cohérent qui concerne également l'eau (et les sols) ; elles étaient proposées en surplus pour éviter les associations automatiques. L'essentiel est que chaque mesure réponde logiquement à la pression : on lutte contre un prélèvement excessif par l'économie d'usage, contre la destruction d'un milieu par sa restauration ou sa protection, contre un prélèvement illégal d'espèces par la protection réglementaire.
\textbf{Points de vigilance :} vérifier la cohérence pression $\rightarrow$ mesure : une bonne mesure s'attaque à la \emph{cause} de la dégradation, pas seulement à ses symptômes.
\end{corrige}

\begin{corrige}[Exercice 7 — Carte des aires protégées du Cameroun]
\begin{enumerate}
\item Localisation et milieux :
\begin{center}
\begin{tabularx}{\linewidth}{|l|X|X|}\hline
\rowcolor{popblueL} \textbf{Aire protégée} & \textbf{Région} & \textbf{Milieu dominant} \\ \hline
Parc national de Waza & Extrême-Nord & Savane sèche (sahélienne) \\ \hline
Parc national de la Bénoué & Nord & Savane arborée (soudanienne) \\ \hline
Parc national de Korup & Sud-Ouest & Forêt dense humide \\ \hline
Réserve de faune du Dja & Est et Sud (Sud-Est) & Forêt dense humide \\ \hline
\end{tabularx}
\end{center}
Sur la carte, Waza se place au nord de Maroua, la Bénoué au sud-est de Garoua, Korup près de la frontière nigériane au nord-ouest de Mundemba, et le Dja entre Sangmélima et Lomie.
\item \textbf{Rôle :} ces aires protégées couvrent les grands types de milieux du Cameroun (savanes sèches et arborées, forêts denses humides). Elles abritent des espèces emblématiques et menacées (éléphants, lions, girafes à Waza et à la Bénoué ; gorilles, chimpanzés et forêts primaires à Korup et au Dja), servent de zones de reproduction et de réservoirs de biodiversité, et préservent les fonctions écologiques des milieux (protection des sols, régulation hydrique et climatique).
\item \textbf{Menaces persistantes (deux) :} le braconnage (ivoire, viande de brousse) ; l'orpaillage et l'exploitation minière ou forestière en périphérie ; l'avancée des fronts agricoles et les feux de brousse ; le sous-financement de la surveillance (écogardes insuffisants).
\end{enumerate}
\textbf{Points de vigilance :} ne pas confondre les régions (la Bénoué est dans la région du Nord, pas dans l'Adamaoua) ; associer correctement le milieu au climat de la région (savane sèche au climat sahélien, forêt dense humide au climat équatorial).
\end{corrige}

\begin{corrige}[Exercice 8 — Exploitation forestière au Sud-Est]
\begin{enumerate}
\item Durée d'épuisement de la concession :
\[ \frac{\num{60000}\ \text{ha}}{\num{3000}\ \text{ha/an}} = 20\ \text{ans}. \]
La concession serait totalement coupée en \textbf{20 ans}.
\item \textbf{Impacts économiques positifs (deux) :} création de 400 emplois locaux ; recettes d'exportation (devises) et taxes pour l'État et les communes ; ouverture de pistes désenclavant les villages. \textbf{Impacts écologiques négatifs (trois) :} destruction des habitats et baisse de la biodiversité ; érosion et lessivage des sols mis à nu ; libération de \ce{CO2} et réduction du puits de carbone forestier ; perturbation du régime des pluies et de l'écoulement des cours d'eau.
\item \textbf{Conditions de durabilité :}
\begin{itemize}
\item \textbf{Limiter le taux de prélèvement} à la capacité de renouvellement de la forêt : coupe sélective, rotation longue de 25 à 30 ans par parcelle, respect d'un diamètre minimal d'exploitabilité pour laisser les jeunes arbres assurer la régénération ;
\item \textbf{Reboiser systématiquement} les parcelles exploitées et mettre en place un plan d'aménagement certifié (par exemple le label FSC), avec suivi écologique indépendant.
\end{itemize}
\end{enumerate}
\textbf{Points de vigilance :} la question 2 exige de \emph{séparer} impacts économiques et écologiques ; la question 3 attend explicitement les notions de taux de prélèvement et de renouvellement de la ressource — c'est le cœur de la gestion durable.
\end{corrige}

\begin{corrige}[Exercice 9 — Étude de cas : la forêt camerounaise en chiffres]
\begin{enumerate}
\item Perte totale 1990--2020 : $24{,}5 - 20{,}6 = 3{,}9$ millions d'hectares. Durée : $2020-1990 = 30$ ans. Taux annuel moyen :
\[ \frac{3{,}9}{30} = 0{,}13\ \text{million d'ha/an}, \]
soit \num{130000} hectares perdus par an en moyenne.
\item Représentation graphique :

\begin{popfigure}
\begin{center}
\begin{tikzpicture}
\begin{axis}[
width=0.85\linewidth, height=6cm,
xlabel={Année}, ylabel={Superficie forestière (millions d'ha)},
xmin=1988, xmax=2022, ymin=20, ymax=25,
xtick={1990,2000,2010,2020}, ytick={20,21,22,23,24,25},
grid=major, grid style={popline!50},
]
\addplot[popcurve, mark=*, mark size=1.5pt] coordinates {(1990,24.5) (2000,23.2) (2010,21.9) (2020,20.6)};
\end{axis}
\end{tikzpicture}
\end{center}
\legende{Évolution de la superficie forestière du Cameroun entre 1990 et 2020 : décroissance régulière d'environ 0,13 million d'hectares par an.}
\end{popfigure}

On observe une décroissance quasi linéaire : la forêt recule de façon régulière, sans ralentissement visible.
\item Trois fonctions menacées (d'après le document 2) : la \textbf{production de bois d'œuvre} (ressource économique) ; la fourniture de \textbf{produits forestiers non ligneux} (nourriture, pharmacopée traditionnelle) ; la \textbf{régulation du régime des pluies} ; la \textbf{protection des sols contre l'érosion} ; la conservation de plus de \num{8000} espèces végétales (réservoir de biodiversité).
\item \textbf{Texte argumenté (corrigé type) :} L'exploitation forestière est indispensable à l'économie camerounaise : elle rapporte des devises, des recettes fiscales et des emplois, notamment dans les régions de l'Est et du Sud. Mais les documents montrent une perte de 3,9 millions d'hectares en trente ans : à ce rythme, la ressource s'épuise et les générations futures en seront privées, ce qui contredit la définition même du développement durable. Concilier exploitation et conservation suppose d'abord de limiter les prélèvements à la capacité de renouvellement de la forêt : plans d'aménagement obligatoires, coupes sélectives, respect des rotations. Il faut ensuite reboiser les parcelles exploitées et développer les plantations, afin de transférer la pression hors des forêts naturelles. La création et la protection effective d'aires protégées (Dja, Korup) garantissent la conservation d'échantillons intacts de biodiversité. Enfin, la transformation locale du bois et la certification (FSC) augmentent la valeur économique de chaque arbre coupé, ce qui permet de gagner plus en exploitant moins. Ainsi, les trois piliers sont réunis : économique (emplois et devises durables), social (revenus des populations riveraines, produits forestiers non ligneux) et environnemental (forêt et biodiversité préservées).
\end{enumerate}
\textbf{Barème indicatif (sur 10) :} calculs justes et détaillés : 2 pts ; graphique correct (axes légendés, points bien placés, lecture) : 2 pts ; trois fonctions citées à partir des documents : 1,5 pt ; texte argumenté structuré (définition du développement durable, au moins trois mesures, articulation des trois piliers) : 4 pts ; qualité de la langue : 0,5 pt.

\textbf{Points de vigilance :} convertir correctement 0,13 million d'ha en \num{130000} ha ; le graphique doit avoir des axes légendés avec unités ; dans le texte, citer explicitement la notion de développement durable et s'appuyer sur les chiffres des documents (exigence de l'approche par les compétences).
\end{corrige}

\begin{corrige}[Exercice 10 — Situation-problème : un village de l'Extrême-Nord]
\begin{enumerate}
\item Chaîne de causes à effets :

\begin{popfigure}
\begin{center}
\begin{tikzpicture}[node distance=0.55cm and 1.2cm,
box/.style={draw=popblue, fill=popblueL, rounded corners, align=center, text width=4.2cm, inner sep=4pt, font=\small},
fleche/.style={-{Stealth[length=2.5mm]}, thick, popdark}]
\node[box] (a) {Feux de brousse annuels et coupe massive du bois de chauffe};
\node[box, below=of a] (b) {Disparition du couvert végétal (arbres et herbes)};
\node[box, below=of b] (c) {Sols nus : érosion par le vent et les pluies, appauvrissement en humus};
\node[box, below=of c] (d) {Baisse des rendements agricoles et assèchement des points d'eau};
\node[box, below=of d] (e) {Désertification, pauvreté et exode des jeunes};
\draw[fleche] (a) -- (b);
\draw[fleche] (b) -- (c);
\draw[fleche] (c) -- (d);
\draw[fleche] (d) -- (e);
\end{tikzpicture}
\end{center}
\legende{Chaîne de causes à effets menant des pratiques villageoises à la désertification.}
\end{popfigure}

\item Le couvert végétal protège le sol de plusieurs façons : le feuillage amortit l'impact des gouttes de pluie et freine le vent au niveau du sol ; les racines retiennent les particules de terre ; la litière et l'humus favorisent l'infiltration de l'eau. Quand la végétation disparaît, le sol nu est directement battu par les pluies et balayé par les vents de la saison sèche : les particules fines sont emportées, l'eau ruisselle au lieu de s'infiltrer, et le sol perd sa couche fertile — c'est l'érosion accélérée.
\item Plan d'action en quatre mesures :
\begin{itemize}
\item \textbf{Énergie :} diffusion de foyers améliorés (qui consomment moins de bois) et plantations villageoises de bois de chauffe à croissance rapide (par exemple le neem, \emph{Azadirachta indica}, en bosquets dédiés) $\rightarrow$ réduit la pression sur les arbres naturels.
\item \textbf{Agriculture :} agroforesterie (association arbres--cultures), culture en bandes suivant les courbes de niveau, semis sous couvert végétal $\rightarrow$ maintient le sol couvert, limite l'érosion et restaure la fertilité.
\item \textbf{Reboisement :} reboisement communautaire des zones dégradées et protection des régénérations naturelles contre les feux (pare-feu, feux précoces contrôlés) $\rightarrow$ reconstitue le couvert végétal et protège sols et points d'eau.
\item \textbf{Sensibilisation :} éducation environnementale au village et à l'école, création d'un comité de gestion des feux de brousse $\rightarrow$ fait évoluer durablement les pratiques.
\end{itemize}
\item Ce plan réunit les trois piliers du développement durable : \textbf{environnemental} (reboisement, sols et eau protégés, désertification enrayée), \textbf{économique} (rendements agricoles restaurés, bois de chauffe produit durablement, revenus maintenus) et \textbf{social} (sécurité alimentaire, implication de la communauté, jeunes incités à rester au village).
\end{enumerate}
\textbf{Barème indicatif (sur 10) :} schéma fléché complet et logique (au moins 5 étapes) : 2,5 pts ; explication du rôle protecteur du couvert végétal : 2 pts ; quatre mesures concrètes, chacune justifiée : 4 pts ; articulation des trois piliers : 1,5 pt.

\textbf{Points de vigilance :} le schéma doit être \emph{fléché} et ordonné (cause $\rightarrow$ effet) ; chaque mesure doit être \emph{réaliste localement} et accompagnée de sa justification, pas d'une simple liste.
\end{corrige}

\begin{corrige}[Exercice 11 — Analyse de dossier : l'eau potable à Yaoundé]
\begin{enumerate}
\item Les deux pressions majeures sont : la \textbf{croissance de la demande en eau} (environ $+5\ \%$ par an, liée à la croissance démographique rapide de la ville) et la \textbf{pollution des sources et rivières} (Mfoundi, Mingoa) par les rejets domestiques et industriels non traités.
\item Lorsque les rivières périurbaines sont polluées, l'eau brute captée contient davantage de matières organiques, de microbes et de produits chimiques : sa potabilisation exige alors des traitements plus nombreux et plus poussés (coagulation--floculation, filtration, désinfection renforcée au chlore), ce qui augmente les coûts de production supportés par la société de distribution et, in fine, par les usagers.
\item Consommation journalière du quartier :
\[ \num{50000}\ \text{hab.} \times 80\ \text{L/j} = \num{4000000}\ \text{L/j}, \]
soit \num{4000} m$^3$ par jour (car $1\ \text{m}^3 = \num{1000}\ \text{L}$). Volume perdu par les fuites (30 \%) :
\[ \num{4000000} \times \frac{30}{100} = \num{1200000}\ \text{L/j}, \]
soit \num{1200} m$^3$ d'eau potable perdus chaque jour — l'équivalent de la consommation quotidienne de \num{15000} habitants ($\num{1200000}/80 = \num{15000}$).
\item Trois mesures de gestion durable :
\begin{itemize}
\item \textbf{Pouvoirs publics :} réhabiliter le réseau de distribution pour réduire les fuites et construire des stations d'épuration traitant les eaux usées avant leur déversement dans le Mfoundi et la Mingoa ;
\item \textbf{Entreprises :} traiter leurs effluents industriels avant rejet et adopter des procédés économes en eau (recyclage de l'eau de process) ;
\item \textbf{Habitants :} adopter les éco-gestes (réparer les robinets qui fuient, ne pas jeter d'ordures dans les rivières, récupérer l'eau de pluie pour l'arrosage) et payer les factures d'eau pour financer l'entretien du réseau.
\end{itemize}
\end{enumerate}
\textbf{Barème indicatif (sur 10) :} deux pressions correctement identifiées à partir du document : 2 pts ; lien pollution--coût expliqué par le mécanisme de traitement : 2 pts ; calculs exacts avec unités et conversions (L $\rightarrow$ m$^3$) : 3 pts ; trois mesures correctement réparties entre acteurs : 3 pts.

\textbf{Points de vigilance :} dans les calculs, donner le résultat en litres \emph{et} en mètres cubes ($\num{1000}\ \text{L} = 1\ \text{m}^3$) ; ne pas confondre la consommation totale et le volume perdu ; attribuer chaque mesure au bon acteur, comme le demande explicitement la consigne.
\end{corrige}

\newpage

\coursec{Fiche bilan}

\courssub{Carte mentale de la leçon}

\begin{popfigure}
\begin{tikzpicture}[mindmap, grow cyclic, every node/.style={concept, execute at begin node=\hspace{0pt}},
root concept/.append style={concept color=popblue, font=\small\bfseries, minimum size=2.6cm},
level 1 concept/.append style={level distance=3.6cm, sibling angle=60, font=\scriptsize},
text=white]
\node[root concept]{Ressources naturelles et développement durable}
  child[concept color=popgreen]{ node {Typologie : eau, sols, forêts, biodiversité, ressources minérales} }
  child[concept color=poporange]{ node {Pressions humaines : déforestation, pollution, surexploitation} }
  child[concept color=poppink]{ node {Conséquences : érosion, désertification, perte de biodiversité} }
  child[concept color=popgold]{ node {Développement durable : 3 piliers (économique, social, environnemental)} }
  child[concept color=poppurple]{ node {Gestion raisonnée : quotas, reboisement, aires protégées} }
  child[concept color=popblue!70]{ node {Cameroun : parcs nationaux, forêts communautaires, lac Nyos} };
\end{tikzpicture}
\legende{Carte mentale de synthèse : les six idées forces de la leçon ND07, à restituer dans tout devoir portant sur la gestion durable des ressources naturelles.}
\end{popfigure}

\courssub{L'essentiel à retenir}

\begin{bilan}
\textbf{Définitions clés}
\begin{itemize}
  \item \cle{Ressource naturelle} : élément de la nature (eau, sol, forêt, espèces vivantes, minerais) utilisable par l'Homme pour satisfaire ses besoins.
  \item \cle{Ressource renouvelable} : ressource qui se régénère à l'échelle humaine si son prélèvement reste inférieur à sa capacité de renouvellement (eau, forêt, poissons). Une ressource renouvelable mal gérée peut s'épuiser.
  \item \cle{Ressource non renouvelable} : ressource dont le stock est fixe ou se reconstitue sur des temps géologiques (pétrole, minerais).
  \item \cle{Biodiversité} : diversité des espèces, des gènes et des écosystèmes.
  \item \cle{Développement durable} : développement qui répond aux besoins du présent sans compromettre la capacité des générations futures à répondre aux leurs (Rapport Brundtland, 1987).
  \item \cle{Gestion raisonnée} : exploitation planifiée d'une ressource qui garantit sa pérennité (prélèvement $\leq$ renouvellement).
\end{itemize}

\textbf{Pressions humaines et conséquences (à connaître par c\oe ur)}
\begin{center}
\begin{tabularx}{\linewidth}{|l|X|X|}
\hline
\rowcolor{popblueL} \textbf{Activité humaine} & \textbf{Ressource atteinte} & \textbf{Conséquence} \\
\hline
Déforestation, agriculture sur brûlis & Forêts, sols & Érosion, perte de biodiversité, désertification \\
\hline
Rejets industriels et domestiques & Eau & Pollution, eutrophisation, maladies hydriques \\
\hline
Surpêche, chasse excessive & Biodiversité & Épuisement des stocks, extinction d'espèces \\
\hline
Extraction minière et pétrolière & Sols, eau, air & Dégradation des paysages, contamination \\
\hline
Urbanisation, barrages & Sols, écosystèmes & Artificialisation, fragmentation des habitats \\
\hline
\end{tabularx}
\end{center}

\textbf{Les trois piliers du développement durable}
\begin{itemize}
  \item \cle{Économique} : activités viables et créatrices de richesses.
  \item \cle{Social} : équité, santé, éducation pour tous.
  \item \cle{Environnemental} : protection des ressources et des écosystèmes.
\end{itemize}

\textbf{Mesures de gestion durable au Cameroun}
\begin{itemize}
  \item Aires protégées : parcs nationaux (Waza, Bénoué, Korup, Campo-Ma'an), réserves de faune.
  \item Reboisement, agroforesterie, forêts communautaires et communales.
  \item Quotas de pêche et de chasse, périodes de fermeture (défense biologique).
  \item Traitement des eaux usées, tri et recyclage des déchets, lutte contre les sacs plastiques.
  \item Dégazage du lac Nyos (prévention des catastrophes limniques), énergies renouvelables (barrages hydroélectriques, solaire).
\end{itemize}

\textbf{Méthodes express}
\begin{itemize}
  \item Analyser un document : identifier la ressource $\rightarrow$ l'activité humaine $\rightarrow$ l'impact $\rightarrow$ proposer une mesure adaptée au contexte local.
  \item Toute mesure proposée doit être concrète, réaliste et reliée à l'impact décrit (jamais de réponse vague comme « protéger la nature »).
\end{itemize}
\end{bilan}

\begin{aretenir}[Checklist avant le Probatoire]
\begin{itemize}
  \item[$\square$] Je sais définir une ressource naturelle et distinguer ressource renouvelable et non renouvelable.
  \item[$\square$] Je cite les quatre grandes ressources étudiées : eau, sols, forêts, biodiversité.
  \item[$\square$] Je sais expliquer pourquoi une ressource renouvelable peut quand même s'épuiser.
  \item[$\square$] Je relie chaque pression humaine (déforestation, pollution, surexploitation) à ses conséquences précises.
  \item[$\square$] Je connais la définition officielle du développement durable (Rapport Brundtland, 1987).
  \item[$\square$] Je cite et j'explique les trois piliers du développement durable.
  \item[$\square$] Je connais au moins trois exemples camerounais d'aires protégées (Waza, Korup, Campo-Ma'an).
  \item[$\square$] Je sais proposer des mesures de gestion adaptées à une situation locale décrite dans un document.
  \item[$\square$] Je sais analyser un dossier documentaire en quatre étapes (ressource, activité, impact, mesure).
  \item[$\square$] J'évite les réponses vagues : chaque mesure proposée est concrète et justifiée.
\end{itemize}
\end{aretenir}

\findecours

\end{document}
